% Consolidation des acquis de la méthodologie du commentaire philosophique — Philosophie 1re Tech — Cours signé OnBuch+
% Programme officiel MINESEC. Compiler avec : tectonic N03-consolidation-acquis-methodologie-commentaire-philosoph.tex
\newcommand{\DOCMATIERE}{Philosophie}
\newcommand{\DOCNIVEAU}{1re Tech}
\newcommand{\DOCMODULE}{Philosophie – classe de Première technique}
\newcommand{\DOCLECON}{N03}
\newcommand{\DOCTITRE}{Consolidation des acquis de la méthodologie du commentaire philosophique}
% OnBuch+ — préambule « cours signés » ENRICHI (compilé avec tectonic / XeLaTeX)
% Système visuel « Pop » d'OnBuch+ : fond crème (jamais blanc), bordures encre
% épaisses, ombres « dures » décalées, titres Archivo Black en capitales.
% Version MATHÉMATIQUES : graphes (pgfplots/tikz), boîtes pédagogiques
% (définition, propriété, méthode, attention, le savais-tu, exercice résolu,
% exercice, TP, bilan), page de garde et signature OnBuch+ ancrée en bas.
%
% Macros attendues dans main.tex AVANT \input{preamble} :
%   \DOCMATIERE  \DOCNIVEAU  \DOCMODULE  \DOCLECON (numéro)  \DOCTITRE
\documentclass[11pt]{article}

\usepackage{fontspec}
\usepackage[french]{babel}
\usepackage[a4paper,margin=2cm,top=3.4cm,bottom=2.8cm,headheight=1.4cm,footskip=1.3cm]{geometry}
\usepackage{amsmath,amssymb,amsfonts}
\usepackage{mathtools} % \xmapsto, \coloneqq... souvent utilisés par les modèles
\usepackage{cancel} % simplifications barrées (bilans de forces)
% \abs{z} (module, valeur absolue) et \norm{u} (norme), utilisés par les modèles
\providecommand{\abs}{}\let\abs\relax\DeclarePairedDelimiter{\abs}{\lvert}{\rvert}
\providecommand{\norm}{}\let\norm\relax\DeclarePairedDelimiter{\norm}{\lVert}{\rVert}
\usepackage[table]{xcolor} % [table] : fournit aussi \rowcolors (alternance de lignes)
\usepackage{pagecolor}
\usepackage{tikz}
\usetikzlibrary{arrows.meta,positioning,calc,shapes.geometric,shapes.misc,shapes.arrows,decorations.pathmorphing,decorations.pathreplacing,decorations.markings,patterns,shadows,mindmap,trees,fit,backgrounds,matrix,intersections,angles,quotes}
\usepackage{pgfplots}
\usepackage[version=4]{mhchem} % équations nucléaires \ce{^{235}_{92}U}
\usepackage[european,siunitx]{circuitikz} % schémas électriques
\pgfplotsset{compat=1.18}
\usepgfplotslibrary{fillbetween,groupplots} % aires entre courbes (calcul intégral)
\usepackage{enumitem}
\usepackage{fancyhdr}
% [most] charge toutes les bibliothèques tcolorbox (listings, minted, raster,
% documentation...) et gonfle inutilement la mémoire TeX sur un document long
% (~300 boîtes) : on ne charge que ce qui est réellement utilisé.
\usepackage[skins,breakable]{tcolorbox}
\usepackage{graphicx}
\usepackage{colortbl}
\usepackage{booktabs}
\usepackage{tabularx}
\usepackage{multirow}
\usepackage{diagbox} % case d'en-tête diagonale des tableaux à double entrée
% Colonnes extensibles centrée / gauche / droite (C, L, R), souvent utilisées par les modèles
\newcolumntype{C}{>{\centering\arraybackslash}X}
\newcolumntype{L}{>{\raggedright\arraybackslash}X}
\newcolumntype{R}{>{\raggedleft\arraybackslash}X}
\usepackage{array}
\usepackage{multicol}
\usepackage{siunitx}
% parse-numbers=false : accepte aussi \SI{2,5 \times 10^{-3}}{...}
\sisetup{locale=FR,output-decimal-marker={,},group-minimum-digits=4,parse-numbers=false}
\DeclareSIUnit\ohm{\ensuremath{\symup{\Omega}}} % U+2126 absent de la police
\DeclareSIUnit\division{div} % réglages d'oscilloscope (ms/div, V/div)
\DeclareSIUnit\tr{tr} % tours (vitesse de rotation, ex. \SI{1500}{\tr\per\minute})
\DeclareSIUnit\molar{M} % concentration molaire (ex. \SI{3}{\molar}, \SI{1}{\micro\molar})
\DeclareSIUnit\an{an} % année (le modèle confond parfois avec \year, un compteur TeX) — ex. \SI{5}{\kilo\meter\per\an}
\DeclareSIUnit\FCFA{FCFA} % franc CFA (ex. \SI{500}{\FCFA\per\kilo\gram})
\DeclareSIUnit\degreeBrix{\textdegree Brix} % teneur en sucre d'un sirop/jus (réfractométrie)
\DeclareSIUnit\month{mois} % le modèle utilise parfois \month, un compteur TeX, comme unité de durée
\DeclareSIUnit\microliter{\micro\liter} % le modèle écrit parfois \microliter en un seul token
\DeclareSIUnit\million{millions} % dénombrement (ex. \SI{15}{\million\per\mL} spermatozoïdes)
\usepackage{qrcode}
% \degree : commande fréquemment utilisée par les modèles pour un angle ou une
% température en dehors de \SI{}{} (non fournie par les paquets chargés).
\providecommand{\degree}{^{\circ}}
% \qed (fin de démonstration), utilisé par les modèles sans amsthm
\providecommand{\qed}{\hfill$\square$}
% \barre : double barre des valeurs interdites dans un tableau de variations (tkz-tab)
\providecommand{\barre}{\ensuremath{\|}}
% environnement proof (amsthm non chargé) utilisé par les modèles
\ifdefined\proof\else\newenvironment{proof}[1][Démonstration]{\par\noindent\textit{#1.}\ }{\qed\par}\fi
\usepackage{lastpage}
\usepackage{hyperref}
\hypersetup{hidelinks}

% ============================================================
% Palette « Pop » OnBuch+ — mêmes hex que la classe P (plus_ui.dart)
% ============================================================
\definecolor{poporange}{HTML}{F59321}
\definecolor{popdark}{HTML}{E07A0C}
\definecolor{popcream}{HTML}{FFF6E8}
\definecolor{popink}{HTML}{1C1714}
\definecolor{poppurple}{HTML}{7A5AE0}
\definecolor{popgreen}{HTML}{2E9E6A}
\definecolor{poppink}{HTML}{E0457B}
\definecolor{popblue}{HTML}{2F7DE1}
\definecolor{popgold}{HTML}{C98A12}
\definecolor{popmuted}{HTML}{8A7F76}
\definecolor{popline}{HTML}{EADFD2}
\definecolor{popgray}{HTML}{8A7F76}
\colorlet{popgrey}{popgray}
% Alias tolérants (le modèle nomme parfois les couleurs différemment)
\colorlet{popwhite}{white}
\colorlet{popblack}{popink}
\colorlet{popred}{poppink}
\colorlet{popyellow}{popgold}
% Teintes claires (fonds de tableaux/encadrés) — le modèle nomme parfois
% les couleurs avec un suffixe L (ex. popblueL) sans les avoir définies.
\colorlet{poporangeL}{poporange!20}
\colorlet{popdarkL}{popdark!20}
\colorlet{poppurpleL}{poppurple!20}
\colorlet{popgreenL}{popgreen!20}
\colorlet{poppinkL}{poppink!20}
\colorlet{popblueL}{popblue!20}
\colorlet{popgoldL}{popgold!20}
\colorlet{popredL}{popred!20}
\colorlet{popgrayL}{popgray!20}
% Teintes claires pour fonds de tableaux / zones
\colorlet{popblueL}{popblue!10!white}
\colorlet{poporangeL}{poporange!14!white}
\colorlet{popgreenL}{popgreen!12!white}
\colorlet{poppurpleL}{poppurple!12!white}
\colorlet{poppinkL}{poppink!12!white}
\colorlet{popgoldL}{popgold!14!white}

\pagecolor{popcream}

% ============================================================
% Typographie
% ============================================================
\defaultfontfeatures{Ligatures=TeX}
\setmainfont{PlusJakartaSans-Medium.ttf}[Path=../../../fonts/,BoldFont=PlusJakartaSans-Bold.ttf]
% Maths en sans-serif (Fira Math) avec chiffres et lettres droites de Plus
% Jakarta, pour que les formules se fondent dans le texte.
\usepackage[math-style=ISO,bold-style=ISO]{unicode-math}
\setmathfont{FiraMath-Regular.otf}[Path=../../../fonts/]
\setmathfont{PlusJakartaSans-Medium.ttf}[Path=../../../fonts/,range={up/num,up/{Latin,latin},\mathup}]
\usepackage{icomma} % virgule décimale sans espace en mode math
% Caractères mathématiques Unicode absents des polices de texte (Plus Jakarta,
% Archivo Black) : rendus par leur équivalent mathématique, en texte comme en
% formule — sinon ils s'affichent en carré vide (ex. « Arithmétique dans ℤ »).
\usepackage{newunicodechar}
\newunicodechar{ℝ}{\ensuremath{\mathbb{R}}}
\newunicodechar{ℤ}{\ensuremath{\mathbb{Z}}}
\newunicodechar{ℂ}{\ensuremath{\mathbb{C}}}
\newunicodechar{ℕ}{\ensuremath{\mathbb{N}}}
\newunicodechar{ℚ}{\ensuremath{\mathbb{Q}}}
\newunicodechar{𝔻}{\ensuremath{\mathbb{D}}}
\newunicodechar{α}{\ensuremath{\alpha}}
\newunicodechar{β}{\ensuremath{\beta}}
\newunicodechar{γ}{\ensuremath{\gamma}}
\newunicodechar{δ}{\ensuremath{\delta}}
\newunicodechar{ε}{\ensuremath{\varepsilon}}
\newunicodechar{θ}{\ensuremath{\theta}}
\newunicodechar{λ}{\ensuremath{\lambda}}
\newunicodechar{μ}{\ensuremath{\mu}}
\newunicodechar{σ}{\ensuremath{\sigma}}
\newunicodechar{τ}{\ensuremath{\tau}}
\newunicodechar{φ}{\ensuremath{\varphi}}
\newunicodechar{ψ}{\ensuremath{\psi}}
\newunicodechar{ω}{\ensuremath{\omega}}
\newunicodechar{Σ}{\ensuremath{\Sigma}}
% majuscules produites par \MakeUppercase dans les titres (α-aminés -> Α)
\newunicodechar{Α}{\ensuremath{\alpha}}
\newunicodechar{Β}{\ensuremath{\beta}}
\newunicodechar{Γ}{\ensuremath{\Gamma}}
\newunicodechar{Δ}{\ensuremath{\Delta}}
\newunicodechar{Ε}{\ensuremath{\varepsilon}}
\newunicodechar{Θ}{\ensuremath{\Theta}}
\newunicodechar{Λ}{\ensuremath{\Lambda}}
\newunicodechar{Μ}{\ensuremath{\mu}}
\newunicodechar{Τ}{\ensuremath{\tau}}
\newunicodechar{Φ}{\ensuremath{\Phi}}
\newunicodechar{Ψ}{\ensuremath{\Psi}}
\newunicodechar{Ω}{\ensuremath{\Omega}}
\newunicodechar{↦}{\ensuremath{\mapsto}}
\newunicodechar{∈}{\ensuremath{\in}}
\newunicodechar{∉}{\ensuremath{\notin}}
\newunicodechar{∧}{\ensuremath{\wedge}}
\newunicodechar{⇔}{\ensuremath{\Leftrightarrow}}
\newunicodechar{⟺}{\ensuremath{\Longleftrightarrow}}
\newunicodechar{⇒}{\ensuremath{\Rightarrow}}
\newunicodechar{⟹}{\ensuremath{\Longrightarrow}}
\newunicodechar{∘}{\ensuremath{\circ}}
\newunicodechar{∪}{\ensuremath{\cup}}
\newunicodechar{∩}{\ensuremath{\cap}}
\newunicodechar{≡}{\ensuremath{\equiv}}
\newunicodechar{⊂}{\ensuremath{\subset}}
\newunicodechar{⊥}{\ensuremath{\perp}}
\newunicodechar{∃}{\ensuremath{\exists}}
\newunicodechar{∀}{\ensuremath{\forall}}
\newunicodechar{∛}{\ensuremath{\sqrt[3]{\,}}}
\newunicodechar{∜}{\ensuremath{\sqrt[4]{\,}}}
\newunicodechar{⃗}{\ensuremath{{}^{\to}}}
\newunicodechar{ⁿ}{\ifmmode^{n}\else\textsuperscript{\ensuremath{n}}\fi}
\newunicodechar{ˣ}{\ifmmode^{x}\else\textsuperscript{\ensuremath{x}}\fi}
\newunicodechar{ᵇ}{\ifmmode^{b}\else\textsuperscript{\ensuremath{b}}\fi}
\newunicodechar{ʳ}{\ifmmode^{r}\else\textsuperscript{\ensuremath{r}}\fi}
\newunicodechar{ᵘ}{\ifmmode^{u}\else\textsuperscript{\ensuremath{u}}\fi}
\newunicodechar{ʸ}{\ifmmode^{y}\else\textsuperscript{\ensuremath{y}}\fi}
\newunicodechar{ᵃ}{\ifmmode^{a}\else\textsuperscript{\ensuremath{a}}\fi}
\newunicodechar{ᵉ}{\ifmmode^{e}\else\textsuperscript{\ensuremath{e}}\fi}
\newunicodechar{ᵏ}{\ifmmode^{k}\else\textsuperscript{\ensuremath{k}}\fi}
\newunicodechar{ᵖ}{\ifmmode^{p}\else\textsuperscript{\ensuremath{p}}\fi}
\newunicodechar{ᴺ}{\ifmmode^{N}\else\textsuperscript{\ensuremath{N}}\fi}
\newunicodechar{ᶜ}{\ifmmode^{c}\else\textsuperscript{\ensuremath{c}}\fi}
\newunicodechar{ʰ}{\ifmmode^{h}\else\textsuperscript{\ensuremath{h}}\fi}
\newunicodechar{ᵠ}{\ifmmode^{q}\else\textsuperscript{\ensuremath{q}}\fi}
\newunicodechar{ₙ}{\ifmmode_{n}\else\textsubscript{\ensuremath{n}}\fi}
\newunicodechar{ₐ}{\ifmmode_{a}\else\textsubscript{\ensuremath{a}}\fi}
\newunicodechar{ᵢ}{\ifmmode_{i}\else\textsubscript{\ensuremath{i}}\fi}
\newunicodechar{ᵣ}{\ifmmode_{r}\else\textsubscript{\ensuremath{r}}\fi}
\newunicodechar{ₖ}{\ifmmode_{k}\else\textsubscript{\ensuremath{k}}\fi}
\newunicodechar{ᵦ}{\ifmmode_{\beta}\else\textsubscript{\ensuremath{\beta}}\fi}
\newunicodechar{ₓ}{\ifmmode_{x}\else\textsubscript{\ensuremath{x}}\fi}
\newfontfamily\popdisplay{ArchivoBlack-Regular.ttf}[Path=../../../fonts/]
\newfontfamily\poplogo{SpaceGrotesk-Bold.ttf}[Path=../../../fonts/]
\newfontfamily\popsemi{PlusJakartaSans-SemiBold.ttf}[Path=../../../fonts/]
\newfontfamily\popxbold{PlusJakartaSans-ExtraBold.ttf}[Path=../../../fonts/]
\newfontfamily\popxboldls{PlusJakartaSans-ExtraBold.ttf}[Path=../../../fonts/,LetterSpace=3.0]

\setlist[enumerate]{leftmargin=1.7em,itemsep=5pt,topsep=4pt}
\setlist[itemize]{leftmargin=1.7em,itemsep=4pt,topsep=4pt,label={\color{poporange}\textbullet}}
\setlength{\parindent}{0pt}
\setlength{\parskip}{5pt}
\renewcommand{\arraystretch}{1.25}

% pgfplots : style Pop par défaut
\pgfplotsset{
  every axis/.append style={axis line style={popink,line width=1pt},tick style={popink},
    grid style={popline,line width=0.5pt},label style={font=\small},tick label style={font=\footnotesize}},
  popcurve/.style={poporange,line width=1.8pt,no markers,smooth},
}
% Même style disponible dans un \draw TikZ simple (hors pgfplots)
\tikzset{popcurve/.style={poporange,line width=1.8pt}}
\tikzset{popgrid/.style={popline,very thin}}
\colorlet{popcurve}{poporange} % le nom du style est parfois utilisé comme couleur

% ============================================================
% Pill / squiggle
% ============================================================
\newcommand{\poppill}[3]{%
  \tikz[baseline=-0.55ex]\node[draw=popink,line width=1.2pt,rounded corners=2.6pt,%
    fill=#2,inner xsep=6.5pt,inner ysep=3pt]{%
    \popxboldls\fontsize{7}{7}\selectfont\color{#3}\MakeUppercase{#1}};%
}
\newcommand{\popsquiggle}[1]{%
  \tikz[baseline=-0.4ex]{
    \draw[#1,line width=2.6pt,line cap=round]
      plot [smooth] coordinates {(0,0.09) (0.34,0.01) (0.68,0.21) (1.02,0.01) (1.36,0.10)};
  }%
}
% Mot-clé surligné (marqueur orange sous le texte)
\newcommand{\cle}[1]{{\popxbold\color{popdark}#1}}

% ============================================================
% En-tête / pied
% ============================================================
\pagestyle{fancy}
\fancyhf{}
\renewcommand{\headrulewidth}{0pt}
\renewcommand{\footrulewidth}{0pt}
\lhead{\raisebox{-0.15ex}{\poplogo\fontsize{13}{13}\selectfont\color{popink}OnBuch\color{poporange}+}}
\rhead{\poppill{\DOCMATIERE\ \textbullet\ \DOCNIVEAU\ \textbullet\ Leçon \DOCLECON}{white}{popink}}
\lfoot{\popxbold\fontsize{7.6}{7.6}\selectfont\color{popmuted}\MakeUppercase{Cours signé OnBuch+}\ \textbullet\ {\popsemi\DOCTITRE}}
\rfoot{\tikz[baseline=-0.6ex]\node[rounded corners=8.5pt,fill=popink,minimum height=17pt,minimum width=17pt,inner xsep=5pt,inner sep=0]{\popxbold\fontsize{8}{8}\selectfont\color{popcream}\thepage\,/\,\pageref*{LastPage}};}

% ============================================================
% Titres
% ============================================================
\newcommand{\chaptitre}[2]{%
  \begin{tcolorbox}[enhanced,colback=poporange,colframe=popink,boxrule=2.6pt,arc=11pt,
      drop shadow=popink,left=15pt,right=15pt,top=12pt,bottom=11pt,before skip=3pt,after skip=18pt]
    \poppill{#2}{popink}{popcream}\par\vspace{9pt}
    {\raggedright\hyphenpenalty=10000\exhyphenpenalty=10000\popdisplay\fontsize{22}{24}\selectfont\color{popcream}\MakeUppercase{#1}\par}
  \end{tcolorbox}
}
% Section numérotée : pastille orange avec le numéro + titre Archivo
\newcounter{popsec}
\newcommand{\coursec}[1]{%
  \stepcounter{popsec}\setcounter{popsub}{0}%
  \par\vspace{16pt}\noindent
  \begin{minipage}{\linewidth}
  \tikz[baseline=-0.7ex]\node[draw=popink,line width=1.6pt,fill=poporange,rounded corners=4pt,minimum size=22pt,inner sep=0]{\popdisplay\fontsize{12}{12}\selectfont\color{popcream}\thepopsec};%
  \hspace{8pt}{\popdisplay\fontsize{15}{17}\selectfont\color{popink}\MakeUppercase{#1}}\par
  \vspace{4pt}\noindent{\color{poporange}\rule{36pt}{3.6pt}}
  \end{minipage}\par\nopagebreak\vspace{8pt}
}
\newcounter{popsub}
\newcommand{\courssub}[1]{%
  \stepcounter{popsub}%
  \par\vspace{9pt}\noindent{\popxbold\fontsize{11.5}{13}\selectfont\color{popink}{\color{poporange}\thepopsec.\thepopsub}\hspace{6pt}#1}\par\nopagebreak\vspace{4pt}
}

% ============================================================
% Boîtes pédagogiques — même recette : fond blanc, bordure épaisse colorée,
% ombre dure encre, badge plein. Titre optionnel [#1].
% ============================================================
\tcbset{popbase/.style={breakable,enhanced,colback=white,boxrule=2.3pt,arc=9pt,
  shadow={3pt}{-3pt}{0pt}{popink},left=14pt,right=14pt,top=11pt,bottom=11pt,before skip=11pt,after skip=15pt,
  fontupper=\popsemi\fontsize{10.6}{14.5}\selectfont\color{popink},
  fontlower=\popsemi\fontsize{10.6}{14.5}\selectfont\color{popink}}}
% Coche dessinée (le glyphe ✓ n'existe pas dans les polices du document)
\newcommand{\popcheck}[1]{\tikz[baseline=-0.5ex]\draw[#1,line width=1.6pt,line cap=round,line join=round](-0.13,0.0)--(-0.04,-0.09)--(0.13,0.1);}
\providecommand{\coche}{\popcheck{popgreen}}
\providecommand{\resultat}[1]{\par\smallskip\noindent\fcolorbox{popgreen}{popgreenL}{\parbox{\dimexpr\linewidth-2\fboxsep-2\fboxrule\relax}{\textbf{Résultat.} #1}}\par\smallskip}
\newcommand{\popboxhead}[3]{\poppill{#1}{#2}{white}\ifx\relax#3\relax\else\hspace{6pt}{\popxbold\fontsize{10.6}{12}\selectfont\color{popink}#3}\fi\par\vspace{7pt}}

\newtcolorbox{aretenir}[1][]{popbase,colframe=popgreen,before upper={\popboxhead{À retenir}{popgreen}{#1}}}
\newtcolorbox{exemplebox}[1][]{popbase,colframe=poporange,before upper={\popboxhead{Exemple}{poporange}{#1}}}
\newtcolorbox{definition}[1][]{popbase,colframe=popblue,before upper={\popboxhead{Définition}{popblue}{#1}}}
\newtcolorbox{propriete}[1][]{popbase,colframe=poppurple,before upper={\popboxhead{Propriété}{poppurple}{#1}}}
\newtcolorbox{methode}[1][]{popbase,colframe=popgold,before upper={\popboxhead{Méthode}{popgold}{#1}}}
\newtcolorbox{attention}[1][]{popbase,colframe=poppink,before upper={\popboxhead{Attention}{poppink}{#1}}}
\newtcolorbox{experience}[1][]{popbase,colframe=popblue,colback=popblueL,before upper={\popboxhead{Expérience}{popblue}{#1}}}
% « Le savais-tu ? » : carte violette pleine (contexte, culture, Cameroun)
\newtcolorbox{savaistu}[1][]{popbase,colback=poppurple,colframe=popink,
  fontupper=\popsemi\fontsize{10.4}{14.2}\selectfont\color{white},
  before upper={\poppill{Le savais-tu\,?}{white}{poppurple}\ifx\relax#1\relax\else\hspace{6pt}{\popxbold\color{white}#1}\fi\par\vspace{7pt}}}
% Exercice résolu : énoncé en haut, \tcblower puis la solution sur fond crème
\newcounter{popexo}\newcounter{popexr}
\newtcolorbox{exoresolu}[1][]{popbase,colframe=poporange,colbacklower=popcream,
  segmentation style={popink,line width=1.2pt,dashed},
  before upper={\stepcounter{popexr}\popboxhead{Exercice résolu \thepopexr}{poporange}{#1}},
  before lower={\poppill{Solution}{popink}{popcream}\par\vspace{6pt}}}
% Exercice (non résolu en place) — la correction est dans la partie Corrigés
\newtcolorbox{exercice}[1][]{popbase,colframe=popink,
  before upper={\stepcounter{popexo}\popboxhead{Exercice \thepopexo}{popink}{#1}}}
% Corrigé
\newtcolorbox{corrige}[1][]{popbase,colframe=popgreen,colback=popgreenL,
  before upper={\popboxhead{Corrigé}{popgreen}{#1}}}
% Objectifs / prérequis / bilan
\newtcolorbox{objectifs}{popbase,colframe=popink,colback=poporangeL,
  before upper={\popboxhead{Objectifs de la leçon}{popink}{}}}
\newtcolorbox{prerequis}{popbase,colframe=popink,colback=popblueL,
  before upper={\popboxhead{Prérequis}{popblue}{}}}
\newtcolorbox{bilan}{popbase,colframe=popink,colback=white,boxrule=2.8pt,
  before upper={\popboxhead{Fiche bilan}{poporange}{L'essentiel à connaître pour l'examen}}}
% Figure encadrée avec légende
\newtcolorbox{popfigure}[1][]{enhanced,breakable=false,colback=white,colframe=popline,boxrule=1.4pt,arc=8pt,
  left=10pt,right=10pt,top=10pt,bottom=8pt,before skip=10pt,after skip=12pt,halign=center,
  lower separated=false,halign lower=center,
  fontlower=\popsemi\fontsize{9}{11.5}\selectfont\color{popmuted},
  #1}
\newcommand{\legende}[1]{\par\vspace{6pt}{\popsemi\fontsize{9}{11.5}\selectfont\color{popmuted}#1}}

% ============================================================
% Signature OnBuch+
% ============================================================
\newcommand{\poplogoinline}[1]{{\poplogo\fontsize{#1}{#1}\selectfont\color{popink}OnBuch\color{poporange}+}}
\newcommand{\signatureonbuch}{%
  \begin{tcolorbox}[enhanced,colback=popink,colframe=popink,boxrule=2.2pt,arc=10pt,
      drop shadow=poporange,left=14pt,right=14pt,top=10pt,bottom=10pt,before skip=0pt,after skip=0pt]
    \tikz[baseline=-0.7ex]\node[circle,fill=popgreen,draw=popcream,line width=1.3pt,minimum size=22pt,inner sep=0]{\popcheck{white}};%
    \hspace{9pt}\begin{minipage}[c]{0.62\linewidth}
      {\popxbold\fontsize{12}{13}\selectfont\color{popcream}Signé\ \textbullet\ {\poplogo OnBuch\color{poporange}+}}\par\vspace{2pt}
      {\raggedright\popsemi\fontsize{8}{10.5}\selectfont\color{popline}\MakeUppercase{Équipe pédagogique OnBuch+}\\\MakeUppercase{Conforme au programme officiel MINESEC}\par}
    \end{minipage}\hfill
    {\poplogo\fontsize{22}{22}\selectfont\color{popcream}OnBuch\color{poporange}+}
  \end{tcolorbox}%
}

% ============================================================
% Page de garde
% #1 titre · #2 matière · #3 niveau · #4 module · #5 n° leçon · #6 durée ·
% #7 description / objectifs (texte ou itemize)
% ============================================================
\newcommand{\pagegarde}[7]{%
  \thispagestyle{empty}
  \newgeometry{a4paper,margin=1.7cm,top=1.6cm,bottom=1.4cm}%
  \begin{tikzpicture}[remember picture,overlay]
    \fill[poporange!18] ([xshift=-3.2cm,yshift=-3.6cm]current page.north east) circle (4.6cm);
    \fill[poppurple!14] ([xshift=2.4cm,yshift=7.6cm]current page.south west) circle (3.8cm);
  \end{tikzpicture}%
  {\poplogo\fontsize{30}{30}\selectfont\color{popink}OnBuch\color{poporange}+}\hfill
  \raisebox{0.6ex}{\poppill{Cours signé \textbullet\ Premium}{popink}{popcream}}\par
  \vspace{1pt}\popsquiggle{poppurple}\par
  \vspace{8pt}
  \begin{tcolorbox}[enhanced,colback=poporange,colframe=popink,boxrule=2.8pt,arc=13pt,
      drop shadow=popink,left=18pt,right=18pt,top=12pt,bottom=13pt,after skip=10pt]
    \poppill{#2 \textbullet\ #3}{popink}{popcream}\hspace{5pt}\poppill{#4}{white}{popink}\par\vspace{8pt}
    {\popdisplay\fontsize{14}{14}\selectfont\color{popink}LEÇON #5}\par\vspace{4pt}
    {\raggedright\hyphenpenalty=10000\exhyphenpenalty=10000\popdisplay\fontsize{25}{27}\selectfont\color{popcream}\MakeUppercase{#1}\par}
    \vspace{8pt}
    {\popxbold\fontsize{9.5}{9.5}\selectfont\color{popink}\MakeUppercase{Durée conseillée : #6}}
  \end{tcolorbox}
  \begin{tcolorbox}[enhanced,colback=white,colframe=popink,boxrule=2.2pt,arc=10pt,
      drop shadow=popink,left=16pt,right=16pt,top=10pt,bottom=10pt,after skip=8pt]
    \poppill{Dans cette leçon}{poppurple}{white}\par\vspace{5pt}
    \popsemi\fontsize{9.3}{12.6}\selectfont\color{popink}#7
  \end{tcolorbox}
  \noindent\begin{minipage}[t]{0.487\linewidth}
    \begin{tcolorbox}[enhanced,colback=popgreen,colframe=popink,boxrule=2.2pt,arc=10pt,
        drop shadow=popink,left=11pt,right=11pt,top=10pt,bottom=10pt,height=3.1cm,valign=center]
      \tcbox[colback=white,colframe=popink,boxrule=1.3pt,arc=5pt,boxsep=0pt,nobeforeafter,
        left=3pt,right=3pt,top=3pt,bottom=3pt,valign=center]{\qrcode[height=1.9cm]{https://play.google.com/store/apps/details?id=com.luvvix.onbuch&hl=fr}}%
      \hspace{8pt}\begin{minipage}[c]{0.5\linewidth}
        {\popdisplay\fontsize{11}{12.5}\selectfont\color{white}\MakeUppercase{Télécharge OnBuch}}\par\vspace{4pt}
        {\raggedright\popsemi\fontsize{8.4}{11}\selectfont\color{white}Gratuit \textbullet\ Google Play\\Tuteur IA, annales, concours, résultats\par}
      \end{minipage}
    \end{tcolorbox}
  \end{minipage}\hfill
  \begin{minipage}[t]{0.487\linewidth}
    \begin{tcolorbox}[enhanced,colback=poppurple,colframe=popink,boxrule=2.2pt,arc=10pt,
        drop shadow=popink,left=11pt,right=11pt,top=10pt,bottom=10pt,height=3.1cm,valign=center]
      \tikz[baseline=-0.7ex]\node[circle,fill=white,draw=popink,line width=1.3pt,minimum size=1.6cm,inner sep=0]{\popdisplay\fontsize{24}{24}\selectfont\color{poppurple}+};%
      \hspace{8pt}\begin{minipage}[c]{0.6\linewidth}
        {\popdisplay\fontsize{11}{12.5}\selectfont\color{white}\MakeUppercase{Passe à OnBuch+}}\par\vspace{4pt}
        {\raggedright\popsemi\fontsize{8.4}{11}\selectfont\color{white}Tous les cours signés, Tuteur IA illimité, fiches PDF — onglet OnBuch+ de l'app\par}
      \end{minipage}
    \end{tcolorbox}
  \end{minipage}
  \vspace{8pt}
  \popcontenu
  \vfill
  \signatureonbuch
  \restoregeometry
  \newpage
}

% Rangée « Ce cours contient » (page de garde)
\newcommand{\poptile}[3]{% #1 couleur #2 gros texte #3 légende
  \begin{tcolorbox}[enhanced,colback=#1,colframe=popink,boxrule=2pt,arc=9pt,shadow={2.5pt}{-2.5pt}{0pt}{popink},
      left=6pt,right=6pt,top=9pt,bottom=9pt,halign=center,valign=center,height=1.75cm,width=\linewidth,nobeforeafter]
    {\popdisplay\fontsize{12.5}{13}\selectfont\color{white}\MakeUppercase{#2}}\par\vspace{3pt}
    {\centering\popsemi\fontsize{7.8}{9.5}\selectfont\color{white}#3\par}
  \end{tcolorbox}}
\newcommand{\popcontenu}{%
  \par\noindent{\popxboldls\fontsize{7.5}{7.5}\selectfont\color{popmuted}CE COURS SIGNÉ CONTIENT}\par\vspace{7pt}
  \noindent\begin{minipage}[t]{0.235\linewidth}\poptile{popblue}{Cours}{complet, illustré, pas à pas}\end{minipage}\hfill
  \begin{minipage}[t]{0.235\linewidth}\poptile{poporange}{Méthodes}{et exercices résolus}\end{minipage}\hfill
  \begin{minipage}[t]{0.235\linewidth}\poptile{poppink}{Exercices}{type examen + corrigés}\end{minipage}\hfill
  \begin{minipage}[t]{0.235\linewidth}\poptile{popgreen}{Fiche bilan}{l'essentiel à retenir}\end{minipage}\par}

% Sommaire
\newtcolorbox{sommaire}{enhanced,colback=white,colframe=popink,boxrule=2.2pt,arc=10pt,
  drop shadow=popink,left=18pt,right=18pt,top=13pt,bottom=13pt,before skip=6pt,after skip=20pt,
  before upper={\poppill{Sommaire}{poppurple}{white}\par\vspace{9pt}\popsemi\fontsize{10.6}{14.5}\selectfont\color{popink}}}

% Signature de fin de cours — poussée tout en bas de la dernière page
\newcommand{\findecours}{%
  \par\vfill\nopagebreak
  \begin{center}\popsquiggle{poporange}\end{center}\vspace{4pt}
  \signatureonbuch
}
\AtBeginDocument{\def\llbracket{\mathopen{[\mkern-3.5mu[}}\def\rrbracket{\mathclose{]\mkern-3.5mu]}}} % intervalles d'entiers (glyphe ⟦ absent de la police)
% Glyphes absents de Fira Math : redéfinis à partir de glyphes disponibles
\AtBeginDocument{%
  \def\setminus{\mathbin{\backslash}}\let\smallsetminus\setminus
  \def\vdots{\mathord{\vbox{\baselineskip3.2pt\lineskiplimit0pt\kern4pt\hbox{.}\hbox{.}\hbox{.}}}}%
  \def\ddots{\mathinner{\mkern1mu\raise6.5pt\hbox{.}\mkern2mu\raise3.6pt\hbox{.}\mkern2mu\raise0.7pt\hbox{.}\mkern1mu}}%
  \def\triangle{\mathord{\scalebox{0.9}{$\symup{\Delta}$}}}%
  \def\nabla{\mathord{\raisebox{\depth}{\scalebox{1}[-1]{$\symup{\Delta}$}}}}%
  \def\checkmark{\popcheck{popdark}}%
  \def\star{\mathbin{\ast}}\def\bigstar{\mathord{\ast}}%
  \def\diamond{\mathbin{\scalebox{0.6}{$\Diamond$}}}%
  \def\bigcup{\mathop{\mathchoice{\raisebox{-0.3em}{\scalebox{1.7}{$\cup$}}}{\scalebox{1.25}{$\cup$}}{\cup}{\cup}}\displaylimits}%
  \def\bigcap{\mathop{\mathchoice{\raisebox{-0.3em}{\scalebox{1.7}{$\cap$}}}{\scalebox{1.25}{$\cap$}}{\cap}{\cap}}\displaylimits}%
  \def\lesseqqgtr{\mathrel{\lessgtr}}\def\bigoplus{\mathop{\oplus}\displaylimits}%
}
\providecommand{\ssi}{si et seulement si} % abréviation fréquente
\AtBeginDocument{\providecommand{\wideparen}{\overparen}} % arc de cercle

%% FIGFIX-BEGIN v5 — correctifs globaux des figures (docs/onbuch-generation/tools/figfix)
\usepackage{adjustbox}\usepackage{etoolbox}\tcbuselibrary{raster}
% toute figure TikZ/pgfplots plus large que la ligne est réduite à la largeur disponible
\BeforeBeginEnvironment{tikzpicture}{\begin{adjustbox}{max width=\linewidth}}
\AfterEndEnvironment{tikzpicture}{\end{adjustbox}}
% légendes pgfplots lisibles par-dessus les courbes
\pgfplotsset{every axis/.append style={legend style={fill=white,fill opacity=0.92,text opacity=1,draw=black!15,inner sep=3pt}}}
% étiquettes d'arêtes (syntaxe edge["..."]) avec fond blanc
\tikzset{every edge quotes/.append style={fill=white,inner sep=1.2pt}}
%% FIGFIX-END
\begin{document}

\pagegarde{Consolidation des acquis de la méthodologie du commentaire philosophique}{Philosophie}{1re Tech}{Philosophie – classe de Première technique}{N03}{2 séances (fiche de progression harmonisée)}{Cette leçon de consolidation réactive et systématise l'ensemble des principes méthodologiques du commentaire philosophique de texte : lecture analytique, dégagement de la thèse, problématisation, plan détaillé et rédaction. À travers des travaux pratiques guidés sur des textes d'auteurs et des situations concrètes de la vie camerounaise, l'élève s'entraîne à construire et justifier une démarche complète en vue du Baccalauréat technique.
\vspace{4pt}
\begin{multicols}{2}\small
\begin{itemize}
\item À la fin de cette leçon, je sais rappeler les étapes et les principes méthodologiques du commentaire philosophique de texte.
\item À la fin de cette leçon, je sais effectuer une lecture analytique d'un texte : repérage du thème, de la thèse, des arguments et de la structure.
\item À la fin de cette leçon, je sais formuler une problématique pertinente à partir d'un texte donné.
\item À la fin de cette leçon, je sais construire une introduction de commentaire en quatre temps (amorce, présentation, problématique, annonce du plan).
\item À la fin de cette leçon, je sais élaborer un plan détaillé cohérent qui suit la progression du texte sans le paraphraser.
\item À la fin de cette leçon, je sais expliquer et discuter une thèse en mobilisant des exemples tirés de la vie courante au Cameroun.
\end{itemize}\end{multicols}}

\begin{sommaire}
\begin{multicols}{2}
\begin{enumerate}
\item Rappel des fondamentaux : qu'est-ce que commenter un texte philosophique ?
\item La lecture analytique du texte
\item La problématisation et l'introduction
\item Le développement : expliquer, illustrer, discuter
\item La conclusion et la relecture
\item Entraînement intégré : commentaire complet guidé (complément Bac)
\item Activité d'intégration
\item Méthodes et exercices résolus
\item Exercices
\item Corrigés des exercices
\item Fiche bilan
\end{enumerate}
\end{multicols}
\end{sommaire}

\begin{objectifs}
\begin{itemize}
\item À la fin de cette leçon, je sais rappeler les étapes et les principes méthodologiques du commentaire philosophique de texte.
\item À la fin de cette leçon, je sais effectuer une lecture analytique d'un texte : repérage du thème, de la thèse, des arguments et de la structure.
\item À la fin de cette leçon, je sais formuler une problématique pertinente à partir d'un texte donné.
\item À la fin de cette leçon, je sais construire une introduction de commentaire en quatre temps (amorce, présentation, problématique, annonce du plan).
\item À la fin de cette leçon, je sais élaborer un plan détaillé cohérent qui suit la progression du texte sans le paraphraser.
\item À la fin de cette leçon, je sais expliquer et discuter une thèse en mobilisant des exemples tirés de la vie courante au Cameroun.
\item À la fin de cette leçon, je sais rédiger une conclusion qui répond à la problématique et ouvre sur une question.
\item À la fin de cette leçon, je sais rédiger et justifier une démarche, une réponse ou une conclusion argumentée.
\item À la fin de cette leçon, je sais repérer et éviter les erreurs fréquentes du commentaire (paraphrase, hors-sujet, récitation de cours).
\end{itemize}
\end{objectifs}

\begin{prerequis}
\begin{itemize}
\item Notion de texte philosophique et distinction texte argumentatif / texte narratif (Seconde)
\item Premiers éléments de méthodologie du commentaire vus en début d'année de Première
\item Notions de thèse, d'argument et d'exemple
\item Capacité à rédiger un paragraphe argumenté en français
\item Notions de l'unité en cours (ex. la conscience, la technique, le devoir) selon la progression
\end{itemize}
\end{prerequis}

\newpage

\coursec{Rappel des fondamentaux : qu'est-ce que commenter un texte philosophique ?}

Lors d'un débat au lycée de Nkongsamba, un élève de Première technique affirme : « La technique nous rend libres ». Son camarade rétorque : « Non, elle nous rend esclaves de nos machines ! » Chacun cite des exemples — le téléphone portable, la moto-taxi, la machine à coudre de l'atelier — mais aucun ne sait organiser sa pensée en un raisonnement rigoureux. Devant un texte de philosophe au Baccalauréat technique, comment passer d'une impression spontanée à une analyse méthodique ? Comment lire, comprendre, problématiser et commenter un texte philosophique sans paraphraser ni délayer ? Cette leçon consolide, par la pratique, toute la démarche du commentaire philosophique.

\courssub{Définition et nature du commentaire philosophique}

\begin{definition}[Commentaire philosophique]
Exercice d'explication et de discussion ordonnée d'un texte d'auteur, qui vise à en dégager la \cle{thèse} (la réponse principale de l'auteur), la \cle{structure} (l'enchaînement logique des arguments) et la \cle{portée} (la signification et les limites de la pensée développée).
\end{definition}

Commenter n'est pas « dire ce que l'on pense du sujet » (ce serait une dissertation), ni « répéter le texte avec ses propres mots » (ce serait une paraphrase), ni « condenser le texte » (ce serait un résumé). Le commentaire est un \textbf{aller-retour constant} entre le texte et la pensée : on part du texte pour y revenir enrichi d'une compréhension plus profonde.

\begin{center}
\begin{popfigure}
\begin{tikzpicture}[
    node distance=1.2cm and 1.5cm,
    box/.style={rectangle, draw=popdark, thick, rounded corners, minimum height=1.8cm, text width=4.5cm, align=center, font=\small, fill=popcream},
    header/.style={rectangle, fill=popblue, text=white, font=\bfseries\small, rounded corners, minimum height=1cm, text width=4.5cm, align=center},
    title/.style={font=\bfseries\large, text=popdark}
]
% Headers
\node[header] (h1) {RÉSUMÉ};
\node[header, right=of h1] (h2) {PARAPHRASE};
\node[header, right=of h2] (h3) {COMMENTAIRE};

% Boxes: Definition
\node[box, below=of h1] (d1) {\textbf{But :} Restituer l'essentiel en moins de mots. \textbf{Action :} Sélectionner les idées principales, supprimer les détails, garder l'ordre de l'auteur.};
\node[box, below=of h2] (d2) {\textbf{But :} Rendre le texte plus accessible. \textbf{Action :} Suivre le texte phrase par phrase, remplacer le vocabulaire difficile par des synonymes, garder la même structure linéaire.};
\node[box, below=of h3] (d3) {\textbf{But :} Expliquer \emph{et} discuter la pensée. \textbf{Action :} Reconstruire la logique interne (thèse, arguments, concepts), montrer \emph{comment} l'auteur démontre, évaluer la solidité du raisonnement.};

% Boxes: Defaut
\node[box, below=of d1, fill=poppinkL] (f1) {\textbf{Défaut au Bac :} Hors-sujet. Le correcteur attend une analyse, pas une synthèse. Aucune valeur ajoutée philosophique.};
\node[box, below=of d2, fill=poppinkL] (f2) {\textbf{Défaut au Bac :} Paraphrase sanctionnée (note < 6/20). L'élève suit le fil du texte sans en dévoiler l'architecture. « Il dit que... puis il dit que... »};
\node[box, below=of d3, fill=popgreenL] (f3) {\textbf{Attendu au Bac :} Preuve de compréhension \emph{et} d'esprit critique. L'élève montre qu'il a « fait parler le texte » en en exposant la mécanique rationnelle.};

% Arrows for flow
\draw[->, thick, popblue] (h1.south) -- (d1.north);
\draw[->, thick, popblue] (h2.south) -- (d2.north);
\draw[->, thick, popblue] (h3.south) -- (d3.north);
\draw[->, thick, poporange] (d1.south) -- (f1.north);
\draw[->, thick, poporange] (d2.south) -- (f2.north);
\draw[->, thick, popgreen] (d3.south) -- (f3.north);

% Title
\node[title, above=0.3cm of h1.north, xshift=6.5cm] {Tableau comparatif : Résumé / Paraphrase / Commentaire};
\end{tikzpicture}
\legende{Distinction opérationnelle entre trois exercices souvent confondus. Seule la dernière colonne correspond à l'exigence du Baccalauréat technique.}
\end{popfigure}
\end{center}

\courssub{Les trois exigences fondamentales}

Un commentaire réussi repose sur un tripode indissociable. L'absence de l'un fait chuter la note.

\begin{propriete}[Fidélité au texte]
On ne fait pas dire au texte ce qu'il ne dit pas. Toute interprétation doit s'ancrer sur une \cle{citation} précise (numéro de ligne ou guillemets) ou sur un renvoi explicite à un passage. L'exégèse précède la critique.
\end{propriete}

\begin{propriete}[Profondeur de l'analyse]
Il ne suffit pas de dire « l'auteur affirme que... ». Il faut dire : « L'auteur \cle{argumente} en montrant que... parce que... ». L'analyse dégage les \cle{concepts} opératoires (ex: \emph{aliénation}, \cle{progrès}, \cle{liberté}), les \cle{articulations logiques} (donc, or, mais, cependant) et les \cle{présupposés} implicites.
\end{propriete}

\begin{propriete}[Esprit critique / Discussion]
Après avoir \emph{expliqué} (le « quoi » et le « comment »), on \emph{discute} (le « combien » et le « pourquoi »). On teste la solidité des arguments : y a-t-il des \cle{angles morts} ? Des \cle{contre-exemples} (au Cameroun ou ailleurs) ? Des \cle{objections} classiques d'autres auteurs ? La discussion n'est pas une opinion gratuite (« je ne suis pas d'accord »), mais un \cle{contre-argument} motivé.
\end{propriete}

\begin{attention}[Erreur fatale : Réciter son cours au lieu d'analyser le texte]
Un candidat voit le mot « technique » dans le texte et écrit trois pages sur « La technique chez Heidegger / Ellul / Marx » sans jamais citer le texte soumis. \textbf{Sanction :} Hors-sujet grave (note souvent < 4/20). Le texte est le \emph{seul} objet d'étude ; vos connaissances ne servent qu'à l'éclairer, jamais à le remplacer.
\end{attention}

\courssub{Vue d'ensemble des cinq étapes de la démarche}

La méthode n'est pas une cage, mais une \cle{boussole}. Voici les cinq marches de l'escalier : on ne monte pas la troisième sans avoir solidement posé le pied sur la deuxième.

\begin{methode}[Les 5 étapes du commentaire philosophique]
\begin{enumerate}
    \item \textbf{Lecture analytique (brouillon) :} 3 lectures minimum. (1) Globale : quel sujet ? quelle thèse ? (2) Détaillée : repérer la structure (arguments, connecteurs, concepts-clés), annoter les marges. (3) Critique : tester la cohérence, noter les questions/objections.
    \item \textbf{Problématisation :} Formuler la \cle{problématique} : la tension philosophique que le texte résout (ex: « Comment la technique, moyen de libération, devient-elle source d'asservissement ? »).
    \item \textbf{Plan détaillé :} Construire l'architecture (2 ou 3 parties principales, 2 sous-parties chacune). Chaque partie = une étape de la démonstration de l'auteur + sa discussion. Rédiger les \cle{phrases d'introduction} et de \cle{transition} de chaque partie.
    \item \textbf{Rédaction :} \emph{Introduction} (accroche, présentation texte, thèse, problématique, annonce du plan). \emph{Développement} (Explication structurée + Discussion intégrée ou séparée). \emph{Conclusion} (bilan, ouverture).
    \item \textbf{Relecture active :} Vérifier : fidélité aux citations, cohérence du fil directeur, orthographe/grammaire, respect du temps imparti.
\end{enumerate}
\end{methode}

\begin{center}
\begin{popfigure}
\begin{tikzpicture}[
    stepnode/.style={rectangle, draw=popdark, thick, rounded corners, minimum width=3.2cm, minimum height=1.5cm, align=center, font=\small\bfseries, fill=popblueL, text=popdark},
    arrownode/.style={font=\Large, text=poporange},
    numnode/.style={circle, fill=popblue, text=white, font=\bfseries, minimum size=0.7cm, inner sep=0}
]
% Nodes
\node[stepnode, align=center] (s1){\textbf{1. LECTURE ANALYTIQUE}\\\emph{Annoter, structurer, questionner}};
\node[numnode, above left=-2pt and -2pt of s1.north west] {1};
\node[stepnode, right=1.5cm of s1, align=center] (s2){\textbf{2. PROBLÉMATISATION}\\\emph{Formuler la tension philosophique}};
\node[numnode, above left=-2pt and -2pt of s2.north west] {2};
\node[stepnode, right=1.5cm of s2, align=center] (s3){\textbf{3. PLAN DÉTAILLÉ}\\\emph{Architecture : parties / sous-parties}};
\node[numnode, above left=-2pt and -2pt of s3.north west] {3};
\node[stepnode, right=1.5cm of s3, align=center] (s4){\textbf{4. RÉDACTION}\\\emph{Intro / Développement / Conclusion}};
\node[numnode, above left=-2pt and -2pt of s4.north west] {4};
\node[stepnode, right=1.5cm of s4, align=center] (s5){\textbf{5. RELECTURE}\\\emph{Contrôle fidélité, cohérence, langue}};
\node[numnode, above left=-2pt and -2pt of s5.north west] {5};

% Arrows (staircase effect)
\draw[->, very thick, poporange] (s1.east) -- (s2.west) node[midway, above, font=\tiny, text=popdark] {Thèse + Structure};
\draw[->, very thick, poporange] (s2.east) -- (s3.west) node[midway, above, font=\tiny, text=popdark] {Question directrice};
\draw[->, very thick, poporange] (s3.east) -- (s4.west) node[midway, above, font=\tiny, text=popdark] {Squelette complet};
\draw[->, very thick, poporange] (s4.east) -- (s5.west) node[midway, above, font=\tiny, text=popdark] {Texte final};

% Global label
\node[above=0.5cm of s3, font=\bfseries\large, text=popdark] {L'escalier du commentaire : chaque marche conditionne la suivante};
\end{tikzpicture}
\legende{Les 5 étapes séquentielles. Sauter l'étape 1 ou 2 rend l'étape 4 (rédaction) hasardeuse et souvent hors-sujet.}
\end{popfigure}
\end{center}

\courssub{Critères d'évaluation au Baccalauréat technique}

Le correcteur dispose d'une grille nationale. Connaître les items, c'est viser les points.

\begin{center}
\begin{tabularx}{\linewidth}{|l|X|>{\centering\arraybackslash}p{2.5cm}|}
\hline
\rowcolor{popblueL}
\textbf{Critère} & \textbf{Indicateurs de réussite} & \textbf{Poids indicatif} \\
\hline
\textbf{Compréhension du texte} & Thèse identifiée exactement ; structure (arguments/objections) restituée ; concepts clés définis dans le contexte du texte. & $\sim$ 4 pts \\
\hline
\textbf{Cohérence du plan} & Plan en 2 ou 3 parties logiques ; chaque partie correspond à un moment de la démonstration de l'auteur ; transitions fluides ; annonces tenues. & $\sim$ 3 pts \\
\hline
\textbf{Qualité de l'explication} & Citations précises (lignes/guillemets) ; analyse des connecteurs logiques ; mise en lumière des présupposés ; pas de paraphrase. & $\sim$ 5 pts \\
\hline
\textbf{Qualité de la discussion} & Objections pertinentes (auteurs, contre-exemples concrets) ; nuances apportées à la thèse ; pas d'opinions gratuites ; retour au texte. & $\sim$ 4 pts \\
\hline
\textbf{Correction de la langue} & Syntaxe correcte ; vocabulaire philosophique précis ; absence de fautes d'orthographe/grammaire majeures ; lisibilité. & $\sim$ 4 pts \\
\hline
\end{tabularx}
\end{center}

\begin{savaistu}[Le commentaire au Bac Technique Camerounais]
Au Probatoire et au Baccalauréat technique, l'épreuve de Philosophie propose \textbf{deux sujets au choix} : une dissertation et un commentaire de texte. Historiquement, le commentaire est souvent \emph{mieux réussi} par les candidats techniques car il offre un « filet de sécurité » : le texte est sous les yeux. Cependant, les correcteurs sanctionnent \textbf{systématiquement} deux écueils majeurs : le \textbf{hors-sujet} (ne pas répondre à la problématique implicite du texte) et la \textbf{paraphrase} (suivre le texte linéairement sans l'analyser). Une copie qui « explique bien le texte » mais ne le « discute » pas plafonne à 10/20. L'excellence (14+) exige l'articulation \textbf{Explication rigoureuse + Discussion argumentée}.
\end{savaistu}

\courssub{Situation concrète : pourquoi le jury rejette la paraphrase}

Imaginez un extrait de \emph{La Condition de l'homme moderne} de Hannah Arendt sur le \cle{monde} et la \cle{terre}. Le texte distingue l'\emph{animal laborans} (enfermé dans la nécessité biologique) et l'\emph{homo faber} (créateur d'un monde durable).

\begin{exoresolu}[Paraphrase vs Commentaire : le même paragraphe, deux traitements]
\textbf{Extrait (fictif, inspiré d'Arendt) :} « Le travail est l'activité qui correspond à la vie biologique de l'homme. Il ne produit rien de durable ; son fruit est consommé aussitôt. L'œuvre, au contraire, édifie un monde d'objets qui résistent au temps. »

\textbf{\textcolor{poppink}{Copie A (Paraphrase — Note : 5/20)} :}
« Arendt dit que le travail correspond à la vie biologique. Elle explique qu'il ne produit rien de durable car le fruit est consommé tout de suite. Ensuite, elle oppose l'œuvre qui, elle, construit un monde d'objets qui résistent au temps. En résumé, le travail c'est la vie et l'œuvre c'est le monde. »

\textit{Analyse :} L'élève suit le texte phrase par phrase (« dit que », « explique que », « oppose »). Aucune analyse de la \cle{distinction conceptuelle} (nécessité vs liberté, cyclique vs linéaire), aucun mot sur la \cle{portée} (condition de l'humanité), aucune discussion.

\tcblower

\textbf{\textcolor{popgreen}{Copie B (Commentaire — Note : 16/20)} :}
« \textbf{Thèse :} Arendt fonde une anthropologie philosophique en distinguant radicalement \emph{Labor} (travail) et \emph{Work} (œuvre).
\textbf{Explication :} Le \emph{Labor} (l. 1-2) est assigné à la \cle{sphère du privé} et à la \cle{nécessité} : il suit le cycle biologique (production/consommation immédiate), ne laissant aucune trace (ex: manger, nettoyer). L'\emph{Work} (l. 3) relève de la \cle{fabrication} : il projette une \cle{fin} (l'objet durable), transforme la matière (le bois, la pierre, le code informatique) et institue un \cle{monde commun} (la table, la route, l'application) qui survive à son créateur.
\textbf{Analyse structurelle :} L'opposition n'est pas descriptive mais \cle{normative} : seul l'\emph{Work} permet l'\cle{immortalité} des choses et la \cle{liberté} politique (agir ensemble dans un monde stable).
\textbf{Discussion :} Cette distinction résiste-t-elle à la \cle{technique moderne} ? La machine à coudre de l'atelier de Douala (exemple local) transforme le travail de couture (répétitif, biologique) en production d'œuvres (vêtements durables, vendus). La frontière s'effrite. De plus, le \emph{digital labor} (modération de contenu, data labeling) produit des « œuvres » (données) immatérielles mais précaires. Arendt sous-estime-t-elle la \cle{technicisation du travail} ? »

\textit{Analyse :} L'élève \textbf{reconstruit la logique} (thèse, concepts, articulation), \textbf{cite les lignes}, \textbf{illustre par un exemple camerounais concret} (machine à coudre, digital labor), \textbf{discute} la pertinence actuelle.
\end{exoresolu}

\begin{aretenir}[La règle d'or pour le Bac]
\textbf{« Suis la logique, pas la ligne. »} Ton plan ne doit \emph{jamais} calquer l'ordre des paragraphes du texte (linéaire). Ton plan doit suivre l'\emph{architecture argumentative} de l'auteur : 1. Ce qu'il pose (thèse/présupposés) $\rightarrow$ 2. Comment il le démontre (arguments/enchaînements) $\rightarrow$ 3. Ce que cela implique / où cela bute (portée/limites). C'est ainsi que tu passes de l'élève qui « donne son avis » au candidat qui « fait de la philosophie ».
\end{aretenir}

La maîtrise de ces fondamentaux acquis, nous entrons maintenant dans le vif du sujet : la \textbf{lecture analytique}, cette étape invisible sur la copie mais décisive dans la tête du correcteur.

\coursec{La lecture analytique du texte}

Après avoir rappelé ce qu'est un commentaire philosophique et pourquoi il ne se réduit pas à un résumé ni à une dissertation déguisée, nous entrons maintenant dans le vif du sujet : \textbf{la lecture analytique}. C'est l'étape fondatrice, celle où l'on « met le texte à nu » pour en comprendre l'architecture interne. Comme un chirurgien qui observe avant d'opérer, ou un ingénieur qui ausculte une structure avant de la renforcer, le commentateur commence par une lecture lente, rigoureuse, outillée. C'est à ce prix seulement que la problématisation et le développement pourront être pertinents.

\courssub{Première lecture : survol et identification du contexte}

La toute première approche ne vise pas le détail, mais \textbf{l'horizon de sens}. On lit d'une traite, sans s'arrêter sur les difficultés, pour saisir l'intention générale : de quoi l'auteur veut-il nous convaincre ? Quel est le ton (polémique, pédagogique, méditatif ?). Simultanément, on identifie les \textbf{indices paratextuels} : le nom de l'auteur, le titre de l'œuvre, la date de publication, la collection éventuelle. Ces indices permettent de situer le \textbf{courant philosophique} (rationalisme, empirisme, existentialisme, phénoménologie, école de Francfort, etc.) et l'\textbf{époque}, ce qui éclaire souvent le vocabulaire technique utilisé.

\begin{exemplebox}[Indices paratextuels]
Un extrait commence par : « \emph{La technique n'est pas un simple moyen...} » (Martin Heidegger, \emph{La question de la technique}, 1954).
\begin{itemize}
    \item \textbf{Auteur :} Martin Heidegger (1889--1976), philosophe allemand majeur du XXe siècle.
    \item \textbf{Courant :} Phénoménologie / Ontologie fondamentale.
    \item \textbf{Contexte :} Après-guerre, essor technologique massif, interrogations sur l'essence de la technique.
    \item \textbf{Attente lexicale :} Mots comme \cle{essence}, \cle{manifestation}, \cle{dévoilement}, \cle{Gestell} (dispositif/encadrement).
\end{itemize}
\end{exemplebox}

\courssub{Deuxième lecture active : repérage lexical et logique}

On reprend le texte, crayon en main (ou surligneur numérique). C'est la \textbf{lecture active}, opératoire. Trois gestes simultanés :

\begin{enumerate}
    \item \textbf{Souligner les \cle{mots-clés} (termes techniques, concepts du programme) :} \emph{liberté, déterminisme, conscience, aliénation, progrès, nature, culture, devoir, État...}
    \item \textbf{Entourer les \cle{connecteurs logiques} :} ce sont les « charnières » du raisonnement. Ils révèlent la structure argumentative sans même lire le contenu des phrases.
    \item \textbf{Noter en marge les questions/objections :} « Pourquoi cet exemple ? », « Ce terme change-t-il de sens ? », « Transition brutale ici. »
\end{enumerate}

\begin{definition}[Connecteur logique]
Un \textbf{connecteur logique} (ou opérateur argumentatif) est un mot ou une expression qui explicite la relation logique entre deux propositions ou deux parties du texte. Il guide le lecteur dans le cheminement de la pensée de l'auteur.
\end{definition}

\begin{center}
\begin{tabularx}{\linewidth}{|l|X|}\hline
\textbf{Type de relation} & \textbf{Principaux connecteurs} \\ \hline
Cause / Explication & \emph{car, car en effet, parce que, puisque, du fait que, étant donné que} \\ \hline
Conséquence / Conclusion & \emph{donc, par conséquent, c'est pourquoi, ainsi, alors, partant, il s'ensuit que} \\ \hline
Opposition / Concession & \emph{mais, cependant, pourtant, toutefois, néanmoins, au contraire, si... du moins} \\ \hline
Illustration / Exemple & \emph{par exemple, ainsi, notamment, c'est le cas de, prenons le cas de} \\ \hline
Énumération / Articulation & \emph{premierement, ensuite, enfin ; d'une part... d'autre part ; de plus, par ailleurs, furthermore} \\ \hline
Reformulation / Résumé & \emph{autrement dit, en somme, bref, cela signifie que, en un mot} \\ \hline
\end{tabularx}
\end{center}

\begin{popfigure}
\legende{Exemple de structure argumentative annotée : repérage des connecteurs (cerclés), mots-clés (soulignés) et découpage en moments (barres verticales).}
\begin{tikzpicture}[
    node distance=0.3cm and 0.5cm,
    connector/.style={circle, draw=poporange, thick, inner sep=2pt, fill=poporange!15, font=\small\bfseries},
    keyword/.style={underline, draw=popblue, thick, text=popdark},
    momentline/.style={dashed, poppurple, thick},
    momentlabel/.style={poppurple, font=\bfseries\small, fill=white, inner sep=2pt},
    textblock/.style={align=left, text width=0.85\linewidth, font=\small}
]
% Text lines as separate nodes for clean annotation
\node[textblock, anchor=north west, align=center] (l1) at (0,0){
\textbf{Extrait : La technique et l'homme (d'après H. Jonas)} \\
\underline{La technique} moderne se distingue \tikz[baseline=-0.5ex]\node[connector]{donc}; de l'artisanat traditionnel par sa \underline{puissance} démesurée.
};
\node[textblock, anchor=north west] (l2) at (0,-1.8) {
\tikz[baseline=-0.5ex]\node[connector]{En effet};, l'outil ancien prolongeait la \underline{main} de l'homme sans la remplacer ; la machine, elle, \tikz[baseline=-0.5ex]\node[connector]{car}; elle possède sa propre \underline{énergie}, s'autonomise. \tikz[baseline=-0.5ex]\node[connector]{Dès lors};, l'homme devient \underline{serviteur} de la machine plus que son maître.
};
\node[textblock, anchor=north west] (l3) at (0,-4.2) {
\tikz[baseline=-0.5ex]\node[connector]{Cependant};, cette \underline{aliénation} n'est pas une fatalité. \tikz[baseline=-0.5ex]\node[connector]{Car}; la technique reste une \underline{production humaine} : elle porte en elle une \underline{finalité} que l'homme a inscrite. \tikz[baseline=-0.5ex]\node[connector]{Il faut donc}; repenser la \underline{responsabilité} éthique face aux objets techniques.
};

% Moment lines on the left
\draw[momentline] (-0.8, 0.3) -- (-0.8, -2.0) node[midway, left=0.3cm, momentlabel] {Moment 1};
\draw[momentline] (-0.8, -2.0) -- (-0.8, -4.5) node[midway, left=0.3cm, momentlabel] {Moment 2};
\draw[momentline] (-0.8, -4.5) -- (-0.8, -6.5) node[midway, left=0.3cm, momentlabel] {Moment 3};

% Frame
\draw[popline, thick, rounded corners=5pt] (-1.2, 0.5) rectangle (0.85\linewidth+0.2, -6.8);
\end{tikzpicture}
\end{popfigure}

\courssub{Repérage du thème et dégagement de la thèse : distinction cruciale}

C'est le point nodal où beaucoup d'élèves trébuchent. Il faut séparer net le \textbf{thème} (le \emph{sujet} du texte) de la \textbf{thèse} (la \emph{position} de l'auteur sur ce sujet).

\begin{definition}[Thème]
Le \textbf{thème} est le sujet général traité dans le texte. Il se formule en \textbf{un mot ou une expression nominale} (ex. : « La technique et la nature », « La liberté face au déterminisme », « Le fondement du devoir »).
\end{definition}

\begin{definition}[Thèse]
La \textbf{thèse} est la \textbf{position précise, affirmative et défendable} que l'auteur soutient sur ce thème. Elle se formule en \textbf{une phrase complète, assertive} (ex. : « La technique moderne aliène l'homme en s'autonomisant, mais demeure une responsabilité humaine »).
\end{definition}

\begin{attention}[Piège classique : Confondre thème et thèse]
\begin{itemize}
    \item \textbf{Thème :} « La liberté. » (Trop vague, ne dit rien de la pensée de l'auteur).
    \item \textbf{Mauvaise thèse :} « L'auteur parle de la liberté. » (C'est une description du thème, pas la thèse).
    \item \textbf{Bonne thèse :} « L'auteur soutient que la liberté n'est pas l'absence de contrainte, mais l'obéissance à une loi qu'on se donne à soi-même (autonomie). »
\end{itemize}
\textbf{Astuce :} La thèse est \emph{contestable}. On peut être d'accord ou pas. Le thème, lui, est neutre.
\end{attention}

\begin{methode}[Trouver la thèse en 3 étapes]
\begin{enumerate}
    \item \textbf{Chercher la phrase « pivot » :} Souvent en début (annonce) ou en fin (conclusion) de texte, parfois au cœur d'un paragraphe charnière. C'est la phrase que \emph{tout le reste du texte cherche à établir}.
    \item \textbf{Tester la couverture :} Est-ce que \emph{tous} les arguments, exemples, distinctions du texte convergent vers cette phrase ? Si un pan du texte ne s'y rattache pas, la thèse est trop étroite ou mal formulée.
    \item \textbf{Formuler en « Je » (mentalement) :} « L'auteur veut me prouver que... » $\rightarrow$ transformer en phrase objective à la 3e personne.
\end{enumerate}
\end{methode}

\courssub{Analyse de la structure argumentative : arguments, exemples et découpage en moments}

Une fois la thèse posée, on reconstitue le \textbf{squelette logique} : comment l'auteur \emph{prouve}-t-il sa thèse ?

\begin{itemize}
    \item \textbf{Les arguments :} Raisons logiques, distinctions conceptuelles, réfutations d'objections, renvois à des principes.
    \item \textbf{Les exemples / illustrations :} Faits historiques, situations concrètes, expériences de pensée, références littéraires ou scientifiques. Ils \emph{ancrent} l'abstrait dans le concret.
    \item \textbf{Le découpage en \cle{moments} (unités de sens) :} Un moment = une étape unifiée du raisonnement (ex. : « Constat », « Analyse de la cause », « Ouverture éthique »). Les connecteurs logiques en sont les frontières naturelles.
\end{itemize}

\begin{aretenir}[Règle d'or du découpage]
Sur un texte de \textbf{15 à 20 lignes} (format Probatoire), un bon découpage compte \textbf{généralement 2 à 4 moments}. Chaque moment deviendra une \textbf{partie principale} (I, II, III...) du développement du commentaire. Moins de 2 moments = texte trop simple ou lecture superficielle ; plus de 4 = découpage trop fin, perte de l'unité argumentative.
\end{aretenir}

\begin{popfigure}
\legende{Carte mentale de la structure argumentative d'un texte philosophique (thème central, thèse, arguments, exemples, moments).}
\begin{tikzpicture}[
    mindmap,
    concept color=popblue,
    level 1 concept/.append style={level distance=4.5cm, sibling angle=90, concept color=popblue!70},
    level 2 concept/.append style={level distance=3.2cm, sibling angle=45, concept color=popgreen!70},
    level 3 concept/.append style={level distance=2.5cm, sibling angle=30, concept color=poporange!70},
    every node/.style={font=\small, text=white},
    text width=2.8cm, align=center
]
\node[concept, align=center]{Thème :\\ Technique \& Humanité}
    child[concept color=poppurple] { node[concept, align=center]{Thèse :\\ L'autonomisation technique\\ exige une responsabilité\\ éthique accrue}
        child { node[concept, align=center]{Moment 1 :\\ Constat de puissance\\ (Argument : énergie propre\\ Exemple : machine vs outil)} }
        child { node[concept, align=center]{Moment 2 :\\ Risque d'aliénation\\ (Argument : inversion maître/serviteur\\ Exemple : ouvrier à la chaîne)} }
        child { node[concept, align=center]{Moment 3 :\\ Réappropriation possible\\ (Argument : finalité humaine\\ Exemple : éthique de la responsabilité)} }
    };
\end{tikzpicture}
\end{popfigure}

\courssub{Repérage des notions du programme mobilisées}

Le texte n'est pas une île : il est ancré dans le \textbf{programme de Première Technique}. Pendant la lecture analytique, on dresse la liste des \cle{notions officielles} explicitement nommées ou implicitement travaillées : \cle{Conscience}, \cle{Liberté}, \cle{Technique}, \cle{Devoir}, \cle{État}, \cle{Société}, \cle{Culture}, \cle{Nature}, \cle{Travail}, \cle{Justice}, \cle{Bonheur}, \cle{Connaissance}, \cle{Vérité}, \cle{Religion}, \cle{Science}, \cle{Art}, \cle{Langage}, \cle{Histoire}, \cle{Éducation}.

\begin{exemplebox}[Lien texte / programme]
Dans l'extrait annoté plus haut (d'après Jonas) :
\begin{itemize}
    \item \textbf{Notion principale :} \cle{Technique} (essence, puissance, autonomie, finalité).
    \item \textbf{Notions connexes :} \cle{Liberté} (aliénation/autonomie), \cle{Devoir} / \cle{Responsabilité} (éthique), \cle{Travail} (ouvrier, machine), \cle{Nature} (énergie propre vs main d'homme).
    \item \textbf{Attente du correcteur :} Le commentaire doit montrer que l'élève maîtrise la définition philosophique de ces notions (ex. : technique $\neq$ technologie, responsabilité $\neq$ culpabilité).
\end{itemize}
\end{exemplebox}

\courssub{TP guidé : lecture analytique collective d'un extrait sur la technique}

\noindent \textbf{Consigne :} Complétez la grille ci-dessous pour l'extrait suivant (durée : 20 min individuel + 15 min mise en commun).

\begin{center}
\begin{tabularx}{\linewidth}{|l|X|}\hline
\textbf{Auteur / Œuvre / Date} & \textbf{Hans Jonas, \emph{Le principe responsabilité}, 1979 (extrait adapté)} \\ \hline
\textbf{Thème (1 mot / expr.)} & \\ \hline
\textbf{Thèse (1 phrase complète)} & \\ \hline
\textbf{Mots-clés soulignés (5--7)} & \\ \hline
\textbf{Connecteurs logiques encerclés (4--5)} & \\ \hline
\textbf{Découpage en moments (2--4)} & \textbf{Moment 1 (lignes ... à ...) :} \newline \textbf{Moment 2 (lignes ... à ...) :} \newline \textbf{Moment 3 (lignes ... à ...) :} \\ \hline
\textbf{Argument principal par moment} & \textbf{M1 :} \newline \textbf{M2 :} \newline \textbf{M3 :} \\ \hline
\textbf{Exemple / Illustration par moment} & \textbf{M1 :} \newline \textbf{M2 :} \newline \textbf{M3 :} \\ \hline
\textbf{Notions du programme repérées} & \\ \hline
\end{tabularx}
\end{center}

\vspace{0.5cm}
\noindent \textbf{Texte support (18 lignes) :}

\emph{« La nature, dans sa totalité, n'est plus un « don » que l'homme reçoit passivement ; elle est devenue un \underline{enjeu} de sa \underline{puissance} technique. \textbf{Car} la technique moderne, par sa capacité d'intervention à l'échelle planétaire et irréversible, a transformé la \underline{relation} homme-nature en une relation de \underline{domination} à haut risque. \textbf{Or}, cette domination n'est pas sans \underline{retour} : les effets différés et cumulés de l'action technique (changement climatique, érosion de la biodiversité, déchets nucléaires) menacent les \underline{conditions mêmes de vie} des générations futures. \textbf{Donc}, l'éthique traditionnelle, fondée sur la proximité spatiale et temporelle (voisin, contemporain), devient \underline{insuffisante}. \textbf{Il faut} un \underline{impératif heuristique} nouveau : « Agis de façon que les effets de ton action soient compatibles avec la permanence d'une vie authentiquement humaine sur terre. » \textbf{C'est pourquoi} la \underline{responsabilité} s'étend désormais à l'\underline{avenir lointain} et à la \underline{biosphère entière}. »}

\begin{exoresolu}[Grille de lecture analytique complétée (corrigé type)]
\textbf{Éléments d'identification et analyse globale :}
\begin{itemize}
    \item \textbf{Auteur / Œuvre / Date :} Hans Jonas, \emph{Le principe responsabilité}, 1979. Courant : Éthique de la responsabilité / Philosophie de la technique.
    \item \textbf{Thème :} \cle{Technique}, nature et responsabilité envers l'avenir.
    \item \textbf{Thèse :} \textbf{La puissance technique planétaire impose une extension inédite de la responsabilité éthique à l'avenir lointain et à la biosphère, via un nouvel impératif heuristique.}
\end{itemize}
\tcblower
\textbf{Analyse détaillée :}
\begin{itemize}
    \item \textbf{Mots-clés :} puissance technique, enjeu, domination, effets différés/irréversibles, conditions de vie, éthique traditionnelle, insuffisante, impératif heuristique, responsabilité, avenir, biosphère.
    \item \textbf{Connecteurs :} \textbf{Car} (cause), \textbf{Or} (articulation/constat), \textbf{Donc} (conséquence), \textbf{Il faut} (prescription), \textbf{C'est pourquoi} (conclusion finale).
    \item \textbf{Découpage (3 moments) :}
    \begin{enumerate}
        \item \textbf{Lignes 1--4 : Constat ontologique.} La technique change la relation homme-nature (don $\rightarrow$ enjeu/domination).
        \item \textbf{Lignes 4--7 : Constat éthique (retour de bâton).} Les effets irréversibles menacent la vie future $\rightarrow$ éthique de proximité obsolète.
        \item \textbf{Lignes 7--fin : Prescription normative.} Nouvel impératif (heuristique) : responsabilité élargie à l'avenir et à la biosphère.
    \end{enumerate}
    \item \textbf{Arguments :} M1 : Argument d'échelle (planétaire/irréversible). M2 : Argument des conséquences (effets différés, cumulés, menace existentielle). M3 : Argument normatif (inadéquation de l'ancien, nécessité du nouveau).
    \item \textbf{Exemples :} M1 : Intervention planétaire. M2 : Changement climatique, biodiversité, déchets nucléaires. M3 : Formulation de l'impératif heuristique.
    \item \textbf{Notions programme :} \cle{Technique} (puissance, irréversibilité), \cle{Nature} (biosphère, don/enjeu), \cle{Devoir}/\cle{Responsabilité} (impératif, extension temporelle), \cle{Avenir} (générations futures), \cle{Éthique} (traditionnelle vs nouvelle).
\end{itemize}
\end{exoresolu}

\begin{methode}[Bilan : Check-list de la lecture analytique réussie]
Avant de passer à la problématisation, vérifiez que vous avez :
\begin{enumerate}
    \item[$\checkmark$] Lu \textbf{deux fois} (survol + active).
    \item[$\checkmark$] Identifié \textbf{auteur, œuvre, courant, date}.
    \item[$\checkmark$] Souligné \textbf{mots-clés} et encerclé \textbf{connecteurs}.
    \item[$\checkmark$] Formulé le \textbf{thème} (expression nominale).
    \item[$\checkmark$] Formulé la \textbf{thèse} (phrase assertive, contestable, couvrante).
    \item[$\checkmark$] Découpé en \textbf{2 à 4 moments} logiques (arguments + exemples).
    \item[$\checkmark$] Listé les \textbf{notions du programme} mobilisées.
\end{enumerate}
\emph{Seule une grille complète et rigoureuse autorise le passage à l'introduction.}
\end{methode}

\coursec{La problématisation et l'introduction}

Après avoir maîtrisé la \emph{lecture analytique} (section précédente), vous possédez désormais la matière première : le thème, la thèse, les concepts, les arguments et la structure interne du texte. Mais ces éléments épars ne forment pas encore un commentaire. Il faut maintenant les unifier sous une tension unique : la \cle{problématique}. C'est elle qui transforme un inventaire d'idées en une démonstration philosophique rigoureuse. Comme le disait Kant, « les concepts sans intuitions sont vides, les intuitions sans concepts sont aveugles » ; de même, une analyse sans problématique est aveugle, une problématique sans analyse est vide.

\courssub{De la lecture analytique à la tension philosophique : naître de la problématique}

\textbf{Intuition : le texte comme réponse à une question cachée.}\par
Tout texte philosophique est une \emph{réponse} apportée par un auteur à une \emph{question} qui le tourmente, explicite ou implicite. Votre tâche n'est pas de répéter la réponse, mais de \emph{reconstruire la question} qui la rend nécessaire. Si le texte de Marx affirme que « le travail aliéné nie l'homme », la question n'est pas « Qu'est-ce que le travail ? » (trop large), mais : \emph{« Comment le travail, essence de l'homme, peut-il devenir sa propre négation ? »} ou \emph{« La technique libère-t-elle l'homme ou l'aliène-t-elle ? »}. C'est cette question, formulée avec précision, qui est la problématique.

\begin{definition}[Problématique]
Question philosophique centrale à laquelle le texte répond, formulée à partir de la tension interne au sujet (opposition, paradoxe, aporie). Elle exprime l'enjeu théorique ou pratique du texte et guide l'intégralité du commentaire : chaque partie du développement doit y apporter un élément de réponse.
\end{definition}

\textbf{Comment passer du thème à la problématique ?}\par
Le thème est un \emph{sujet large} (ex. : « La technique », « La liberté », « L'État »). La problématique est une \emph{question resserrée} qui fait surgir la difficulté.
\begin{enumerate}
    \item \textbf{Identifiez la thèse de l'auteur} (grâce à la lecture analytique).
    \item \textbf{Repérez la tension} : l'auteur s'oppose-t-il à une opinion commune (doxa) ? Résout-il un paradoxe ? Dépasse-t-il une alternative fausse ?
    \item \textbf{Formulez la question} qui rend cette thèse \emph{nécessaire} et \emph{non évidente}.
\end{enumerate}

\begin{exemplebox}[Exemple : La technique chez Ellul vs Heidegger]
\textbf{Texte support (imaginaire, style Bac) :} « La technique n'est pas un outil neutre ; elle devient un système autonome qui dicte ses lois à la société. L'efficacité devient la seule valeur. »
\begin{itemize}
    \item \textbf{Thème :} La technique.
    \item \textbf{Thèse :} La technique est un système totalitaire (autonomie, efficacité).
    \item \textbf{Tension :} Opposition entre la représentation courante (technique = outil, moyen neutre) et la thèse de l'auteur (technique = fin en soi, milieu).
    \item \textbf{Problématique forte :} \emph{« En quoi l'autonomie du système technique transforme-t-elle le rapport de l'homme à son œuvre, au point de faire de l'efficacité une fin sans finalité humaine ? »}
    \item \textbf{Problématique faible :} \emph{« Qu'est-ce que la technique ? »} (trop général, ne guide pas le commentaire de \emph{ce} texte précis).
\end{itemize}
\end{exemplebox}

\begin{popfigure}
\begin{tikzpicture}[scale=0.9, every node/.style={font=\small}]
    % Funnel shape using a path
    \def\funnelwidth{10}
    \def\funnelheight{12}
    
    % Background funnel
    \draw[popblueL, line width=2pt, fill=popblueL!10] 
        (-\funnelwidth/2, 0) -- (-\funnelwidth/4, \funnelheight/3) -- (\funnelwidth/4, \funnelheight/3) -- (\funnelwidth/2, 0) -- cycle;
    \draw[popblueL, line width=2pt, fill=popblueL!10] 
        (-\funnelwidth/4, \funnelheight/3) -- (-\funnelwidth/8, 2*\funnelheight/3) -- (\funnelwidth/8, 2*\funnelheight/3) -- (\funnelwidth/4, \funnelheight/3) -- cycle;
    \draw[popblueL, line width=2pt, fill=popblueL!10] 
        (-\funnelwidth/8, 2*\funnelheight/3) -- (-0.5, \funnelheight) -- (0.5, \funnelheight) -- (\funnelwidth/8, 2*\funnelheight/3) -- cycle;

    % Labels inside funnel
    \node[align=center, text width=9cm] at (0, 1.0) {\textbf{1. AMORCE (Large)} \\ Fait social, expérience vécue, idée reçue, citation d'actualité. \\ \textit{Ex. : « À Yaoundé, 95\% des lycéens possèdent un smartphone... »}};
    \node[align=center, text width=6cm] at (0, 4.5) {\textbf{2. PRÉSENTATION DU TEXTE} \\ Auteur, œuvre, date, thèse, concepts clés, structure. \\ \textit{« Dans \emph{La Technique ou l'Enjeu du siècle}, Ellul soutient que... »}};
    \node[align=center, text width=4cm, text=popdark] at (0, 8.0) {\textbf{3. PROBLÉMATIQUE (Étroit)} \\ Question unique, précise, philosophique. \\ \textit{« La technique est-elle devenue un milieu totalitaire ? »}};
    \node[align=center, text width=3cm, text=poppurple] at (0, 11.0) {\textbf{4. ANNONCE DU PLAN} \\ 2 ou 3 moments logiques. \\ \textit{« Nous verrons d'abord... ensuite... enfin... »}};

    % Arrows on side
    \draw[poporange, thick, ->] (-6, 1) -- (-6, 4);
    \draw[poporange, thick, ->] (-6, 4) -- (-6, 7);
    \draw[poporange, thick, ->] (-6, 7) -- (-6, 10);
    \node[poporange, left] at (-6.5, 2.5) {Général};
    \node[poporange, left] at (-6.5, 5.5) {Contextualiser};
    \node[poporange, left] at (-6.5, 8.5) {Problématiser};
    \node[poporange, left] at (-6.5, 11.5) {Structurer};
\end{tikzpicture}
\legende{Schéma en entonnoir : la structure de l'introduction. On part du large (amorce) pour aboutir au précis (problématique + plan).}
\end{popfigure}

\courssub{Les qualités d'une bonne problématique : le triptyque « Précise – Philosophique – Textuelle »}

Une problématique n'est pas un titre, ni un thème, ni une suite de questions. Elle se juge à trois critères indispensables.

\begin{propriete}[Critères de validité d'une problématique]
\begin{enumerate}
    \item \textbf{Précise :} Elle circonscrit le débat. Elle évite les termes vagues (« l'homme », « la société ») au profit des concepts opératoires du texte (\cle{aliénation}, \cle{technicité}, \cle{subjectivité}, \cle{contractualisme}).
    \item \textbf{Philosophique :} Elle porte sur un concept, une valeur, une condition de possibilité, un fondement, une limite. Elle ne demande pas « Comment ? » (technique) ni « Combien ? » (statistique), mais « Pourquoi ? », « En quel sens ? », « Jusqu'où ? », « Sous quelle condition ? ».
    \item \textbf{Directement liée au texte :} Elle doit être \emph{la} question à laquelle \emph{ce} texte répond. Si on change de texte (même auteur, même thème), la problématique doit changer.
\end{enumerate}
\end{propriete}

\begin{popfigure}
\begin{tikzpicture}[scale=0.85]
    % Tableau comparatif
    \def\colA{3.5cm}
    \def\colB{5.5cm}
    \def\colC{5.5cm}
    \def\rowh{1.6cm}
    
    % Header
    \fill[popblue] (0,0) rectangle (\colA+\colB+\colC, 1);
    \node[white, font=\bfseries] at (\colA/2, 0.5) {Critère};
    \node[white, font=\bfseries] at (\colA+\colB/2, 0.5) {Problématique FAIBLE (À éviter)};
    \node[white, font=\bfseries] at (\colA+\colB+\colC/2, 0.5) {Problématique FORTE (Modèle)};
    
    % Row 1 : Généralité
    \fill[popred!10] (0,-\rowh) rectangle (\colA+\colB+\colC, 0);
    \node[align=center] at (\colA/2, -\rowh/2) {Trop générale};
    \node[align=center, text=popred] at (\colA+\colB/2, -\rowh/2) {\textit{« Qu'est-ce que la liberté ? »}};
    \node[align=center, text=popgreen] at (\colA+\colB+\colC/2, -\rowh/2) {\textit{« La liberté s'identifie-t-elle à l'autonomie rationnelle ou suppose-t-elle l'absence de contrainte extérieure ? »}};
    
    % Row 2 : Hors texte
    \fill[poporange!10] (0,-2*\rowh) rectangle (\colA+\colB+\colC, -\rowh);
    \node[align=center] at (\colA/2, -1.5*\rowh) {Hors texte};
    \node[align=center, text=popred] at (\colA+\colB/2, -1.5*\rowh) {\textit{« La technique nous rend-elle heureux ? » (Texte sur l'aliénation ouvrière)}};
    \node[align=center, text=popgreen] at (\colA+\colB+\colC/2, -1.5*\rowh) {\textit{« Comment la division technique du travail transforme-t-elle l'ouvrier en appendice de la machine ? » (Marx, \emph{Le Capital})}};
    
    % Row 3 : Fermée / Binaire
    \fill[popred!10] (0,-3*\rowh) rectangle (\colA+\colB+\colC, -2*\rowh);
    \node[align=center] at (\colA/2, -2.5*\rowh) {Fermée (Oui/Non)};
    \node[align=center, text=popred] at (\colA+\colB/2, -2.5*\rowh) {\textit{« L'État est-il nécessaire ? »}};
    \node[align=center, text=popgreen] at (\colA+\colB+\colC/2, -2.5*\rowh) {\textit{« En quel sens l'État, né de la contrainte, devient-il le garant de la liberté civile ? » (Hobbes / Rousseau)}};
    
    % Row 4 : Liste de questions
    \fill[poporange!10] (0,-4*\rowh) rectangle (\colA+\colB+\colC, -3*\rowh);
    \node[align=center] at (\colA/2, -3.5*\rowh) {Liste de questions};
    \node[align=center, text=popred] at (\colA+\colB/2, -3.5*\rowh) {\textit{« Qu'est-ce que le bonheur ? Est-il accessible ? Comment l'atteindre ? »}};
    \node[align=center, text=popgreen] at (\colA+\colB+\colC/2, -3.5*\rowh) {\textit{« La poursuite du bonheur comme fin morale ne conduit-elle pas paradoxalement à le manquer ? » (Kant, \emph{Fondements})}};
    
    % Grid lines
    \draw[popline] (0,0) -- (0,-4*\rowh);
    \draw[popline] (\colA,0) -- (\colA,-4*\rowh);
    \draw[popline] (\colA+\colB,0) -- (\colA+\colB,-4*\rowh);
    \draw[popline] (\colA+\colB+\colC,0) -- (\colA+\colB+\colC,-4*\rowh);
    \foreach \i in {0,1,2,3,4} {
        \draw[popline] (0,-\i*\rowh) -- (\colA+\colB+\colC,-\i*\rowh);
    }
\end{tikzpicture}
\legende{Tableau comparatif : problématiques faibles vs fortes. La forte utilise les concepts du texte, ouvre une tension dialectique et exclut la réponse binaire.}
\end{popfigure}

\begin{attention}[Piège n°1 : La problématique « parapluie »]
Ne posez \textbf{jamais} une question trop vaste du type : \emph{« Qu'est-ce que la justice ? », « La technique est-elle bonne ou mauvaise ? », « L'homme est-il libre ? »}.
\begin{itemize}
    \item \textbf{Conséquence :} Le développement devient un cours général, hors texte. Le correcteur note : « Hors sujet / Sujet non problématisé ».
    \item \textbf{Remède :} Revenez au texte. Soulignez le concept clé (ex. \emph{contractualisme}, \emph{aliénation}, \emph{volonté générale}) et formulez : \emph{« Comment le contractualisme fonde-t-il la légitimité politique sans nier la liberté naturelle ? »} (Rousseau).
\end{itemize}
\end{attention}

\courssub{L'architecture de l'introduction : les quatre temps immuables}

L'introduction est la \emph{porte d'entrée} de votre copie. Elle se lit en 30 secondes. Elle doit être un \emph{cheminement logique} sans saut, sans répétition, sans trou. Elle suit un rituel codifié en quatre temps, dans cet ordre strict.

\begin{methode}[Construire l'introduction en 4 temps (Ordre obligatoire)]
\begin{enumerate}
    \item \textbf{L'Amorce (Accroche) :} 1 à 2 phrases. Partir d'un fait concret, d'une expérience commune, d'une actualité, d'une idée reçue (doxa) ou d'une citation courte \emph{liée au thème}. \textbf{Interdit :} « Depuis la nuit des temps... », « De tout temps les hommes se demandent... », définitions du dictionnaire.
    \item \textbf{La Présentation du texte :} Auteur (nom + qualité/époque), Titre de l'œuvre (en italique), Date, Thèse principale (en une phrase), Concepts clés, Structure du texte (2 ou 3 moments).
    \item \textbf{La Problématique :} La question unique, formulée distinctement (souvent sur une ligne séparée ou en italique). Elle émerge \emph{naturellement} de la présentation (l'auteur répond \emph{à cela}).
    \item \textbf{L'Annonce du plan :} Formulation claire des 2 ou 3 grandes parties du développement. Utilisez des connecteurs logiques : « Dans un premier temps, nous analyserons... / Dans un second temps, nous montrerons... / Enfin, nous discuterons... ». \textbf{Ne numérotez pas} « I, II, III » dans l'intro (réservé au développement), mais annoncez le \emph{contenu} des parties.
\end{enumerate}
\end{methode}

\textbf{L'amorce : l'art de l'ancrage concret (Contexte camerounais).}\par
L'amorce prouve que la philosophie ne flotte pas dans les nuages, qu'elle éclaire \emph{ici et maintenant}. Pour un texte sur la technique/aliénation :
\begin{exemplebox}[Amorces types selon le thème]
\begin{itemize}
    \item \textbf{Technique / Aliénation :} « Dans les rues de Douala ou de Yaoundé, le smartphone est devenu une prothèse cognitive : notification, géolocalisation, algorithmes dictent nos déplacements, nos achats, nos relations. L'outil technique a-t-il cessé d'être un moyen pour devenir notre milieu ? »
    \item \textbf{État / Loi :} « Lorsque le code de la route est respecté par peur de la caméra et non par civisme, la contrainte légale a-t-elle éduqué la liberté ou simplement contraint le corps ? »
    \item \textbf{Travail / Technique :} « Dans les plantations du Moungo ou les chantiers de Kribi, la machine impose son rythme à l'ouvrier. Le progrès technique libère-t-il l'homme de la peine ou l'enferme-t-il dans la répétition ? »
    \item \textbf{Conscience / Inconscient :} « Un étudiant rate son oral alors qu'il maîtrise son cours : un lapsus, un blanc. La conscience est-elle vraiment maîtresse en sa propre maison ? » (Freud).
\end{itemize}
\end{exemplebox}

\textbf{L'annonce du plan : clarté et logique.}\par
Le plan annoncé doit refléter la \emph{structure dialectique} du texte (Thèse – Antithèse – Synthèse / Analyse – Critique – Dépassement / Concept – Argumentation – Portée).
\begin{aretenir}[Formulation type de l'annonce du plan]
\begin{itemize}
    \item \textbf{Plan en 2 parties (Analyse / Discussion) :} « Nous expliquerons d'abord comment l'auteur démontre que [Thèse + arguments principaux]. Nous discuterons ensuite les limites de cette thèse en montrant que [Objection / Nuance / Portée]. »
    \item \textbf{Plan en 3 parties (Dialectique) :} « Nous verrons dans un premier temps que [Thèse 1]. Nous analyserons ensuite comment [Antithèse / Complexification]. Enfin, nous montrerons que [Synthèse / Dépassement de l'auteur]. »
\end{itemize}
\textbf{Astuce :} Les titres de vos parties du développement doivent être les \emph{mêmes mots-clés} que ceux annoncés ici.
\end{aretenir}

\begin{savaistu}[Le regard du correcteur]
Les correcteurs du Probatoire / Baccalauréat technique lisent \textbf{d'abord l'introduction} (et la conclusion). Une introduction réussie (problématique ciselée, plan clair, présentation précise) crée un \emph{a priori favorable} : le correcteur entre dans le développement en confiance, il suit le fil. Une introduction floue, hors texte, sans problématique ou avec un plan « catalogue » (I. Les arguments, II. Les exemples, III. Mon avis) conditionne une note médiocre, même si le développement contient de bonnes choses. \textbf{Soignez-la comme une vitrine.}
\end{savaistu}

\begin{exemplebox}[Introduction type complète (Modèle Bac — ~10 lignes)]
\textbf{Sujet :} Commentaire d'un extrait de \emph{La Condition de l'homme moderne} (Hannah Arendt, 1958) sur la distinction \emph{Labor / Work / Action}.

\textit{[AMORCE]} À l'heure où le « \emph{burn-out} » touche jusqu'aux jeunes diplômés camerounais, la frontière entre travail nécessaire et réalisation de soi s'effondre dans l'urgence productive.
\textit{[PRÉSENTATION]} Dans \emph{La Condition de l'homme moderne} (1958), Hannah Arendt, philosophe politique allemande, distingue trois activités fondamentales : le \emph{labor} (travail cyclique, survie biologique), l'\emph{work} (\emph{œuvre}, fabrication d'un monde durable) et l'\emph{action} (parole et acte politique, révélation de l'unicité). Elle soutient que la société moderne a réduit l'humain à \emph{animal laborans}, sacrifiant l'action au labor.
\textit{[PROBLÉMATIQUE]} \textbf{En quoi la prédominance du labor sur l'action dans la société moderne menace-t-elle la condition politique de l'homme ?}
\textit{[PLAN]} Nous analyserons d'abord comment Arendt fonde la hiérarchie des activités \emph{vita activa} par la critique de la réduction marxiste du politique au social. Nous examinerons ensuite le processus de « levée » du social au rang de préoccupation publique comme perte du monde commun. Nous discuterons enfin la pertinence de cette distinction face à l'économie numérique contemporaine.
\end{exemplebox}

\courssub{TP guidé : rédaction collective d'une introduction (Entraînement Bac)}

\textbf{Consigne :} À partir du texte suivant (extrait court, style Probatoire), rédigez collectivement une introduction complète respectant les 4 temps. Durée : 20 min.

\begin{exoresolu}[TP Guidé : Introduction sur le Travail et la Technique]
\textbf{Texte support :} « La machine n'est pas seulement un instrument que l'homme use pour transformer la nature ; elle devient le modèle selon lequel s'organise le travail humain. L'ouvrier s'adapte au rythme de la machine, non l'inverse. Le temps de travail devient du temps machine. » (Inspiré de Günther Anders, \emph{L'Obsolescence de l'homme}, 1956).

\textbf{Étape 1 : Analyse express (5 min) — Rappel Section 2}
\begin{itemize}
    \item \textbf{Thème :} Technique / Travail / Condition ouvrière.
    \item \textbf{Thèse :} La machine s'autonomise et soumet l'homme (inversion moyen/fin).
    \item \textbf{Concepts :} hétéronomie, rythme, aliénation, modèle machinique.
    \item \textbf{Structure :} 1. Définition naïve (instrument) / 2. Renversement (modèle) / 3. Conséquence (temps machine).
\end{itemize}

\textbf{Étape 2 : Problématisation (5 min)}
\textit{Tension :} Croyance courante : « L'homme est maître de la machine » vs Thèse : « La machine dicte sa loi à l'homme ».
$\rightarrow$ \textbf{Problématique retenue :} \emph{« Comment la machine, conçue comme moyen de libération par la technique, devient-elle le principe d'hétéronomie qui soumet le temps et le geste de l'ouvrier ? »}
\tcblower
\textbf{Étape 3 : Rédaction des 4 temps (10 min)}

\begin{center}
\begin{tabularx}{\linewidth}{|X|}
\hline
\textbf{INTRODUCTION RÉDIGÉE (Équipe 3 - Note : 18/20)} \\ \hline
\textit{[Amorce]} Dans les usines de textile de Douala-Bassa comme dans les centres d'appel de Yaoundé, le chronomètre et l'algorithme imposent une cadence que l'opérateur ne choisit pas. L'outil promis à la libération des corps devient le métronome de leurs gestes. \\
\textit{[Présentation]} C'est ce renversement que décrit Günther Anders, philosophe autrichien de la technique, dans \emph{L'Obsolescence de l'homme} (1956). Il montre que la machine cesse d'être un simple instrument au service d'une fin humaine pour devenir un « modèle » organisationnel : l'ouvrier doit s'adapter au rythme de la machine, transformant le temps de travail en « temps machine ». Les concepts clés sont l'hétéronomie, l'inversion moyen/fin et l'obsolescence de l'humain. \\
\textit{[Problématique]} \textbf{En quel sens l'autonomisation du système technique transforme-t-elle le travailleur en fonction d'un appareil dont il a perdu la maîtrise ?} \\
\textit{[Plan]} Nous expliquerons d'abord comment Anders analyse le passage de l'outil à l'appareil comme inversion de la finalité technique. Nous montrerons ensuite que cette soumission du geste au rythme machinique constitue une aliénation spécifique du temps vécu. Nous discuterons enfin si la « désobéissance technique » ou la réappropriation du temps reste possible face à l'obsolescence programmée de l'homme. \\
\hline
\end{tabularx}
\end{center}
\end{exoresolu}

\begin{corrige}[Analyse pédagogique de l'introduction type]
\begin{itemize}
    \item \textbf{Amorce :} Ancrage camerounais concret (usines, centres d'appel), image forte (métronome), lien direct avec la thèse.
    \item \textbf{Présentation :} Auteur, œuvre, date, thèse résumée en 1 phrase, concepts nommés, structure esquissée. Complet et synthétique.
    \item \textbf{Problématique :} Unique, philosophique (hétéronomie, finalité, maîtrise), issue de la tension texte/doxa. Parfaite.
    \item \textbf{Plan :} 3 parties logiques (Explication thèse / Analyse concept clé / Discussion portée). Verbes d'action (expliquer, montrer, discuter). Annonce le contenu, pas juste « I, II, III ».
    \item \textbf{Longueur :} ~11 lignes. Idéal.
\end{itemize}
\end{corrige}

\begin{attention}[Erreurs fréquentes relevées lors du TP]
\begin{enumerate}
    \item \textbf{Amorce hors sujet :} « Le travail c'est la santé » (proverbe) sans lien avec la machine. $\rightarrow$ L'amorce doit contenir le \emph{nœud du problème} (machine/rythme).
    \item \textbf{Présentation incomplète :} Oubli de la date, du titre en italique, ou de la thèse. $\rightarrow$ Check-list obligatoire : Auteur / Titre / Date / Thèse / Concepts / Structure.
    \item \textbf{Problématique « liste de courses » :} « Qu'est-ce que la machine ? L'homme est-il libre ? Le travail est-il aliénant ? » $\rightarrow$ Une SEULE question, point final.
    \item \textbf{Plan « Catalogue » :} « I. Les arguments de l'auteur. II. Mes exemples. III. Mon opinion. » $\rightarrow$ Le plan doit refléter la \emph{logique de la démonstration} (Thèse / Tension / Dépassement).
    \item \textbf{Confusion Intro/Développement :} Développer un argument dans l'intro. $\rightarrow$ L'intro \emph{annonce}, elle ne \emph{prouve} pas. 8-12 lignes max.
\end{enumerate}
\end{attention}

\begin{methode}[Check-list finale avant de passer au développement (Auto-évaluation)]
Cochez chaque case avant d'écrire « \textbf{Développement} » :
\begin{itemize}
    \item[$\square$] Mon amorce est concrète, brève (1-2 lignes) et liée au thème du texte.
    \item[$\square$] J'ai cité : Auteur, Titre (\emph{italique}), Date, Thèse (1 phrase), 2-3 concepts, Structure du texte.
    \item[$\square$] Ma problématique est UNE question, philosophique, précise, issue de la tension du texte.
    \item[$\square$] Mon plan annonce 2 ou 3 parties avec des verbes (Analyser, Expliquer, Montrer, Discuter, Nuancer) et le \emph{contenu} de chaque partie.
    \item[$\square$] L'ensemble tient en 8 à 12 lignes manuscrites (environ 150-200 mots).
    \item[$\square$] Il n'y a AUCUN argument développé, AUCUN exemple personnel, AUCUN « Je pense que » dans l'introduction.
\end{itemize}
\end{methode}

\textbf{Transition vers la Section 4.}\par
L'introduction est signée, le cap est fixé. La boussole (problématique) et la carte (plan) sont en main. Il reste à \emph{naviguer} : c'est l'objet de la Section 4 — \emph{Le développement : expliquer, illustrer, discuter} — où chaque partie annoncée devient un argument rigoureux, nourri du texte et prolongé par la réflexion.

\coursec{Le développement : expliquer, illustrer, discuter}

Dans la section précédente, nous avons appris à construire la problématique et à rédiger une introduction qui annonce clairement le plan du commentaire. L'introduction est une promesse : le développement est le moment où cette promesse est tenue. C'est le cœur de la copie, la partie la plus longue et la plus décisive pour la note. Nous allons maintenant apprendre à construire ce développement selon trois gestes complémentaires : \cle{expliquer} la pensée de l'auteur, l'\cle{illustrer} par des exemples, puis la \cle{discuter} avec rigueur.

\courssub{Le principe de construction : suivre l'ordre du texte}

\textbf{Une partie = un moment du texte}\par

Le principe fondamental du commentaire philosophique est simple : le texte de l'auteur est notre guide. Un texte philosophique n'est pas un amas d'idées : c'est un \emph{raisonnement ordonné}, qui avance par étapes, comme un technicien qui assemble une machine pièce par pièce. Chaque étape du raisonnement de l'auteur s'appelle un \cle{moment} du texte.

\begin{definition}[Moment du texte]
Un \cle{moment} du texte est une étape du raisonnement de l'auteur, repérable par une articulation logique (<< d'abord >>, << en effet >>, << cependant >>, << donc >>) ou par un changement d'idée. À chaque moment du texte correspond une partie du développement.
\end{definition}

Concrètement : si le texte comporte trois moments (par exemple : l'auteur énonce une thèse, puis la justifie, puis en tire une conséquence), le développement comportera trois parties, rédigées \emph{dans le même ordre}. On ne commence jamais par la fin du texte, et on ne mélange pas les moments.

\begin{exemplebox}[Repérer les moments d'un texte]
Soit l'extrait suivant (inspiré de Descartes) : << Ceux qui ne marchent que fort lentement peuvent avancer beaucoup davantage, s'ils suivent toujours le droit chemin, que ne font ceux qui courent et qui s'en éloignent. >> On y repère deux moments : (1) une comparaison entre la lenteur droite et la rapidité errante ; (2) une leçon implicite : en recherche de la vérité, la méthode importe plus que la vitesse. Le développement aura donc deux parties, dans cet ordre.
\end{exemplebox}

\begin{exemplebox}[Un développement équilibré en chiffres]
Au Probatoire et au Baccalauréat technique, un développement équilibré compte \textbf{2 à 3 parties} de \textbf{2 à 3 paragraphes} chacune, soit environ \textbf{les deux tiers de la copie}. Si votre copie fait 4 pages, le développement en occupe environ 2 pages et demie à 3 pages. Une partie d'un seul paragraphe est presque toujours insuffisante ; une partie de cinq paragraphes trahit un défaut de synthèse.
\end{exemplebox}

\courssub{La structure d'un paragraphe de commentaire}

Chaque paragraphe du développement obéit à une architecture stable, qu'il faut s'entraîner à reproduire jusqu'à ce qu'elle devienne un réflexe. Un bon paragraphe de commentaire enchaîne cinq temps :

\begin{enumerate}
\item \textbf{L'idée du moment} : on annonce, en une phrase, ce que l'auteur dit à cette étape du texte (avec une courte citation entre guillemets).
\item \textbf{L'explication des concepts} : on définit les notions clés employées par l'auteur.
\item \textbf{L'analyse des arguments} : on montre \emph{comment} l'auteur justifie son idée, on dégage la logique de son raisonnement.
\item \textbf{L'exemple} : on illustre par un cas concret, de préférence tiré de la vie courante.
\item \textbf{La transition} : on résume ce qui vient d'être établi et on annonce ce qui reste à examiner.
\end{enumerate}

\begin{popfigure}
\begin{center}
\begin{tikzpicture}[node distance=0.35cm,
bloc/.style={draw=popblue, fill=popblueL, rounded corners=2pt, text width=10.5cm, minimum height=0.85cm, align=left, font=\small, inner sep=5pt},
fleche/.style={-{Stealth[length=2.5mm]}, thick, poporange}]
\node[bloc] (b1) {\textbf{1. Idée du moment} : << citation courte >> + reformulation de ce que dit l'auteur};
\node[bloc, below=of b1] (b2) {\textbf{2. Explication} : définition des concepts, présupposés dégagés};
\node[bloc, below=of b2] (b3) {\textbf{3. Analyse des arguments} : la logique du raisonnement est mise au jour};
\node[bloc, below=of b3] (b4) {\textbf{4. Exemple} : cas concret (atelier, marché, vie scolaire, réseaux sociaux)};
\node[bloc, below=of b4] (b5) {\textbf{5. Transition} : bilan du moment + annonce de la suite du texte};
\draw[fleche] (b1) -- (b2);
\draw[fleche] (b2) -- (b3);
\draw[fleche] (b3) -- (b4);
\draw[fleche] (b4) -- (b5);
\end{tikzpicture}
\end{center}
\legende{Schéma de la structure d'un paragraphe type de commentaire philosophique : cinq temps chaînés, de l'idée du texte jusqu'à la transition.}
\end{popfigure}

\begin{attention}[Piège fréquent : le paragraphe << catalogue >>]
Beaucoup de copies alignent des citations sans les expliquer, ou accumulent des exemples sans les rattacher au texte. Un paragraphe qui ne contient ni définition de concept ni analyse d'argument n'est pas un paragraphe de commentaire : c'est un résumé. Chaque citation doit être \emph{travaillée}, jamais simplement recopiée.
\end{attention}

\courssub{Expliquer : définir, dégager les présupposés, montrer la logique}

\textbf{Définir les notions}\par

Expliquer, c'est d'abord rendre clair ce que l'auteur dit. La première tâche consiste à \textbf{définir les notions} employées dans le texte. Si l'auteur parle de << devoir >>, de << liberté >> ou de << bonheur >>, le commentateur doit préciser le sens exact que ces mots prennent \emph{dans ce texte}, car un même mot peut avoir des sens différents selon les auteurs. Définir, c'est répondre à la question : << que faut-il entendre ici par... ? >>

\textbf{Dégager les présupposés}\par

Mais un raisonnement ne repose pas seulement sur ce qui est dit : il repose aussi sur ce qui est \emph{tu}. Tout raisonnement s'appuie sur des idées admises sans démonstration.

\begin{definition}[Présupposé]
Un \cle{présupposé} est une idée admise sans démonstration, sur laquelle repose le raisonnement de l'auteur. Il n'est pas écrit dans le texte, mais le texte ne peut pas fonctionner sans lui. Dégager un présupposé, c'est rendre visible cette fondation cachée.
\end{definition}

\begin{exemplebox}[Repérer un présupposé]
Quand Descartes affirme que ceux qui << suivent le droit chemin >> avancent davantage que ceux qui courent, il présuppose sans le démontrer dans cette phrase qu'\emph{il existe un droit chemin}, c'est-à-dire une méthode correcte de recherche de la vérité, et qu'on peut le distinguer des mauvais chemins. C'est ce présupposé que tout le Discours de la méthode va ensuite justifier. Le repérer montre qu'on a compris la structure profonde du texte.
\end{exemplebox}

\textbf{Montrer la logique du raisonnement}\par

Enfin, expliquer c'est montrer \emph{comment} l'auteur passe d'une idée à l'autre. Il faut identifier le type de raisonnement à l'œuvre :

\begin{itemize}
\item \textbf{Raisonnement déductif} : l'auteur tire une conséquence nécessaire de principes posés (<< si... alors... donc... >>).
\item \textbf{Argument par analogie} : l'auteur compare deux situations (la marche et la recherche de la vérité chez Descartes).
\item \textbf{Argument par l'absurde} : l'auteur montre que la thèse adverse conduit à une contradiction.
\item \textbf{Appel à l'expérience} : l'auteur s'appuie sur un fait observable.
\end{itemize}

Nommer précisément le procédé logique (<< l'auteur procède ici par analogie >>) est une marque de qualité immédiatement repérée par le correcteur.

\courssub{Illustrer : la force de l'exemple concret}

Expliquer rend le texte clair ; illustrer le rend \emph{vivant}. L'exemple a trois fonctions : il vérifie que vous avez réellement compris (on ne peut illustrer que ce qu'on a compris), il rend la thèse sensible au lecteur, et il prépare la discussion en testant la thèse sur le réel.

Un bon exemple de commentaire possède trois qualités : il est \textbf{pertinent} (il illustre exactement l'idée du texte, pas une idée voisine), \textbf{précis} (il est situé, détaillé, pas vague) et \textbf{bref} (deux ou trois phrases, jamais un paragraphe entier).

\begin{exemplebox}[Des exemples tirés de la vie courante camerounaise]
\begin{itemize}
\item \emph{Atelier technique} : pour illustrer l'idée que l'outil transforme le travail humain, le tourneur-fraiseur de Douala qui, grâce à sa machine, produit en une heure ce qu'il façonnait en une journée à la main.
\item \emph{Marché} : pour illustrer la confiance comme fondement des échanges, la vendeuse du marché Mokolo qui accorde du crédit à ses clientes habituelles : sans confiance réciproque, aucun échange durable n'est possible.
\item \emph{Réseaux sociaux} : pour illustrer la distinction entre opinion et vérité, une rumeur partagée des milliers de fois sur WhatsApp ne devient pas vraie pour autant : le nombre de partages mesure la croyance, non la vérité.
\item \emph{Vie scolaire} : pour illustrer l'idée que la contrainte peut libérer, le règlement intérieur du lycée technique, qui limite certaines conduites mais rend possible l'apprentissage de tous.
\end{itemize}
\end{exemplebox}

\begin{savaistu}[Un exemple concret vaut mieux qu'une longue théorie abstraite]
Un exemple bien choisi éclaire instantanément une thèse entière. Prenons la machine à coudre du tailleur de Bafoussam : elle multiplie sa production par cinq ou six, ce qui illustre la puissance libératrice de la technique ; mais elle le rend en même temps dépendant de l'électricité, des pièces détachées importées et du réparateur, ce qui illustre la nouvelle servitude que la technique crée. En une seule image concrète, vous avez montré les deux faces de la thèse << la technique libère et asservit >>, là qu'un long discours abstrait resterait confus. Les correcteurs du Baccalauréat technique valorisent toujours l'exemple précis et maîtrisé.
\end{savaistu}

\courssub{Discuter : évaluer la pertinence et les limites de la thèse}

Le commentaire philosophique n'est pas une paraphrase admirative. Après avoir expliqué et illustré, il faut \cle{discuter}, c'est-à-dire évaluer la pensée de l'auteur. Discuter, ce n'est pas << être d'accord ou pas d'accord >> : c'est soumettre la thèse à trois questions précises :

\begin{enumerate}
\item \textbf{L'argument est-il convaincant ?} La preuve avancée suffit-elle à établir la conclusion ? Le raisonnement comporte-t-il un saut logique ?
\item \textbf{Un contre-exemple est-il possible ?} Peut-on trouver un cas réel où la thèse échoue ? (Par exemple, face à << la technique libère l'homme >>, le téléphone portable qui rend l'apprenti mécanicien distrait et dépendant constitue un contre-exemple partiel.)
\item \textbf{Les présupposés sont-ils solides ?} Ce que l'auteur admet sans démonstration est-il lui-même acceptable ?
\end{enumerate}

\begin{methode}[Discuter une thèse : l'ordre obligatoire]
\begin{enumerate}
\item \textbf{Exposer fidèlement} d'abord : présenter la thèse de l'auteur telle qu'il la soutient, avec ses arguments, sans la déformer ni l'affaiblir.
\item \textbf{Évaluer ensuite} seulement : peser la force des arguments, proposer éventuellement un contre-exemple, interroger les présupposés.
\item \textbf{Nuancer enfin} : une discussion aboutit rarement à un rejet total ; elle délimite le plus souvent le domaine de validité de la thèse (<< la thèse de l'auteur vaut pour..., mais elle se heurte à... >>).
\end{enumerate}
Jamais l'inverse : évaluer avant d'avoir exposé revient à critiquer une pensée qu'on n'a pas encore comprise.
\end{methode}

\begin{attention}[L'erreur la plus fréquente : l'opinion personnelle prématurée]
La faute la plus sanctionnée au commentaire consiste à donner son opinion personnelle dès le début (<< je pense que l'auteur a tort >>, << selon moi, la technique est mauvaise >>) au lieu d'expliquer d'abord la pensée de l'auteur. Le correcteur attend d'abord que vous fassiez parler le texte ; votre jugement critique n'a de valeur que s'il s'appuie sur une explication exacte. Règle d'or : \textbf{d'abord comprendre, ensuite juger}.
\end{attention}

\courssub{Les transitions : relier les parties}

Entre deux parties du développement, il faut une \cle{transition} : une ou deux phrases qui remplissent une double fonction.

\begin{definition}[Transition]
La transition est une articulation rédigée qui (1) \textbf{résume ce qui vient d'être établi} dans la partie qui s'achève et (2) \textbf{annonce ce qui reste à examiner}, en montrant le lien logique entre les deux moments du texte.
\end{definition}

\begin{exemplebox}[Modèles de transitions]
\begin{itemize}
\item << Nous avons vu que l'auteur établit d'abord que la méthode prime sur la rapidité. Mais cette exigence suppose qu'on sache reconnaître le droit chemin : c'est précisément ce que la suite du texte se propose de montrer. >>
\item << Ainsi, la confiance apparaît comme le fondement de l'échange. Reste à savoir si elle suffit à le garantir : c'est la question que l'auteur aborde dans le dernier moment du texte. >>
\end{itemize}
Une bonne transition cite le fil conducteur : elle montre que le texte progresse, et que votre commentaire progresse avec lui.
\end{exemplebox}

\courssub{Commentaire suivi ou commentaire thématique ?}

Il existe deux manières d'organiser le développement :

\begin{itemize}
\item Le \cle{commentaire suivi} (ou linéaire) respecte l'ordre du texte : il suit l'auteur pas à pas, moment après moment, du début à la fin de l'extrait.
\item Le \cle{commentaire thématique} regroupe les idées par thèmes, en rapprochant des passages éloignés du texte qui traitent du même concept, quitte à bouleverser l'ordre de l'extrait.
\end{itemize}

\begin{popfigure}
\begin{center}
\begin{tikzpicture}[font=\small,
titre/.style={font=\bfseries},
bloc/.style={draw, rounded corners=2pt, text width=2.1cm, minimum height=0.7cm, align=center, inner sep=3pt},
fleche/.style={-{Stealth[length=2mm]}, thick}]
% Commentaire suivi
\node[titre, popblue] at (2.6,3.6) {Commentaire suivi};
\node[bloc, draw=popblue, fill=popblueL, align=center] (s1) at (0.6,2.6){Moment 1\\(début du texte)};
\node[bloc, draw=popblue, fill=popblueL, align=center] (s2) at (3.4,2.6){Moment 2\\(milieu)};
\node[bloc, draw=popblue, fill=popblueL, align=center] (s3) at (6.2,2.6){Moment 3\\(fin du texte)};
\draw[fleche, popblue] (s1) -- (s2);
\draw[fleche, popblue] (s2) -- (s3);
\node[text width=7cm, align=center, font=\footnotesize\itshape, popdark] at (3.4,1.6) {On suit l'ordre de l'auteur :\\fidélité maximale au texte};
% Commentaire thématique
\node[titre, poporange] at (10.4,3.6) {Commentaire thématique};
\node[bloc, draw=poporange, fill=poporangeL, text width=2.4cm] (t1) at (8.6,2.6) {Thème A\\(passages 1 et 3)};
\node[bloc, draw=poporange, fill=poporangeL, text width=2.4cm] (t2) at (11.6,2.6) {Thème B\\(passages 2 et 4)};
\draw[fleche, poporange] (t1) -- (t2);
\node[text width=5.4cm, align=center, font=\footnotesize\itshape, popdark] at (10.1,1.6) {On regroupe par idées :\\risque de perdre le fil du texte};
% Verdict
\node[draw=popgreen, fill=popgreenL, rounded corners=3pt, text width=12.6cm, align=center, font=\small\bfseries, inner sep=5pt] at (6.6,0.3) {Au Probatoire et au Baccalauréat technique : le commentaire suivi est recommandé};
\end{tikzpicture}
\end{center}
\legende{Comparaison des deux organisations du développement. Le commentaire suivi respecte l'ordre du raisonnement de l'auteur ; le commentaire thématique regroupe les passages par idées. Le commentaire suivi est recommandé à l'examen.}
\end{popfigure}

\begin{aretenir}[Pourquoi le commentaire suivi est recommandé au Bac technique]
Le commentaire suivi garantit trois choses : la \textbf{fidélité} au texte (on ne risque pas de déformer la pensée de l'auteur), la \textbf{couverture complète} de l'extrait (aucun moment n'est oublié) et la \textbf{clarté} du plan (le correcteur suit sans effort). Le commentaire thématique exige une maîtrise très grande du texte et expose à deux dangers : ignorer des passages et imposer au texte des catégories étrangères. Au Probatoire et au Baccalauréat technique, choisissez toujours le commentaire suivi.
\end{aretenir}

\courssub{TP guidé : rédaction d'un paragraphe de développement}

Mettons immédiatement en pratique la structure étudiée, à partir d'un moment de texte.

\begin{exoresolu}[Rédiger un paragraphe de commentaire complet]
\textbf{Énoncé.} Voici le premier moment d'un texte inspiré d'Épicure : << Le plaisir est le commencement et la fin de la vie heureuse. C'est lui que nous reconnaissons comme le premier bien, car c'est de lui que nous partons pour tout choix et tout rejet. >> Rédigez un paragraphe de développement complet (idée, explication, analyse, exemple, transition) commentant ce moment.

\tcblower
\textbf{Solution détaillée.} On suit les cinq temps du paragraphe type.

\emph{1. Idée du moment.} Dans ce premier moment, Épicure pose sa thèse fondatrice : << le plaisir est le commencement et la fin de la vie heureuse >>. Autrement dit, le plaisir est à la fois le point de départ et le but de toute existence réussie.

\emph{2. Explication des concepts.} Il faut entendre ici par \emph{plaisir} non pas n'importe quelle sensation agréable, mais ce que l'auteur reconnaît comme << le premier bien >> : un bien originel, qui n'a pas besoin d'être justifié par autre chose. Le présupposé du raisonnement est que tout choix humain suppose un critère, et que ce critère ne peut être trouvé que dans ce que nous éprouvons nous-mêmes.

\emph{3. Analyse des arguments.} L'auteur procède par un constat d'expérience : << c'est de lui que nous partons pour tout choix et tout rejet >>. Il observe que, de fait, nous choisissons ce qui nous procure du plaisir et rejetons ce qui nous fait souffrir ; le plaisir fonctionne donc comme la boussole naturelle de nos décisions.

\emph{4. Exemple.} Ainsi, l'élève d'un lycée technique de Yaoundé qui choisit la filière électrotechnique parce qu'il éprouve du plaisir à comprendre les circuits illustre ce propos : son choix d'orientation part bien d'un plaisir ressenti, avant tout calcul rationnel.

\emph{5. Transition.} Le plaisir apparaît donc comme le critère premier de nos choix. Mais tous les plaisirs se valent-ils ? C'est la question à laquelle l'auteur répond dans le moment suivant du texte, qu'il nous faut maintenant examiner.
\end{exoresolu}

\begin{exercice}[Entraînement : correction collective]
Rédigez un paragraphe de développement (10 à 15 lignes) commentant ce moment de texte inspiré de Kant : << Agis de telle sorte que tu traites l'humanité, en ta personne comme en celle de tout autre, toujours en même temps comme une fin, et jamais simplement comme un moyen. >> Consignes : annoncez l'idée avec une courte citation, définissez les notions de \emph{fin} et de \emph{moyen}, dégagez le présupposé, illustrez par un exemple de la vie courante camerounaise, puis rédigez une transition. La correction se fera collectivement au tableau, en comparant plusieurs productions.
\end{exercice}

\begin{corrige}[Exercice : éléments de correction attendus]
\emph{Idée} : Kant énonce un impératif moral fondamental, la formule de l'humanité. \emph{Définitions} : traiter quelqu'un comme une \emph{fin}, c'est le respecter pour lui-même ; le traiter comme un simple \emph{moyen}, c'est l'utiliser comme un instrument au service de nos intérêts. \emph{Présupposé} : tout être humain possède une dignité, c'est-à-dire une valeur qui ne se mesure pas à son utilité. \emph{Exemple} : l'employeur d'un atelier de soudure à Bafoussam qui paie correctement son apprenti et veille à sa sécurité le traite comme une fin ; celui qui l'exploite sans protection ni salaire le réduit à un moyen. \emph{Transition} : le respect de la personne apparaît ainsi comme la condition de toute morale ; reste à savoir si ce principe peut s'appliquer à toutes les situations concrètes, ce que la suite du texte devra préciser.
\end{corrige}

\begin{aretenir}[L'essentiel du développement]
\begin{itemize}
\item Une partie du développement correspond à un moment du texte, dans l'ordre du raisonnement de l'auteur.
\item Chaque paragraphe suit cinq temps : idée, explication, analyse de l'argument, exemple, transition.
\item Expliquer = définir les notions, dégager les présupposés, montrer la logique.
\item Illustrer = un exemple précis, pertinent et bref, tiré si possible de la vie courante.
\item Discuter = évaluer après avoir exposé, jamais l'inverse.
\item Au Baccalauréat technique, on pratique le commentaire suivi.
\end{itemize}
\end{aretenir}

Le développement ainsi construit, il restera à apprendre, dans la prochaine section, à conclure le commentaire et à relire efficacement la copie avant de la rendre.

\coursec{La conclusion et la relecture}

\courssub{Transition depuis le développement}

Après avoir conduit le développement — expliquer le texte pas à pas, l'illustrer d'exemples concrets, en discuter les thèses et les limites — vous arrivez au terme de votre devoir. Mais le travail n'est pas fini : la conclusion et la relecture sont les deux dernières marches qui transforment une copie correcte en une copie excellente. Rappelons-le : en commentaire philosophique, \textbf{la conclusion n'est pas un résumé}, elle est la \textbf{réponse explicite à la problématique} que vous avez posée en introduction. La relecture, quant à elle, garantit la cohérence globale et la qualité de l'expression. Cette section vous donne les clés pour les maîtriser.

\begin{popfigure}
\begin{tikzpicture}[scale=0.9, every node/.style={font=\small}]
  % Entonnoir inversé : bilan -> réponse -> ouverture
  \def\wtop{10}
  \def\wmid{6}
  \def\wbot{2}
  \def\h{2.2}
  \def\gap{0.3}
  
  % Forme entonnoir (trois trapèzes empilés)
  % Niveau 1 : Bilan (large)
  \fill[popblueL] (-\wtop/2, 3*\h+2*\gap) -- (\wtop/2, 3*\h+2*\gap) 
    -- (\wmid/2, 2*\h+2*\gap) -- (-\wmid/2, 2*\h+2*\gap) -- cycle;
  \draw[popblue, thick] (-\wtop/2, 3*\h+2*\gap) -- (\wtop/2, 3*\h+2*\gap) 
    -- (\wmid/2, 2*\h+2*\gap) -- (-\wmid/2, 2*\h+2*\gap) -- cycle;
  \node[text width=\wtop*0.85cm, align=center, font=\bfseries] at (0, 2.5*\h+2*\gap) {1. BILAN DE LA DÉMARCHE};
  \node[text width=\wtop*0.85cm, align=center] at (0, 2.2*\h+2*\gap) {Rappel synthétique des grandes étapes du développement (thèses explicitées, discussion, exemples)};
  
  % Flèche de transition
  \draw[popblue, thick, ->] (0, 2*\h+2*\gap) -- (0, 2*\h+\gap);
  
  % Niveau 2 : Réponse à la problématique (moyen)
  \fill[popgreenL] (-\wmid/2, 2*\h+\gap) -- (\wmid/2, 2*\h+\gap) 
    -- (\wbot/2, \h+\gap) -- (-\wbot/2, \h+\gap) -- cycle;
  \draw[popgreen, thick] (-\wmid/2, 2*\h+\gap) -- (\wmid/2, 2*\h+\gap) 
    -- (\wbot/2, \h+\gap) -- (-\wbot/2, \h+\gap) -- cycle;
  \node[text width=\wmid*0.85cm, align=center, font=\bfseries] at (0, 1.5*\h+\gap) {2. RÉPONSE À LA PROBLÉMATIQUE};
  \node[text width=\wmid*0.85cm, align=center] at (0, 1.2*\h+\gap) {Phrase claire, directe, qui reprend les termes de la question posée en introduction};
  
  % Flèche de transition
  \draw[popgreen, thick, ->] (0, \h+\gap) -- (0, \h);
  
  % Niveau 3 : Ouverture (étroit)
  \fill[poporangeL] (-\wbot/2, \h) -- (\wbot/2, \h) 
    -- (0, 0) -- cycle;
  \draw[poporange, thick] (-\wbot/2, \h) -- (\wbot/2, \h) -- (0, 0) -- cycle;
  \node[text width=\wbot*1.2cm, align=center, font=\bfseries] at (0, 0.5*\h) {3. OUVERTURE};
  \node[text width=\wbot*1.2cm, align=center] at (0, 0.2*\h) {Nouvelle question philosophique prolongeant le thème};
  
  % Légendes latérales
  \node[left, text width=3.5cm, align=right, font=\itshape\small] at (-6, 2.5*\h+2*\gap) {Large : synthèse};
  \node[left, text width=3.5cm, align=right, font=\itshape\small] at (-6, 1.5*\h+\gap) {Moyen : trancher};
  \node[left, text width=3.5cm, align=right, font=\itshape\small] at (-6, 0.5*\h) {Étroit : ouvrir};
\end{tikzpicture}
\legende{Schéma de la conclusion en entonnoir inversé : du bilan large à la réponse précise, puis à l'ouverture ciblée.}
\end{popfigure}

\courssub{La fonction de la conclusion : répondre à la problématique}

\begin{definition}[Fonction de la conclusion]
La conclusion est le moment où le candidat \textbf{répond explicitement à la problématique} posée en introduction, à la lumière de l'analyse du texte et de la discussion menée dans le développement. Elle ne résume pas : elle \textbf{tranche}.
\end{definition}

Intuitivement, imaginez un procès : l'introduction pose la question au jury (« Le texte de Descartes prouve-t-il l'existence de Dieu ? »), le développement instruit le dossier (arguments pour, arguments contre, nuances), la conclusion rend le verdict (« Oui, mais sous la condition que l'on accepte le critère de clarté et distinction »). Sans verdict, le procès n'a pas de sens. De même, une copie sans réponse finale laisse le correcteur sur sa faim.

\courssub{Structure de la conclusion : trois temps obligatoires}

Une conclusion réussie suit invariablement la structure en trois temps (l'entonnoir inversé ci-dessus) :

\begin{enumerate}
  \item \textbf{Le bilan de la démarche (2--3 lignes)} : rappel synthétique du chemin parcouru. On ne répète pas le plan mot pour mot, on en donne la logique : « Nous avons vu que le texte soutient X, puis que cette thèse se heurte à l'objection Y, ce qui nous a conduits à nuancer... »
  \item \textbf{La réponse à la problématique (1--2 lignes, le cœur)} : une phrase nette, qui reprend les termes de la question. Exemple : « Ainsi, le texte de Kant montre que la morale ne peut être fondée sur le bonheur, mais sur le devoir seul. »
  \item \textbf{L'ouverture (1--2 lignes)} : une question philosophique nouvelle, liée au thème, qui prolonge la réflexion sans la clore.
\end{enumerate}

\begin{methode}[Rédiger la conclusion en 3 étapes]
\begin{enumerate}
  \item \textbf{Relisez votre introduction} et soulignez la problématique exacte.
  \item \textbf{Écrivez la réponse en une phrase}, en reprenant le vocabulaire de la problématique.
  \item \textbf{Précédez-la du bilan} (ce que le développement a établi) et \textbf{suivez-la de l'ouverture} (question nouvelle).
\end{enumerate}
\end{methode}

\courssub{L'ouverture : une question nouvelle, pas un avis personnel}

\begin{definition}[Ouverture philosophique]
L'ouverture est une \textbf{question philosophique nouvelle}, en lien direct avec le thème du texte, qui \textbf{prolonge la réflexion} sans prétendre la clore. Elle montre que le problème traité dépasse le cadre du texte et invite à d'autres analyses.
\end{definition}

\begin{attention}[Pièges fréquents sur l'ouverture]
\begin{itemize}
  \item \textbf{Avis personnel} : « Je pense que la liberté est importante. » $\rightarrow$ \emph{Hors sujet : le commentaire n'est pas une dissertation d'opinion.}
  \item \textbf{Idée totalement nouvelle} : introduire un auteur ou un concept absent du développement. $\rightarrow$ \emph{Incohérence : l'ouverture doit naître de ce qui précède.}
  \item \textbf{Question fermée ou factuelle} : « Quand Kant a-t-il écrit ce texte ? » $\rightarrow$ \emph{Pas philosophique.}
  \item \textbf{Formule toute faite} : « La philosophie ne finit jamais de questionner. » $\rightarrow$ \emph{Vide de sens.}
\end{itemize}
\end{attention}

\begin{exemplebox}[Ouvertures réussies vs ratées]
\textbf{Texte :} Extrait de \emph{L'Existence précède l'essence} (Sartre). \textbf{Problématique :} « L'homme est-il condamné à être libre ? »

\begin{center}
\begin{tabularx}{\linewidth}{|l|X|}\hline
\textbf{Ratée} & \textbf{Réussie} \\ \hline
« Chacun a son avis sur la liberté. » & « Si l'existence précède l'essence, la responsabilité de l'homme est-elle sans limite, ou existe-t-il des situations où la liberté est concrètement impossible ? » \\ \hline
« La liberté, c'est important dans la vie. » & « La condamnation à la liberté ne risque-t-elle pas de nier l'influence des conditions sociales et historiques sur nos choix ? » \\ \hline
« Est-ce que Dieu existe ? » & « La liberté radicale sartrienne peut-elle fonder une éthique universelle, ou conduit-elle au relativisme ? » \\ \hline
\end{tabularx}
\end{center}
\end{exemplebox}

\courssub{Justifier sa conclusion : montrer la filiation avec l'analyse}

La conclusion ne tombe pas du ciel. Elle doit \textbf{découler visiblement} de ce qui a été écrit dans le développement. Le correcteur doit pouvoir tracer un fil continu : introduction $\rightarrow$ développement $\rightarrow$ conclusion.

\begin{propriete}[Critères de justification de la conclusion]
Une conclusion est justifiée si :
\begin{enumerate}
  \item Elle \textbf{reprend les termes de la problématique} (mots-clés identiques ou synonymes précis).
  \item Elle \textbf{s'appuie sur les résultats de la discussion} : « Puisque l'objection de X a montré que... », « Au terme de l'analyse des deux thèses... ».
  \item Elle \textbf{ne contredit pas le développement} : pas de revirement brutal non argumenté.
\end{enumerate}
\end{propriete}

\begin{exemplebox}[Conclusion justifiée vs non justifiée]
\textbf{Problématique :} « Le langage trahit-il la pensée ? » (texte de Bergson).

\textbf{Conclusion non justifiée :} « En conclusion, le langage ne trahit pas la pensée, il la révèle. C'est important pour communiquer. \textit{Ouverture :} Le langage des animaux est-il une langue ? »

\textbf{Conclusion justifiée :} « Au terme de cette analyse, nous pouvons répondre que le langage, selon Bergson, \textbf{trahit la pensée} en la fixant dans des concepts généraux qui en gomment la fluidité vivante — mais cette trahison est \textbf{nécessaire} pour la communication et l'action. \textit{Ouverture :} Cette nécessaire imperfection du langage ne fait-elle pas de toute philosophie une traduction approximative de l'intuition immédiate ? »
\end{exemplebox}

\courssub{La relecture : une étape stratégique (5 minutes sur 4 heures)}

\begin{aretenir}[Règle d'or : 5 minutes de relecture]
Sur une épreuve de 4 heures (Probatoire / Baccalauréat technique), réservez \textbf{5 minutes minimum} à la relecture. C'est le meilleur investissement temps/points : corriger une maladresse syntaxique, ajouter un connecteur manquant, vérifier l'accord d'un participe passé peut faire basculer une note de 11 à 13.
\end{aretenir}

La relecture ne se fait pas « en diagonale ». Elle suit une \textbf{check-list précise} :

\begin{popfigure}
\begin{tikzpicture}[scale=0.95, every node/.style={font=\small}]
  % Check-list visuelle : 5 cases à cocher
  \def\xstart{-7}
  \def\ystart{3}
  \def\dy{1.3}
  \def\boxsize{0.55}
  
  % Titre
  \node[font=\bfseries\large, text=popdark] at (0, 4.5) {CHECK-LIST DE RELECTURE (5 points de contrôle)};
  
  % 1. Cohérence Introduction / Développement / Conclusion
  \draw[popblue, thick] (\xstart, \ystart) rectangle (\xstart+\boxsize, \ystart+\boxsize);
  \node[right, text width=12cm, anchor=west] at (\xstart+0.8, \ystart+\boxsize/2) 
    {\textbf{1. Cohérence globale} : La conclusion répond-elle \emph{exactement} à la problématique de l'introduction ? Le développement a-t-il tenu les promesses du plan annoncé ?};
  
  % 2. Problématique présente et traitée
  \draw[popgreen, thick] (\xstart, \ystart-\dy) rectangle (\xstart+\boxsize, \ystart-\dy+\boxsize);
  \node[right, text width=12cm, anchor=west] at (\xstart+0.8, \ystart-\dy+\boxsize/2) 
    {\textbf{2. Problématique} : Est-elle clairement formulée en introduction ? Est-elle reprise mot à mot (ou synonymes précis) dans la conclusion ?};
  
  % 3. Plan suivi et connecteurs
  \draw[poporange, thick] (\xstart, \ystart-2*\dy) rectangle (\xstart+\boxsize, \ystart-2*\dy+\boxsize);
  \node[right, text width=12cm, anchor=west] at (\xstart+0.8, \ystart-2*\dy+\boxsize/2) 
    {\textbf{3. Plan \& connecteurs} : Les grandes parties sont-elles signalées (Premièrement, Ensuite, Enfin) ? Les transitions logiques (donc, cependant, par conséquent) sont-elles présentes ?};
  
  % 4. Exemples et références
  \draw[poppurple, thick] (\xstart, \ystart-3*\dy) rectangle (\xstart+\boxsize, \ystart-3*\dy+\boxsize);
  \node[right, text width=12cm, anchor=west] at (\xstart+0.8, \ystart-3*\dy+\boxsize/2) 
    {\textbf{4. Exemples \& références} : Chaque thèse discutée est-elle illustrée (exemple concret, contre-exemple, référence d'auteur) ? Les citations du texte sont-elles exactes et entre guillemets ?};
  
  % 5. Langue : orthographe, syntaxe, ponctuation
  \draw[poppink, thick] (\xstart, \ystart-4*\dy) rectangle (\xstart+\boxsize, \ystart-4*\dy+\boxsize);
  \node[right, text width=12cm, anchor=west] at (\xstart+0.8, \ystart-4*\dy+\boxsize/2) 
    {\textbf{5. Langue} : Accords (sujet/verbe, participe passé), ponctuation (virgules avant \og mais\fg{}, \og car\fg{}), pas de phrases trop longues (> 2 lignes), vocabulaire précis (pas de \og chose\fg{}, \og truc\fg{}).};
  
  % Cases à cocher symboliques (dessinées par le nœud TikZ lui-même)
  \foreach \i in {0,...,4} {
    \node[draw=popgray, thick, minimum size=0.4cm] at (\xstart+0.275, \ystart-\i*\dy+0.275) {};
  }
\end{tikzpicture}
\legende{Check-list visuelle de relecture : 5 cases à cocher systématiquement avant de rendre la copie.}
\end{popfigure}

\courssub{Grille d'auto-évaluation : suis-je prêt pour le Bac ?}

Avant de rendre votre copie, posez-vous ces 7 questions. Si vous répondez « non » à l'une d'elles, corrigez immédiatement.

\begin{center}
\begin{tabularx}{\linewidth}{|c|X|X|}\hline
\textbf{N°} & \textbf{Question} & \textbf{Action si NON} \\ \hline
1 & Ai-je une \textbf{problématique} explicite en introduction (question philosophique, pas un sujet) ? & Reformulez l'introduction. \\ \hline
2 & Mon \textbf{plan} (2 ou 3 parties) est-il annoncé et \textbf{suivi} rigoureusement ? & Realignez le développement sur le plan. \\ \hline
3 & Chaque partie \textbf{explique le texte} (pas de hors-texte) et \textbf{discute} (objection/nuance) ? & Ajoutez l'explication manquante ou la discussion. \\ \hline
4 & Ai-je des \textbf{exemples concrets} (camerounais si possible) et des \textbf{références d'auteurs} ? & Insérez un exemple ou une référence par partie. \\ \hline
5 & La \textbf{conclusion} répond-elle à la problématique en une phrase claire ? & Réécrivez la phrase de réponse. \\ \hline
6 & Y a-t-il une \textbf{ouverture} philosophique (question nouvelle, pas avis) ? & Remplacez par une vraie question philosophique. \\ \hline
7 & Ai-je \textbf{relue} : cohérence, connecteurs, orthographe, citations exactes ? & Passez 5 min sur la check-list. \\ \hline
\end{tabularx}
\end{center}

\courssub{TP guidé : correction collective d'une conclusion défectueuse}

\begin{exoresolu}[Correction d'une conclusion : « Le travail libère-t-il l'homme ? » (texte de Marx)]
\textbf{Problématique :} « Le travail, chez Marx, est-il aliénation ou réalisation de l'homme ? »

\textbf{Conclusion de l'élève (défectueuse) :}
\begin{quote}
\emph{« Pour conclure, on a vu que le travail c'est dur. Mais aussi ça peut être bien. Le travail libère l'homme si on aime son travail. Moi je pense que le travail c'est la vie. \textbf{Ouverture :} Est-ce que les robots vont remplacer les hommes au travail ? »}
\end{quote}
\tcblower
\textbf{Analyse des défauts (grille) :}
\begin{enumerate}
  \item \textbf{Pas de bilan structuré} : « on a vu que... mais aussi... » — pas de synthèse des étapes du développement.
  \item \textbf{Réponse floue} : « Le travail libère l'homme si on aime son travail » — condition subjective, pas de reprise de la problématique (aliénation/réalisation).
  \item \textbf{Avis personnel} : « Moi je pense que... » — interdit en commentaire.
  \item \textbf{Ouverture factuelle} : question sur les robots, pas philosophique.
  \item \textbf{Langue familière} : « c'est dur », « c'est la vie ».
\end{enumerate}

\textbf{Correction collective (réécriture guidée) :}

\begin{quote}
\emph{« Au terme de cette analyse, le texte de Marx montre que le travail est \textbf{aliénation} dans le capitalisme (produit confisqué, activité contrainte), mais qu'il porte en lui la \textbf{possibilité d'une réalisation humaine} dans une société sans classes où l'activité devient libre et consciente. \textbf{Aussi, le travail n'est pas en soi libération ni aliénation : il devient l'un ou l'autre selon les rapports sociaux de production qui l'encadrent.} \textit{Ouverture :} Cette analyse ne risque-t-elle pas de sous-estimer la dimension anthropologique du travail, c'est-à-dire le besoin naturel de l'homme de transformer la nature, indépendamment de tout mode de production ? »}
\end{quote}

\textbf{Pourquoi ça marche :}
\begin{itemize}
  \item Bilan : deux phrases qui résument la dialectique aliénation/réalisation.
  \item Réponse : phrase unique, termes de la problématique repris (\og aliénation\fg{}, \og réalisation\fg{}, \og rapports sociaux\fg{}).
  \item Ouverture : question philosophique sur l'anthropologique vs l'historique, née de la discussion.
  \item Langue : soutenue, connecteurs (\og aussi\fg{}, \og mais\fg{}).
\end{itemize}
\end{exoresolu}

\courssub{Exercice d'entraînement : à vous de jouer}

\begin{exercice}[Réécrire une conclusion]
\textbf{Texte :} Extrait du \emph{Contrat social} (Rousseau) : « L'obéissance à la loi qu'on s'est prescrite, c'est la liberté. »
\textbf{Problématique :} « La loi est-elle contrainte ou condition de la liberté ? »

\textbf{Conclusion proposée (à corriger) :}
\begin{quote}
\emph{« En conclusion, Rousseau dit que la loi c'est la liberté. C'est vrai parce qu'on vote les lois. Mais parfois les lois sont injustes. \textbf{Ouverture :} Est-ce qu'il faut désobéir aux lois injustes ? »}
\end{quote}

\textbf{Consignes :}
\begin{enumerate}
  \item Identifiez les 4 défauts majeurs (utilisez la grille d'analyse ci-dessus).
  \item Réécrivez la conclusion en respectant la structure en 3 temps (bilan, réponse, ouverture).
  \item Vérifiez votre production avec la check-list de relecture (5 points).
\end{enumerate}
\end{exercice}

\begin{corrige}[Exercice 1]
\textbf{Défauts identifiés :}
\begin{enumerate}
  \item Pas de bilan (pas de rappel de la discussion : loi générale vs loi particulière, volonté générale vs volonté de tous).
  \item Réponse approximative : « Rousseau dit que... » — pas de réponse tranchée à la problématique (contrainte OU condition).
  \item Argument d'autorité (« C'est vrai parce que... ») sans nuance.
  \item Ouverture : question de droit positif (désobéissance civile) pas directement philosophique sur le concept de liberté/loi.
\end{enumerate}

\textbf{Proposition de réécriture :}
\begin{quote}
\emph{« Notre analyse a montré que pour Rousseau, la loi n'est contrainte que lorsqu'elle émane d'une volonté particulière (volonté de tous), mais devient \textbf{condition de la liberté} lorsqu'elle exprime la \textbf{volonté générale} — loi que le citoyen se donne à lui-même en tant que membre du souverain. \textbf{Ainsi, la loi n'est pas l'opposée de la liberté : elle en est le fondement, pourvu qu'elle soit l'expression de la volonté générale et non la somme des intérêts privés.} \textit{Ouverture :} Cette identification de la liberté à l'obéissance à la loi qu'on s'est prescrite ne suppose-t-elle pas une transparence de la volonté générale que l'histoire politique a souvent démentie ? »}
\end{quote}

\textbf{Vérification check-list :}
\begin{itemize}
  \item[$\square$] Cohérence intro/développement/conclusion : OUI
  \item[$\square$] Problématique reprise : OUI (\og contrainte/condition\fg{}, \og volonté générale\fg{})
  \item[$\square$] Plan suivi + connecteurs : OUI (\og ainsi\fg{}, \og pourvu que\fg{})
  \item[$\square$] Exemples/références : OUI (volonté générale vs volonté de tous)
  \item[$\square$] Langue : OUI (soutenue, précise)
\end{itemize}
\end{corrige}

\courssub{Synthèse : les 3 commandements de la conclusion et de la relecture}

\begin{aretenir}[Les 3 commandements pour le Bac]
\begin{enumerate}
  \item \textbf{Commandement de la réponse} : « Tu répondras à la problématique en une phrase nette, reprenant ses termes, avant d'ouvrir. »
  \item \textbf{Commandement de la filiation} : « Ta conclusion ne contredira pas ton développement ; elle en sera le fruit visible. »
  \item \textbf{Commandement des 5 minutes} : « Tu reliras ta copie avec la check-list : cohérence, problématique, plan/connecteurs, exemples, langue. »
\end{enumerate}
\end{aretenir}

\begin{savaistu}[Le commentaire au Probatoire technique : chiffres clés]
\begin{itemize}
  \item \textbf{Durée :} 4 heures (240 minutes).
  \item \textbf{Temps idéal :} 30 min lecture/analyse, 30 min brouillon/plan, 150 min rédaction, \textbf{5 min conclusion}, \textbf{5 min relecture}, 20 min marge.
  \item \textbf{Longueur conclusion :} 5 à 8 lignes manuscrites (≈ 60--100 mots).
  \item \textbf{Notation :} Sur 20 pts, la conclusion + relecture valent \textbf{3 à 4 pts} (cohérence, réponse, ouverture, langue). Une conclusion manquante ou hors-sujet coûte cher.
  \item \textbf{Spécificité technique :} Les exemples concrets tirés du milieu professionnel (atelier, chantier, bureau d'études, exploitation agricole) sont \textbf{très valorisés} par les correcteurs MINESEC.
\end{itemize}
\end{savaistu}

\courssub{Conclusion de la leçon (méta-pédagogique)}

Vous disposez maintenant de la \textbf{boîte à outils complète} pour réussir un commentaire philosophique en Première technique : de la lecture analytique à la relecture finale, en passant par la problématisation, le développement et la conclusion. La méthode n'est pas une cage : c'est un \textbf{filet de sécurité} qui libère votre pensée en lui donnant une forme rigoureuse. Entraînez-vous sur des sujets variés (technique, travail, liberté, savoir, politique, morale), chronométrez-vous, corrigez-vous collectivement. Le jour de l'épreuve, vous ne découvrirez pas la méthode — vous la \textbf{maîtrisez}.

\begin{methode}[Rituel de révision hebdomadaire (15 min)]
\begin{enumerate}
  \item Choisissez un extrait court (10--15 lignes) au programme.
  \item Faites la lecture analytique (5 min) : thèse, concepts, structure, présupposés.
  \item Formulez une problématique (2 min).
  \item Écrivez \textbf{uniquement l'introduction et la conclusion} (8 min) — sans développement.
  \item Vérifiez avec la check-list : la conclusion répond-elle à la problématique ? Y a-t-il une ouverture ?
\end{enumerate}
\end{methode}

\emph{Prochaine étape : mise en situation complète (Épreuve blanche n°1) — sujet type Probatoire, 4 heures, correction collective.}

\coursec{Entraînement intégré : commentaire complet guidé (complément Bac)}

Après avoir maîtrisé l'art de conclure et de relire — cette ultime vigilance qui transforme un bon devoir en une copie excellente —, nous passons maintenant à l'exercice synthétique : le commentaire complet dans les conditions réelles de l'épreuve du Probatoire. C'est ici que toutes les compétences acquises (lecture, problématisation, planification, rédaction, relecture) s'articulent en une démarche unique, chronométrée et évaluée. Cette section constitue votre répétition générale avant l'examen.

\courssub{Présentation du texte support}

Pour cet entraînement, nous travaillons sur un extrait court, dense et typique des sujets du Bac technique, croisant les notions de \cle{technique}, de \cle{liberté} et de \cle{responsabilité} — au cœur de l'unité 1.

\begin{exemplebox}[Texte à commenter — Extrait inspiré de Hans Jonas, \emph{Le Principe responsabilité} (1979)]
« La technique moderne a modifié la nature de l'action humaine à un point tel que l'éthique traditionnelle, fondée sur la proximité spatiale et temporelle des effets, ne suffit plus à la guider. Jadis, l'agir humain était borné par sa propre puissance limitée ; aujourd'hui, la puissance technique démultipliée projette nos actes dans un avenir lointain et sur une échelle planétaire. La responsabilité doit donc s'étendre à la mesure de cette puissance : elle ne saurait se borner au cercle étroit des contemporains, mais doit embrasser les générations futures et la biosphère tout entière. Ainsi, le commandement éthique nouveau s'énonce : \emph{Agis de sorte que les effets de ton action soient compatibles avec la permanence d'une vie authentiquement humaine sur terre.} »
\end{exemplebox}

\courssub{Étape 1 : Lecture analytique en binômes avec grille}

\emph{Durée : 30 minutes (inclus dans la phase d'analyse globale).} En binôme, vous remplissez la grille ci-dessous. Chaque case doit être justifiée par une citation précise (numéro de ligne ou reprise de mots).

\begin{experience}[Grille de lecture analytique — Travail en binômes]
\begin{center}
\begin{tabularx}{\linewidth}{|l|X|}\hline
\textbf{Élément} & \textbf{Repérage et citations} \\ \hline
\textbf{Thème général} &  \\ \hline
\textbf{Thèse défendue} &  \\ \hline
\textbf{Arguments principaux} (3 minimum) & 1. \\ & 2. \\ & 3. \\ \hline
\textbf{Moments du texte} (découpage en 3--4 séquences logiques) & §1 (l. 1--...) : \\ & §2 (l. ...--...) : \\ & §3 (l. ...--...) : \\ \hline
\textbf{Concepts-clés à définir} & 1. \emph{Technique moderne} -- 2. \emph{Éthique traditionnelle} -- 3. \emph{Responsabilité} -- 4. \emph{Puissance technique} -- 5. \emph{Vie authentiquement humaine} \\ \hline
\textbf{Type de raisonnement} & (déductif, analogique, \emph{a fortiori}, etc.) \\ \hline
\textbf{Portée / Enjeu} & Pourquoi ce texte nous concerne-t-il, nous, futurs techniciens ? \\ \hline
\end{tabularx}
\end{center}
\par\textit{Légende : Grille à photocopier ou reproduire au brouillon. Chaque case exige une preuve textuelle.}
\end{experience}

\begin{aretenir}[Résultats attendus de l'analyse (corrigé indicatif)]
\begin{itemize}
\item \textbf{Thème} : La mutation de la responsabilité éthique à l'ère de la technique moderne.
\item \textbf{Thèse} : La responsabilité doit s'étendre spatialement (planétaire) et temporellement (générations futures) à la mesure de la puissance technique.
\item \textbf{Arguments} : (1) Rupture anthropologique : la technique change la nature de l'action. (2) Inadéquation de l'éthique traditionnelle (proximité, effets proches). (3) Nécessité d'un nouvel impératif éthique universel et prospectif.
\item \textbf{Moments} : §1 = Constat de rupture (technique moderne vs agir d'autrefois). §2 = Conséquence : extension de la responsabilité. §3 = Formulation du nouvel impératif (principe responsabilité).
\item \textbf{Concepts} : \cle{Technique moderne} = ensemble des moyens artificiels démultipliant la puissance humaine ; \cle{Éthique traditionnelle} = morale de la proximité (voisinage, temps présent) ; \cle{Responsabilité} = obligation de répondre des conséquences de ses actes ; \cle{Puissance technique} = capacité d'agir à grande échelle et long terme ; \cle{Vie authentiquement humaine} = vie digne, durable, non aliénée.
\item \textbf{Raisonnement} : Argument \emph{a fortiori} (si l'éthique traditionnelle valait pour peu de puissance, \emph{a fortiori} faut-il une éthique renforcée pour une puissance immense) + analogie (nouvel impératif kantien).
\item \textbf{Enjeu technique} : En tant que futurs ingénieurs, techniciens, gestionnaires, nos décisions techniques (choix matériaux, procédés, implantation) engagent la responsabilité au sens de Jonas.
\end{itemize}
\end{aretenir}

\begin{attention}[Piège fréquent en binôme]
Ne pas confondre \emph{résumé} et \emph{analyse}. La grille n'attend pas « l'auteur dit que... » mais « l'argument 1 est : ... ; il repose sur le passage "..." ». Évitez les cases vides : si un argument vous semble implicite, écrivez « Implicite : l'auteur suppose que... » et citez l'indice textuel.
\end{attention}

\courssub{Étape 2 : Formulation de la problématique}

\emph{Durée : 15 minutes (individuel puis mise en commun).} Chaque élève rédige \textbf{une seule phrase interrogative} qui : (1) naît de la tension interne au texte, (2) appelle une réponse argumentée, (3) dépasse le texte sans le trahir.

\begin{methode}[Formuler une problématique en 3 temps]
\begin{enumerate}
\item \textbf{Identifiez la tension} : ici, tension entre \emph{puissance technique illimitée} et \emph{éthique traditionnelle bornée}.
\item \textbf{Transformez en question philosophique} : « Comment l'éthique peut-elle répondre au décalage entre la puissance technique et sa responsabilité ? » ou « La responsabilité technique exige-t-elle une refonte de l'éthique ? »
\item \textbf{Testez la pertinence} : La question permet-elle d'organiser un plan en 3 parties ? Ouvre-t-elle à la discussion (thèses contraires possibles) ? Reste-t-elle ancrée dans le texte ?
\end{enumerate}
\end{methode}

\begin{exemplebox}[Problématiques proposées par la classe — Choix collectif]
\begin{enumerate}
\item \emph{« Jusqu'où la responsabilité du technicien doit-elle s'étendre face à la puissance de la technique moderne ? »} (Centrée sur l'agent technique — très pertinente pour le Bac Tech)
\item \emph{« L'éthique traditionnelle est-elle devenue obsolète face aux défis de la technique contemporaine ? »} (Centrée sur la théorie éthique)
\item \emph{« Comment fonder une responsabilité à la hauteur de la puissance technique ? »} (Centrée sur la fondation — plus abstraite)
\end{enumerate}
\textbf{Choix retenu pour la suite} : Problématique n°1 (ancrage concret, lisibilité, faisabilité en 3 parties).
\end{exemplebox}

\courssub{Étape 3 : Construction du plan détaillé}

\emph{Durée : 30 minutes (au brouillon, collectivement guidé par le professeur).} Le plan doit répondre \textbf{exactement} à la problématique, suivre la logique du texte (explicatif) puis l'ouvrir (discussion), et présenter des titres de parties et sous-parties \emph{argumentatifs} (phrases nominales complètes, pas de mots seuls).

\begin{propriete}[Structure-type d'un plan de commentaire au Bac (3 parties, 3 sous-parties chacune)]
\begin{itemize}
\item \textbf{Partie I — Explication :} Suivre le fil du texte, montrer \emph{comment} l'auteur établit sa thèse.
\item \textbf{Partie II — Discussion / Approfondissement :} Confronter la thèse à d'autres points de vue, en préciser la portée, en tester les limites.
\item \textbf{Partie III — Dépassement / Perspectives :} Apporter une réponse personnelle argumentée à la problématique, en s'appuyant sur la discussion.
\end{itemize}
\end{propriete}

\begin{exoresolu}[Plan détaillé pour la problématique n°1]
\textbf{Introduction} (rédigée in extenso à l'étape 4)

\textbf{I. La technique moderne bouleverse les conditions de l'action et de la responsabilité}
\begin{enumerate}
\item A. La démultiplication technique brise les bornes de l'agir traditionnel (proximité, réversibilité, prévisibilité).
\item B. L'éthique de la proximité devient inadéquate : elle ne prévoit ni l'échelle planétaire, ni le temps long.
\item C. La responsabilité doit donc s'étendre à la mesure de la nouvelle puissance : spatialement (biosphère) et temporellement (générations futures).
\end{enumerate}

\textbf{II. Les défis et les limites de cette extension de la responsabilité}
\begin{enumerate}
\item A. Difficulté épistémique : comment prévoir les effets à très long terme d'innovations techniques complexes ? (Ex. : CFC, nucléaire, IA).
\item B. Risque d'une responsabilité diluée : si tout le monde est responsable de tout, personne ne l'est vraiment (problème de l'agent collectif).
\item C. Tension avec la liberté technique et l'innovation : une responsabilité trop lourde ne paralyse-t-elle pas le progrès nécessaire ?
\end{enumerate}

\textbf{III. Vers une éthique de la responsabilité technique située et partagée}
\begin{enumerate}
\item A. Le technicien comme premier responsable : intégrer l'évaluation d'impact \emph{amont} (conception responsable, \emph{privacy by design}, éco-conception).
\item B. Responsabilité partagée : cadres réglementaires (lois, normes ISO), comités d'éthique, vigilance citoyenne — la responsabilité ne repose pas sur l'individu seul.
\item C. Le principe responsabilité comme boussole, non comme paralysie : agir \emph{sous le régime de l'hypothèse} (Jonas), réviser ses choix, privilégier la réversibilité technique.
\end{enumerate}

\textbf{Conclusion} (synthèse + ouverture)
\end{exoresolu}

\begin{attention}[Erreur fatale : recopier le plan au propre sans rédiger]
Au brouillon, votre plan est \emph{détaillé} (arguments, exemples, références, transitions). Au propre, vous \textbf{rédigez des paragraphes complets}. Ne copiez pas les titres de parties comme des étiquettes vides : chaque sous-partie = un paragraphe (ou deux) avec phrase d'amorce, argument, exemple/référence, phrase de transition. Le correcteur note la \emph{rédaction}, pas le squelette.
\end{attention}

\courssub{Étape 4 : Rédaction de l'introduction et d'un paragraphe de développement}

\emph{Durée : 45 minutes (introduction 15 min, paragraphe 30 min).} Vous rédigez \textbf{au propre} (sur la copie d'examen) l'introduction complète et le premier paragraphe de la partie I (I.A). Le reste sera fait en devoir maison ou en classe complète.

\begin{methode}[Rédaction de l'introduction — 4 mouvements obligatoires]
\begin{enumerate}
\item \textbf{Accroche / Contextualisation} : Situer le texte (auteur, ouvrage, date, courant) et le thème général.
\item \textbf{Présentation du texte} : Résumer en 2--3 lignes le cheminement de la pensée (thèse + arguments principaux).
\item \textbf{Problématique} : La question unique, issue de l'analyse.
\item \textbf{Annonce du plan} : « Nous verrons d'abord que... (I), puis nous examinerons... (II), pour enfin montrer... (III). »
\end{enumerate}
\end{methode}

\begin{exoresolu}[Introduction rédigée — Modèle]
\textbf{Hans Jonas}, philosophe allemand naturalisé américain (1903--1993), publie en 1979 \emph{Le Principe responsabilité}, ouvrage fondateur de l'éthique de l'environnement et de la technique. Dans l'extrait proposé, il dresse le constat d'une rupture anthropologique majeure : la technique moderne a démultiplié la puissance d'agir de l'humanité au point de rendre caduque l'éthique traditionnelle, cantonnée à la proximité des effets. Il en déduit que la responsabilité doit s'étendre à l'échelle planétaire et temporelle, et formule un nouvel impératif catégorique : agir pour que les effets de nos actions soient compatibles avec la permanence d'une vie authentiquement humaine sur terre.

Ce texte soulève une question cruciale pour nous, futurs techniciens : \textbf{jusqu'où la responsabilité du technicien doit-elle s'étendre face à la puissance de la technique moderne ?}

Pour y répondre, nous montrerons d'abord que la technique moderne bouleverse les conditions de l'action et de la responsabilité (I). Nous analyserons ensuite les défis et les limites de cette extension (II). Enfin, nous esquisserons les contours d'une éthique de la responsabilité technique située et partagée (III).
\end{exoresolu}

\begin{exoresolu}[Paragraphe I.A rédigé — Modèle]
\textbf{I. A. La démultiplication technique brise les bornes de l'agir traditionnel (proximité, réversibilité, prévisibilité).}

Jonas ouvre son texte par un constat radical : « La technique moderne a modifié la nature de l'action humaine à un point tel que l'éthique traditionnelle [...] ne suffit plus à la guider. » Cette phrase initiale pose la thèse directrice : il y a \emph{mutation de l'agir}, non simple amplification. Jadis, l'action humaine était « bornée par sa propre puissance limitée » : un artisan, un paysan, un bâtisseur agissait ici et maintenant, avec des outils dont les effets étaient immédiats, localisés, prévisibles et souvent réversibles. L'éthique traditionnelle (aristotélicienne, kantienne, utilitariste) s'est construite sur cette \emph{proximité} : je suis responsable de ce que je vois, touche, prévois à courte échéance.

Or, la technique moderne — pensons à l'industrie chimique, au nucléaire, aux algorithmes d'IA, aux biotechnologies — projette nos actes « dans un avenir lointain et sur une échelle planétaire ». L'exemple des CFC (chlorofluorocarbures), conçus comme fluides frigorigènes inertes et sûrs, illustre ce décalage : leur effet destructeur sur la couche d'ozone n'a été découvert que des décennies plus tard, à l'échelle globale. La \cle{proximité} (spatiale et temporelle), condition de l'imputation morale classique, est brisée. Dès lors, l'éthique de la responsabilité (au sens de Weber : répondre des conséquences prévisibles) doit elle-même muter : elle ne peut plus se borner au « cercle étroit des contemporains ».
\end{exoresolu}

\begin{aretenir}[Critères de réussite du paragraphe]
\begin{itemize}
\item \textbf{Ancrage textuel} : citations intégrées, mots-clés repris (« nature de l'action », « bornée », « proximité »).
\item \textbf{Explicitation} : le sens des concepts est déployé (proximité, réversibilité, prévisibilité).
\item \textbf{Illustration concrète} : exemple des CFC (technique, daté, pertinent).
\item \textbf{Articulation} : le paragraphe montre \emph{comment} l'argument fonctionne (rupture -> inadéquation -> nécessité d'extension).
\item \textbf{Transition implicite} : la dernière phrase ouvre sur la sous-partie B (l'inadéquation de l'éthique traditionnelle).
\end{itemize}
\end{aretenir}

\courssub{Étape 5 : Évaluation croisée avec la grille critériée du Bac}

\emph{Durée : 30 minutes.} Chaque binôme échange sa copie avec un autre binôme. À l'aide de la grille officielle (adaptée ci-dessous), vous attribuez une note sur 20 et justifiez chaque item par un commentaire précis. L'enseignant collecte ensuite les grilles pour une analyse collective.

\begin{popfigure}
\begin{tikzpicture}[scale=0.85, every node/.style={font=\small}]
% Couleurs
\definecolor{popblueL}{RGB}{200,220,255}
\definecolor{popgreenL}{RGB}{200,255,200}
\definecolor{poporangeL}{RGB}{255,230,200}
\definecolor{poppinkL}{RGB}{255,200,220}
\definecolor{popgoldL}{RGB}{255,245,200}

% En-têtes
\fill[popblueL] (0,0) rectangle (14,0.8);
\node[anchor=west] at (0.3,0.4) {\textbf{CRITÈRE}};
\node[anchor=center] at (4,0.4) {\textbf{INDICATEURS OBSERVABLES}};
\node[anchor=center] at (10,0.4) {\textbf{BAREME}};
\node[anchor=center] at (12.5,0.4) {\textbf{NOTE / OBS.}};

% Lignes
% Compréhension
\fill[popblueL!30] (0,0.8) rectangle (14,2.2);
\node[anchor=west] at (0.3,1.5) {\textbf{1. Compréhension du texte}};
\node[anchor=west, align=left] at (4.2,1.5) {• Thème, thèse, arguments identifiés avec citations\\• Concepts-clés définis et articulés\\• Enjeu philosophique perçu};
\node[anchor=center] at (10,1.5) {\textbf{/4}};
\node[anchor=center] at (12.5,1.5) {.......};

% Plan / Structure
\fill[popgreenL!30] (0,2.2) rectangle (14,3.6);
\node[anchor=west] at (0.3,2.9) {\textbf{2. Plan / Structure}};
\node[anchor=west, align=left] at (4.2,2.9) {• Plan en 3 parties répondant à la problématique\\• Titres argumentatifs, équilibrés, progressifs\\• Transitions logiques annoncées/réalisées};
\node[anchor=center] at (10,2.9) {\textbf{/4}};
\node[anchor=center] at (12.5,2.9) {.......};

% Argumentation / Contenu
\fill[poporangeL!30] (0,3.6) rectangle (14,5.0);
\node[anchor=west] at (0.3,4.3) {\textbf{3. Argumentation / Contenu}};
\node[anchor=west, align=left] at (4.2,4.3) {• Explication fidèle et approfondie du texte (I)\\• Discussion critique pertinente, exemples variés (II)\\• Dépassement personnel argumenté, cohérent (III)\\• Références philosophiques exactes (auteurs, notions)};
\node[anchor=center] at (10,4.3) {\textbf{/6}};
\node[anchor=center] at (12.5,4.3) {.......};

% Expression / Langue
\fill[poppinkL!30] (0,5.0) rectangle (14,6.4);
\node[anchor=west] at (0.3,5.7) {\textbf{4. Expression / Langue}};
\node[anchor=west, align=left] at (4.2,5.7) {• Français correct (syntaxe, orthographe, vocabulaire)\\• Style philosophique (précision, connecteurs logiques)\\• Ponctuation, paragraphage, lisibilité\\• Respect du temps (copie complète)};
\node[anchor=center] at (10,5.7) {\textbf{/6}};
\node[anchor=center] at (12.5,5.7) {.......};

% Total
\fill[popgoldL] (0,6.4) rectangle (14,7.2);
\node[anchor=west] at (0.3,6.8) {\textbf{TOTAL}};
\node[anchor=center] at (10,6.8) {\textbf{/20}};
\node[anchor=center] at (12.5,6.8) {.......};

% Grille
\draw[thin, gray] (0,0) grid[step=0.8] (14,7.2);
\draw[thick] (0,0) rectangle (14,7.2);
\draw[thick] (3.5,0) -- (3.5,7.2);
\draw[thick] (9.5,0) -- (9.5,7.2);
\draw[thick] (12,0) -- (12,7.2);
\draw[thick] (0,0.8) -- (14,0.8);
\draw[thick] (0,2.2) -- (14,2.2);
\draw[thick] (0,3.6) -- (14,3.6);
\draw[thick] (0,5.0) -- (14,5.0);
\draw[thick] (0,6.4) -- (14,6.4);
\end{tikzpicture}
\legende{Grille critériée d'évaluation du commentaire philosophique (adaptée du barème officiel Probatoire / Bac Tech). Chaque critère est noté sur la valeur indiquée ; le total fait 20.}
\end{popfigure}

\begin{methode}[Évaluation croisée efficace — 5 étapes]
\begin{enumerate}
\item \textbf{Lisez la copie en entier} sans annoter (2 min) pour saisir la cohérence globale.
\item \textbf{Appliquez la grille critère par critère} : soulignez dans la copie les éléments qui justifient la note (citations, connecteurs, exemples).
\item \textbf{Notez et commentez} : « Compréhension : 3/4 — thèse bien vue, mais argument 2 mal explicité (pas de citation). »
\item \textbf{Échangez les grilles} : l'auteur lit l'évaluation, pose des questions, note les points à améliorer.
\item \textbf{Synthèse collective} : l'enseignant relève les forces/faiblesses récurrentes (ex. : « Beaucoup de 2/4 en argumentation : exemples trop généraux, pas de références précises »).
\end{enumerate}
\end{methode}

\courssub{Application aux situations concrètes : du texte philosophique au réel camerounais}

La méthodologie du commentaire n'est pas un exercice scolaire enfermé : c'est une \cle{compétence transférable} pour tout technicien. Analyser un texte argumentatif, en extraire la thèse, les arguments, les présupposés, en évaluer la portée — cette rigueur s'applique directement à votre futur professionnel.

\begin{savaistu}[La lecture analytique dans la vie technique au Cameroun]
\begin{itemize}
\item \textbf{Cahier des charges / Appel d'offres} : Comme un texte de philosophie, un Cdc contient une thèse (le besoin), des arguments (contraintes techniques, normes, délais), des concepts-clés (spécifications), des implicites (budget réel, priorités politiques). Le lire « philosophiquement » = repérer les contradictions, les zones floues, les non-dits \emph{avant} de chiffrer.
\item \textbf{Contrat de travail / Convention collective} : Structure argumentative (droits $\leftrightarrow$ obligations), définitions juridiques (concepts-clés), portée (champ d'application). La même grille (thème, thèse, arguments, enjeux) évite les pièges.
\item \textbf{Règlement de sécurité / Norme ISO} : Texte normatif à portée universelle. L'analyse en « moments » (champ d'application, exigences, annexes) et la problématisation (« Cette norme est-elle adaptée à notre contexte local ? ») guident la mise en conformité.
\item \textbf{Débat radiophonique / Publication réseaux sociaux} : Discours souvent émotionnel, fallacieux, non structuré. Appliquer la grille : quel est le thème ? La thèse implicite ? Les arguments (autorité, analogie, peur, chiffre) ? Les concepts flous ? L'enjeu réel ? Cela forme l'esprit critique citoyen et technique.
\end{itemize}
\end{savaistu}

\begin{exemplebox}[Exercice d'application — Analyse d'une publication Facebook]
\textbf{Publication} : « Arrêtez de construire des barrages ! Le barrage de Lom Pangar a déplacé des villages, détruit la forêt, et l'électricité ne profite qu'aux usines. Le développement durable, c'est du mensonge. Partagez si vous êtes d'accord. »

\textbf{Travail (10 min, individuel) :} Remplissez la grille simplifiée :
\begin{enumerate}
\item Thème :
\item Thèse implicite :
\item Arguments (explicites / implicites) :
\item Concepts-clés mal définis ou instrumentalisés :
\item Type de raisonnement (fallacieux ?) :
\item Enjeu réel pour un technicien :
\item Réponse argumentée courte (3 lignes) :
\end{enumerate}
\end{exemplebox}

\begin{corrige}[Exercice — Publication Facebook]
\begin{enumerate}
\item \textbf{Thème} : Acceptabilité sociale des grands ouvrages hydrauliques au Cameroun.
\item \textbf{Thèse implicite} : Les barrages sont nuisibles globalement ; le « développement durable » est un leurre.
\item \textbf{Arguments} : Explicites : déplacements, déforestation, inégalité de répartition de l'électricité. Implicites : toute industrialisation = prédation ; pas d'alternative viable.
\item \textbf{Concepts} : « Développement durable » réduit à « mensonge » (définition caricaturale) ; « profit » limité aux usines (occulte l'électrification rurale, l'irrigation, la régulation de crue).
\item \textbf{Raisonnement} : Generalisation abusive (un cas $\rightarrow$ tous les barrages) ; appel à l'émotion (déforestation, déplacés) ; faux dilemme (barrage OU rien).
\item \textbf{Enjeu technicien} : Savoir répondre techniquement (bilan coûts-avantages, mesures compensatoires, mix énergétique) et éthiquement (justice distributive, consentement libre et éclairé des populations).
\item \textbf{Réponse} : « Lom Pangar a des impacts réels (30 000 ha inondés, 1 200 déplacés), mais alimente le réseau interconnecté Sud (stabilité), régule la Sanaga (crue/étiage), et finance l'électrification rurale. Le développement durable n'est pas un slogan : c'est une méthode d'arbitrage entre piliers économique, social, environnemental — perfectible, pas mensonger. »
\end{enumerate}
\end{corrige}

\courssub{Gestion du temps : la frise chronologique de l'épreuve (4 heures)}

Le jour de l'épreuve, le temps est votre ressource la plus rare. Une gestion rigoureuse, répétée à l'entraînement, devient un automatisme. Voici la répartition officielle recommandée, visualisée sur une frise.

\begin{popfigure}
\begin{tikzpicture}[scale=0.95, every node/.style={font=\small}]
% Couleurs
\definecolor{phase1}{RGB}{100,200,150}
\definecolor{phase2}{RGB}{100,150,250}
\definecolor{phase3}{RGB}{250,200,80}
\definecolor{phase4}{RGB}{250,120,120}
\definecolor{phase1L}{RGB}{200,245,230}
\definecolor{phase2L}{RGB}{200,220,255}
\definecolor{phase3L}{RGB}{255,245,200}
\definecolor{phase4L}{RGB}{255,220,220}

% Axes et graduations (4h = 240 min)
\draw[thick, ->] (0,0) -- (24.5,0) node[right] {\textbf{Temps (minutes)}};
\foreach \x/\label in {0/0, 3/30, 6/60, 9/90, 12/120, 15/150, 18/180, 21/210, 24/240} {
  \draw (\x, -0.15) -- (\x, 0.15) node[below=2pt] {\label};
}

% Barres de phases
% Phase 1 : Lecture + Analyse (30 min)
\fill[phase1L] (0, 0.5) rectangle (3, 1.5);
\draw[phase1, thick] (0, 0.5) rectangle (3, 1.5);
\node[anchor=center, text=phase1] at (1.5, 1.0) {\textbf{LECTURE \& ANALYSE (30')}};
\node[anchor=center, text=phase1, font=\tiny] at (1.5, 0.3) {3 lectures + Grille + Problématique};

% Phase 2 : Plan détaillé (30 min)
\fill[phase2L] (3, 0.5) rectangle (6, 1.5);
\draw[phase2, thick] (3, 0.5) rectangle (6, 1.5);
\node[anchor=center, text=phase2] at (4.5, 1.0) {\textbf{PLAN DÉTAILLÉ (30')}};
\node[anchor=center, text=phase2, font=\tiny] at (4.5, 0.3) {3 parties $\times$ 3 sous-parties + arguments/exemples};

% Phase 3 : Rédaction (150 min = 2h30)
\fill[phase3L] (6, 0.5) rectangle (21, 1.5);
\draw[phase3, thick] (6, 0.5) rectangle (21, 1.5);
\node[anchor=center, text=phase3] at (13.5, 1.0) {\textbf{RÉDACTION AU PROPRE (2h30 = 150')}};
\node[anchor=center, text=phase3, font=\tiny] at (13.5, 0.3) {Intro (15') -- Développement (2h) -- Conclusion (15')};

% Phase 4 : Relecture (30 min)
\fill[phase4L] (21, 0.5) rectangle (24, 1.5);
\draw[phase4, thick] (21, 0.5) rectangle (24, 1.5);
\node[anchor=center, text=phase4] at (22.5, 1.0) {\textbf{RELECTURE (30')}};
\node[anchor=center, text=phase4, font=\tiny] at (22.5, 0.3) {Cohérence + Langue + Complétude + Propreté};

% Repères visuels
\draw[dashed, gray] (3,0) -- (3,-0.5);
\draw[dashed, gray] (6,0) -- (6,-0.5);
\draw[dashed, gray] (21,0) -- (21,-0.5);

% Légende
\node[anchor=west, align=left, font=\footnotesize] at (0, -1.8) {
\textbf{Légende :} \textcolor{phase1}{■} Analyse (lecture, grille, problématique) \quad
\textcolor{phase2}{■} Planification (brouillon complet) \quad
\textcolor{phase3}{■} Rédaction (copie finale) \quad
\textcolor{phase4}{■} Relecture (correction, vérification)
};
\node[anchor=west, font=\footnotesize\itshape] at (0, -2.4) {Total : 240 minutes = 4 heures exactement. Aucune marge : chaque minute compte.};
\end{tikzpicture}
\legende{Frise chronologique de gestion du temps pour l'épreuve de commentaire philosophique (4 heures). Respectez les bornes : ne débordez jamais sur la phase suivante.}
\end{popfigure}

\begin{methode}[Gérer son temps au Bac — Règle d'or]
\begin{enumerate}
\item \textbf{Jamais de rédaction avant plan complet.} Si à 1h00 (fin phase 2) votre plan n'est pas fini (3 parties, 9 sous-parties, arguments, exemples, transitions notés au brouillon), \textbf{ne commencez pas à recopier}. Prenez 10 min de plus sur la rédaction : un plan solide gagne 30 min de rédaction fluide.
\item \textbf{Chronomètre visible.} Posez votre montre ou minuteur sur la table. Points de passage obligatoires : 1h00 = plan fini ; 3h30 = fin rédaction ; 4h00 = copie rendue.
\item \textbf{Rédaction par blocs.} N'écrivez pas l'introduction, puis I.A, puis I.B... Écrivez \textbf{bloc par bloc} : Intro (15 min) $\rightarrow$ Partie I complète (45 min) $\rightarrow$ Partie II (50 min) $\rightarrow$ Partie III (40 min) $\rightarrow$ Conclusion (15 min). Si un bloc dépasse, coupez net, passez au suivant, revenez en relecture.
\item \textbf{Relecture en 3 passes.} (1) Cohérence globale : thèse tenue ? plan respecté ? transitions ? (2) Langue : fautes, lourdeurs, vocabulaire précis ? (3) Forme : paragraphes, lisibilité, ratures propres, numérotation des pages.
\end{enumerate}
\end{methode}

\begin{exemplebox}[Exemple chiffré : répartition type sur 4 heures]
\begin{center}
\begin{tabularx}{\linewidth}{|l|c|X|}\hline
\textbf{Phase} & \textbf{Durée} & \textbf{Objectif concret} \\ \hline
Lecture 1 (découverte) & 5 min & Imprégnation globale, ton, structure apparente \\ \hline
Lecture 2 (crayon) & 10 min & Surlignage thèse, arguments, concepts, moments \\ \hline
Lecture 3 (grille) & 15 min & Remplissage grille analytique + problématique \\ \hline
\textbf{Total Analyse} & \textbf{30 min} & \textbf{Plan prêt à être construit} \\ \hline
Plan détaillé brouillon & 30 min & 9 titres argumentatifs + 2--3 mots-clés/sous-partie + exemples + refs \\ \hline
Introduction & 15 min & 4 mouvements rédigés au propre \\ \hline
Partie I (3 $\times$ 15 min) & 45 min & Explication textuelle rigoureuse \\ \hline
Partie II (3 $\times$ 17 min) & 50 min & Discussion critique + exemples variés \\ \hline
Partie III (3 $\times$ 13 min) & 40 min & Dépassement personnel argumenté \\ \hline
Conclusion & 15 min & Synthèse + ouverture (1 phrase) \\ \hline
\textbf{Total Rédaction} & \textbf{2h30} & \textbf{Copie complète, environ 4--5 pages} \\ \hline
Relecture cohérence & 10 min & Thèse / plan / transitions / exemples \\ \hline
Relecture langue & 15 min & Fautes, syntaxe, connecteurs, vocabulaire \\ \hline
Relecture forme & 5 min & Paragraphes, ratures, numérotation, signature \\ \hline
\textbf{Total Relecture} & \textbf{30 min} & \textbf{Copie propre, lisible, sans trous} \\ \hline
\textbf{TOTAL} & \textbf{4h00} &  \\ \hline
\end{tabularx}
\end{center}
\end{exemplebox}

\begin{attention}[Pièges temporels mortels]
\begin{itemize}
\item \textbf{« Je vais lire une 4e fois pour être sûr. »} $\rightarrow$ Perte de 10 min, plan bâclé, rédaction hachée.
\item \textbf{« J'écris l'intro au brouillon, je la recopierai. »} $\rightarrow$ Double travail, temps volé au développement. \textbf{Rédigez l'intro directement au propre.}
\item \textbf{« Je finis la partie III en relecture. »} $\rightarrow$ Partie III tronquée = note plafonnée à 10/20 max. Mieux vaut une partie III courte mais présente + conclusion.
\item \textbf{« Je n'ai pas le temps de relire. »} $\rightarrow$ 3--4 points perdus sur « Expression » (sur 6) + incohérences non corrigées. La relecture \emph{est} la différence entre 12 et 15.
\end{itemize}
\end{attention}

\begin{aretenir}[Check-list finale avant de rendre la copie]
\begin{itemize}
\item [ ] Introduction : 4 mouvements présents, problématique claire, plan annoncé.
\item [ ] Plan : 3 parties, 9 sous-parties, titres argumentatifs, respectés dans le développement.
\item [ ] Développement : Chaque sous-partie = paragraphe(s) avec amorce, argument, exemple/référence, transition.
\item [ ] Explication (I) : Fidèle au texte, citations intégrées, concepts définis.
\item [ ] Discussion (II) : Thèses contraires nommées (auteurs), exemples concrets, limites montrées.
\item [ ] Dépassement (III) : Réponse personnelle à la problématique, argumentée, cohérente.
\item [ ] Conclusion : Synthèse (réponse) + Ouverture (nouvelle question / perspective).
\item [ ] Langue : 0 faute évitable, connecteurs logiques (\emph{donc, toutefois, en effet, par conséquent}), vocabulaire précis.
\item [ ] Forme : Paragraphes visibles (alinéa ou saut de ligne), ratures nettes (trait unique), pages numérotées.
\end{itemize}
\end{aretenir}

\begin{savaistu}[L'éthique de la responsabilité technique : un atout professionnel au Cameroun]
Au-delà de l'examen, la démarche du commentaire — \emph{analyser, problématiser, argumenter, conclure} — est la signature du \textbf{technicien supérieur responsable}. Au Cameroun, les grands chantiers (barrage de Nachtigal, autoroute Yaoundé-Douala, fibre optique, plateformes numériques) exigent des ingénieurs et techniciens capables de :
\begin{itemize}
\item Lire un cahier des charges comme un texte philosophique : en extraire les présupposés, les zones d'ombre, les enjeux éthiques.
\item Argumenter un choix technique (matériau, procédé, implantation) face à une hiérarchie, un client, une population riveraine.
\item Anticiper les effets à long terme (maintenance, environnement, social) — l'esprit du \emph{Principe responsabilité} de Jonas.
\item Rédiger des rapports clairs, structurés, sans ambiguïté — la même rigueur que la dissertation ou le commentaire.
\end{itemize}
Maîtriser le commentaire philosophique, c'est déjà exercer votre métier de technicien citoyen.
\end{savaistu}

\begin{exercice}[Entraînement complet — Sujet type Bac blanc (à faire en 4h chrono)]
\textbf{Texte} : Extrait de \emph{La Condition humaine} d'Hannah Arendt (1958) sur le \emph{homo faber} et la fabrication du monde artificiel.
\textbf{Consigne} : Commentez ce texte.
\textbf{Modalités} : Appliquez intégralement la frise chronologique. Remettez : grille d'analyse, plan détaillé, copie complète, grille d'auto-évaluation remplie.
\textbf{Évaluation} : Votre professeur utilisera la grille critériée (Figure 2) + grille d'auto-évaluation pour une note sur 20 + commentaires personnalisés.
\end{exercice}

\begin{corrige}[Exercice — Éléments attendus pour l'auto-évaluation]
\begin{itemize}
\item \textbf{Grille d'analyse} : Thème (fabrication / monde artificiel), Thèse (l'homo faber produit la permanence mais instrumentalise le monde), Arguments (utilité, durabilité, réification), Moments (définition $\rightarrow$ caractère violent de la fabrication $\rightarrow$ monde humain), Concepts (\emph{homo faber}, \emph{instrumentalité}, \emph{monde}, \emph{réification}).
\item \textbf{Problématique type} : « La fabrication technique humanise-t-elle le monde ou le réduit-elle à n'être qu'un stock de moyens ? »
\item \textbf{Plan} : I. Le \emph{homo faber} comme producteur de la permanence humaine. II. L'instrumentalisation du monde : la nature réduite à ressource. III. Vers une technique poétique ? Réhabiliter le \emph{faire} comme soin du monde.
\item \textbf{Auto-évaluation} : Soyez honnête : « J'ai passé 45 min sur l'analyse $\rightarrow$ plan incomplet $\rightarrow$ partie III expédiée. Note estimée : 11/20. Prochaine fois : respecter les 30 min analyse. »
\end{itemize}
\end{corrige}

\newpage

\coursec{Activité d'intégration}

\coursec{Leçon 3 : Consolidation des acquis de la méthodologie du commentaire philosophique}

\courssub{Objectifs de l'activité}

Cette activité d'intégration clôt la leçon de consolidation du commentaire philosophique. Elle vise à faire mobiliser, dans une situation unique et complète, l'ensemble des savoir-faire travaillés au cours de l'unité. À l'issue de l'atelier, l'élève doit être capable de :

\begin{itemize}
\item conduire une \cle{lecture analytique} d'un texte philosophique : repérer la thèse, dégager les concepts clés, identifier la structure argumentative ;
\item formuler une \cle{problématique} pertinente à partir de la tension interne du texte ;
\item construire un \cle{plan détaillé} en deux ou trois parties, cohérent avec la problématique ;
\item rédiger une \cle{introduction} complète (accroche, présentation du texte, problématique, annonce du plan) et un \cle{paragraphe de développement} intégrant un exemple concret camerounais ;
\item rédiger une \cle{conclusion justifiée}, qui répond à la problématique sans simple répétition ;
\item \cle{évaluer} la production d'un autre groupe à l'aide d'une grille critériée et \cle{justifier} chaque point attribué.
\end{itemize}

\courssub{Situation}

\textbf{Contexte.} Dans le cadre de la semaine de la philosophie organisée par le lycée technique de Douala-Bassa, le club de philosophie propose un concours interne de commentaire de texte. Le thème retenu cette année est : « La technique : libération ou aliénation de l'homme ? ». Chaque groupe de quatre élèves dispose de 2 heures (deux séances) pour produire un commentaire complet, qui sera ensuite évalué par un groupe concurrent.

\textbf{Le texte à commenter.}

\begin{quote}
« L'outil prolonge la main de l'homme, la machine prolonge son corps, et voici que l'ordinateur prolonge son intelligence. À chaque étape, l'homme s'est cru libéré d'une fatigue : fatigue du bras, fatigue du dos, fatigue de la mémoire. Mais chaque libération a un prix. Celui qui ne sait plus compter sans machine a perdu quelque chose de sa propre puissance. Celui qui confie toutes ses décisions à l'outil technique cesse peu à peu de décider lui-même. La technique donne le pouvoir, mais elle crée en même temps la dépendance. Elle est donc ambivalente : ce qui libère peut aussi asservir. Toute la question est de savoir si l'homme reste le maître de l'outil qu'il a créé, ou s'il devient le serviteur de son propre serviteur. »
\end{quote}

\textbf{Documents d'accompagnement.}

\begin{itemize}
\item \textbf{Document 1 (photographie) :} un atelier de mécanique au quartier Deïdo à Douala. Un jeune mécanicien répare le moteur d'une moto à l'aide d'outils simples ; grâce à son savoir-faire technique, il gagne sa vie et fait vivre sa famille.
\item \textbf{Document 2 (donnée chiffrée) :} selon une enquête récente, environ 85 \% des jeunes urbains camerounais de 15 à 24 ans possèdent un téléphone portable ; beaucoup déclarent passer plus de 5 heures par jour sur les réseaux sociaux, et certains reconnaissent ne plus pouvoir s'en passer.
\item \textbf{Document 3 (témoignage) :} « Avant, je connaissais tous les quartiers de Yaoundé par cœur. Maintenant, sans le GPS de mon téléphone, je me perds. Mais grâce à l'application de transport, je trouve plus de clients qu'avant. » (Moussa, moto-taximan, Yaoundé.)
\end{itemize}

\textbf{Tâche.} Rédiger un commentaire philosophique complet du texte (introduction, plan détaillé, un paragraphe de développement rédigé, conclusion), en intégrant au moins un des trois documents comme exemple, puis évaluer la production d'un autre groupe.

\courssub{Consignes}

\begin{enumerate}
\item \textbf{Lecture analytique.} Remplissez la grille de lecture : quel est le thème du texte ? Quelle est la thèse de l'auteur ? Quels sont les concepts clés ? Quelle est la structure du raisonnement (relevez les étapes) ?
\item \textbf{Problématisation.} Formulez une problématique sous forme d'une question qui exprime la tension centrale du texte. Comparez les problématiques des groupes et retenez la plus pertinente en justifiant votre choix.
\item \textbf{Construction du plan.} Proposez un plan détaillé en deux ou trois parties, avec les sous-parties et les exemples prévus pour chaque partie.
\item \textbf{Rédaction de l'introduction.} Rédigez l'introduction complète : phrase d'accroche, présentation du texte et de sa thèse, problématique, annonce du plan.
\item \textbf{Rédaction d'un paragraphe de développement.} Rédigez un paragraphe complet (explication d'une idée du texte, argument, exemple tiré d'au moins un document camerounais, analyse de l'exemple, transition).
\item \textbf{Conclusion.} Rédigez une conclusion qui répond explicitement à la problématique, bilan nuancé, et ouverture justifiée.
\item \textbf{Évaluation croisée.} Échangez votre production avec un autre groupe. Attribuez les points selon la grille critériée ci-dessous et \textbf{justifiez par écrit} chaque point accordé ou refusé.
\end{enumerate}

\begin{center}
\begin{tabularx}{\linewidth}{|X|c|}
\hline
\rowcolor{popblueL} \textbf{Critère d'évaluation} & \textbf{Points} \\
\hline
Thèse et concepts du texte correctement identifiés & /2 \\
\hline
Problématique claire, sous forme de question, fidèle au texte & /2 \\
\hline
Introduction complète (accroche, texte, problématique, plan) & /4 \\
\hline
Plan cohérent, parties équilibrées et ordonnées & /3 \\
\hline
Paragraphe : explication + argument + exemple camerounais analysé & /4 \\
\hline
Conclusion : réponse à la problématique, bilan, ouverture & /3 \\
\hline
Langue claire, vocabulaire philosophique correct & /2 \\
\hline
\textbf{Total} & \textbf{/20} \\
\hline
\end{tabularx}
\end{center}

\courssub{Ressources à mobiliser}

\begin{aretenir}[Les outils du commentaire philosophique]
\begin{itemize}
\item \textbf{Lecture analytique :} thème (de quoi parle le texte ?), thèse (que soutient l'auteur ?), concepts clés, structure argumentative.
\item \textbf{Problématique :} une question qui naît d'une tension ou d'un paradoxe du texte ; jamais une question fermée dont la réponse serait évidente.
\item \textbf{Introduction en quatre temps :} accroche $\rightarrow$ présentation du texte et de la thèse $\rightarrow$ problématique $\rightarrow$ annonce du plan.
\item \textbf{Paragraphe de développement :} idée expliquée $\rightarrow$ argument $\rightarrow$ exemple concret $\rightarrow$ analyse de l'exemple $\rightarrow$ transition.
\item \textbf{Conclusion :} réponse explicite à la problématique, bilan nuancé, ouverture en lien avec le sujet.
\item \textbf{Notions de l'unité :} la \cle{technique}, l'\cle{outil}, la \cle{libération}, l'\cle{aliénation}, la \cle{dépendance}, le \cle{progrès}.
\end{itemize}
\end{aretenir}

\begin{attention}[Pièges fréquents à éviter]
Ne paraphrasez pas le texte : commenter, c'est expliquer et discuter, pas répéter avec d'autres mots. Ne citez pas un exemple sans l'analyser : l'exemple doit servir l'argument. N'introduisez pas dans la conclusion une idée totalement nouvelle. Enfin, la problématique ne doit pas être une simple question de cours (« Qu'est-ce que la technique ? ») mais doit exprimer le problème propre au texte.
\end{attention}

\courssub{Démarche et corrigé commenté}

\begin{exoresolu}[Corrigé guidé de l'atelier d'intégration]
\textbf{Étape 1 — Grille de lecture analytique.}
\begin{itemize}
\item \emph{Thème :} la technique et ses effets sur la liberté humaine.
\item \emph{Thèse de l'auteur :} la technique est \cle{ambivalente} : elle libère l'homme de certaines fatigues, mais elle crée en même temps une dépendance qui peut l'asservir.
\item \emph{Concepts clés :} outil, machine, libération, dépendance, pouvoir, maître/serviteur.
\item \emph{Structure :} (a) constat : la technique prolonge les capacités humaines ; (b) renversement : chaque libération a un prix ; (c) conclusion provisoire : la technique est ambivalente ; (d) question finale : l'homme reste-t-il maître de son outil ?
\end{itemize}

\textbf{Étape 2 — Problématique.} La tension du texte oppose libération et asservissement par le même instrument. Problématique retenue : \emph{« La technique, qui libère l'homme de ses fatigues, ne risque-t-elle pas de le rendre dépendant au point de l'asservir ? »} Cette question est préférable à « Qu'est-ce que la technique ? », car elle reprend exactement le paradoxe du texte.

\textbf{Étape 3 — Plan détaillé.}
\begin{itemize}
\item \textbf{I. La technique libère l'homme.} (a) Elle prolonge le corps et l'intelligence (explication du texte) ; (b) exemple : le mécanicien de Deïdo (document 1) et le moto-taximan qui trouve plus de clients (document 3).
\item \textbf{II. Mais la technique rend dépendant.} (a) Celui qui délègue tout à la machine perd une part de sa puissance ; (b) exemples : le GPS qui fait perdre la mémoire des quartiers (document 3), les 5 heures quotidiennes sur les réseaux sociaux (document 2).
\item \textbf{III. Rester maître de l'outil.} (a) L'ambivalence n'est pas une fatalité : tout dépend de l'usage ; (b) condition : éducation et esprit critique.
\end{itemize}

\textbf{Étape 4 — Introduction rédigée.} « De la roue au téléphone portable, l'homme n'a cessé de fabriquer des outils pour étendre ses pouvoirs. Le texte soumis à notre étude traite de la technique et de ses effets sur la liberté humaine. L'auteur y soutient une thèse forte : la technique est ambivalente, car ce qui libère peut aussi asservir. En effet, si l'outil soulage le corps et l'esprit, il engendre une dépendance qui menace l'autonomie de celui qui s'y abandonne. Dès lors, une question s'impose : la technique, qui libère l'homme de ses fatigues, ne risque-t-elle pas de le rendre dépendant au point de l'asservir ? Pour y répondre, nous montrerons d'abord que la technique libère l'homme, puis qu'elle peut le rendre dépendant, et enfin que tout dépend de l'usage que l'homme en fait. »

\textbf{Étape 5 — Paragraphe de développement rédigé (partie II).} « L'auteur affirme que chaque libération a un prix : celui qui ne sait plus compter sans machine a perdu quelque chose de sa propre puissance. En d'autres termes, déléguer une faculté à l'outil, c'est risquer de la voir s'affaiblir en nous. L'argument est que la facilité technique dispense de l'effort qui entretient nos capacités. Le témoignage de Moussa, moto-taximan à Yaoundé, l'illustre parfaitement : lui qui connaissait tous les quartiers par cœur déclare se perdre sans le GPS de son téléphone. Cet exemple montre que la mémoire spatiale, non exercée, s'est atrophiée au profit de la machine. De même, les 85 \% de jeunes urbains camerounais équipés de téléphones portables passent parfois plus de 5 heures par jour sur les réseaux sociaux, au point de ne plus pouvoir s'en passer : l'outil de communication devient alors un lien de dépendance. Ainsi, la libération technique se paie d'une perte d'autonomie. Faut-il pour autant rejeter la technique ? »

\textbf{Étape 6 — Conclusion rédigée.} « Notre analyse a montré que la technique libère l'homme de la fatigue et élargit ses pouvoirs, mais qu'elle engendre aussi une dépendance capable de l'asservir. La réponse à notre problématique est donc nuancée : la technique n'est ni bonne ni mauvaise en elle-même ; c'est l'usage que l'homme en fait qui décide de sa valeur. Le mécanicien de Deïdo maîtrise son outil et s'émancipe ; l'usager qui s'abandonne aux écrans devient le serviteur de son serviteur. Reste alors une question : l'éducation peut-elle former des usagers capables de rester maîtres de leurs outils ? »

\textbf{Étape 7 — Évaluation croisée.} Chaque groupe attribue les points critère par critère et rédige une justification, par exemple : « Nous accordons 2/2 à la problématique car elle reprend le paradoxe libération/aliénation sous forme de question ; nous retirons 1 point au paragraphe car l'exemple du document 2 est cité mais insuffisamment analysé. »
\tcblower
\textbf{Justification de la démarche.} La grille de lecture garantit que le commentaire part du texte et non d'idées générales. La problématique est construite à partir de la tension interne (libération/dépendance), ce qui est la condition d'un plan dialectique. Le plan en trois parties suit la logique thèse--antithèse--dépassement, adaptée à un texte ambivalent. Le paragraphe rédigé respecte la chaîne explication--argument--exemple--analyse--transition : les documents camerounais ne sont pas décoratifs, ils prouvent. La conclusion répond à la question posée et l'ouverture prolonge la réflexion sans introduire de notion étrangère. L'évaluation croisée, enfin, développe l'esprit critique : justifier un point oblige à formuler les critères d'un bon commentaire.
\end{exoresolu}

\courssub{Bilan de l'activité}

Cet atelier a permis de mettre en évidence trois acquis essentiels. D'abord, le commentaire philosophique est une \cle{démarche ordonnée} : on ne rédige pas avant d'avoir lu analytiquement et problématisé. Ensuite, l'\cle{exemple concret} — ici les documents camerounais — n'a de valeur que s'il est analysé et rattaché à l'argument : c'est ce qui distingue le commentaire de la simple paraphrase. Enfin, l'\cle{évaluation croisée} a révélé les erreurs récurrentes (problématique trop générale, exemple plaqué, conclusion répétitive) et a montré que savoir évaluer un texte, c'est déjà savoir mieux le rédiger. L'élève est désormais prêt à affronter, dans les conditions du Probatoire et du Baccalauréat technique, l'épreuve du commentaire de texte.

\newpage

\coursec{Méthodes et exercices résolus}

\courssub{Les fondamentaux du commentaire philosophique}

\begin{methode}[Lecture analytique et prise de notes (15--20 min)]
\begin{enumerate}
    \item \textbf{Lecture 1 (découverte) :} lire le texte intégralement sans stylo pour saisir le ton, le thème général et la thèse principale.
    \item \textbf{Lecture 2 (analyse structurale) :} numéroter les paragraphes. Identifier les \textbf{mots-clés} (concepts philosophiques) et les \textbf{connecteurs logiques} (donc, or, mais, cependant, en effet) qui révèlent l'enchaînement des idées.
    \item \textbf{Lecture 3 (détail) :} souligner les \textbf{arguments} (prémisses), les \textbf{exemples} ou \textbf{références} (auteurs, faits), les \textbf{définitions} implicites ou explicites. Repérer les \textbf{oppositions} (thèse/antithèse) et les \textbf{nuances}.
    \item \textbf{Fiche de lecture :} construire un tableau synthétique : \emph{Paragraphe / Idée force / Argument(s) / Concept(s) clé(s) / Fonction dans le texte (thèse, objection, exemple, distinction, conclusion)}.
    \item \textbf{Formule à retenir :} \cle{Texte = Thèse + Arguments + Structure logique}. Ne jamais commenter mot à mot, mais idée par idée.
\end{enumerate}
\end{methode}

\begin{exoresolu}[Application : lecture analytique d'un extrait (type Bac)]
\textbf{Extrait :} « On confond souvent liberté et indépendance. L'indépendance est l'absence de contrainte extérieure ; la liberté est une puissance intérieure, une maîtrise de soi. L'esclave affranchi est indépendant ; mais s'il reste esclave de ses passions, il n'est pas libre. Inversement, Socrate en prison était libre, car il obéissait à sa raison, non à la peur. Être libre, ce n'est pas faire ce qu'on veut, mais vouloir ce qu'on fait, c'est-à-dire agir selon la raison. » \emph{(Texte d'inspiration spinoziste, d'après l'\emph{Éthique})}

\tcblower
\textbf{Solution détaillée :}

\noindent \textbf{1. Thème et thèse :}
\begin{itemize}
    \item \textbf{Thème :} la distinction entre liberté et indépendance.
    \item \textbf{Thèse :} la vraie liberté est autonomie (obéissance à sa propre raison), et non simple indépendance (absence de contrainte extérieure).
\end{itemize}

\noindent \textbf{2. Structure logique (découpage en trois moments) :}
\begin{center}
\begin{tabularx}{\linewidth}{|c|X|X|X|}\hline
\textbf{§} & \textbf{Idée force} & \textbf{Arguments / Exemples} & \textbf{Concepts clés} \\ \hline
1 & Distinction conceptuelle & Définition par genre et différence spécifique : indépendance = absence de contrainte extérieure ; liberté = puissance intérieure, maîtrise de soi. & Indépendance, liberté, contrainte, puissance, maîtrise. \\ \hline
2 & Illustration par l'opposition (esclave / Socrate) & Exemple 1 : l'esclave affranchi est indépendant mais non libre (esclave de ses passions). Exemple 2 : Socrate en prison est non indépendant mais libre (il obéit à la raison). & Affranchissement, passion, raison, autonomie, hétéronomie. \\ \hline
3 & Définition positive (paradoxe) & « Vouloir ce qu'on fait » contre « faire ce qu'on veut ». Agir selon la raison = définition de la liberté. & Volonté, action, raison, détermination interne. \\ \hline
\end{tabularx}
\end{center}

\noindent \textbf{3. Problématique émergente :} \emph{la liberté se réduit-elle à l'absence de chaînes extérieures ou exige-t-elle une contrainte intérieure, celle de la raison ?}

\noindent \textbf{4. Conclusion de l'analyse :} le texte opère un \textbf{renversement paradoxal} : la contrainte physique (la prison) ne nuit pas à la liberté si l'esprit est autonome ; l'absence de contrainte physique ne garantit pas la liberté si l'esprit est asservi aux passions.
\end{exoresolu}

\begin{methode}[Problématiser et rédiger l'introduction (10--15 min)]
\begin{enumerate}
    \item \textbf{Accroche (1 phrase) :} citation brève, fait d'actualité camerounaise, ou énoncé du sens commun à renverser. \emph{Exemple : « Au Cameroun, on revendique souvent la "liberté" comme le fait de ne pas avoir de chef ; or le texte montre... »}
    \item \textbf{Présentation du texte :} auteur (si connu), œuvre, date, contexte (1 phrase). \emph{« Dans cet extrait de l'\emph{Éthique} (1677), Spinoza... »}
    \item \textbf{Résumé de la thèse :} reformuler la thèse de l'auteur en une phrase claire (pas de citation brute).
    \item \textbf{Problématique :} une question \textbf{ouverte}, \textbf{philosophique}, qui met en tension la thèse du texte et une objection possible. Structure type : \emph{« Dans quelle mesure [thèse du texte]... alors que / sachant que [objection ou limite]... ? »}
    \item \textbf{Annonce du plan :} \emph{« Nous verrons d'abord que... (explication), puis nous montrerons que... (discussion, nuance), pour enfin comprendre que... (dépassement, synthèse). »}
    \item \textbf{Réflexe Bac :} l'introduction se rédige \textbf{en dernier} au brouillon, mais se place \textbf{en premier} sur la copie. Elle ne doit pas dépasser 15 à 20 lignes.
\end{enumerate}
\end{methode}

\begin{exoresolu}[Rédaction d'une introduction type Bac]
\textbf{Sujet :} même extrait que précédemment (texte d'inspiration spinoziste, liberté et indépendance).

\tcblower
\textbf{Proposition d'introduction rédigée :}

« "L'homme est né libre, et partout il est dans les fers", écrit Rousseau. Cette phrase résonne particulièrement dans nos sociétés camerounaises, où l'aspiration à l'indépendance politique ou économique est souvent confondue avec l'accès à la liberté véritable. C'est à cette confusion que s'attaque Spinoza dans l'\emph{Éthique} (1677), œuvre majeure du rationalisme classique. Dans l'extrait proposé, le philosophe montre que l'indépendance, entendue comme absence de contrainte extérieure, ne saurait se confondre avec la liberté, définie comme une puissance intérieure de maîtrise de soi par la raison. Il illustre ce paradoxe par l'opposition entre l'esclave affranchi, indépendant mais asservi à ses passions, et Socrate en prison, physiquement captif mais spirituellement libre.

\textbf{Dès lors, la liberté consiste-t-elle simplement à briser ses chaînes extérieures, ou exige-t-elle, plus radicalement, la maîtrise de ses propres déterminations intérieures ?}

Pour répondre, nous expliquerons d'abord en quoi la distinction entre indépendance (contrainte extérieure) et liberté (maîtrise intérieure) repose sur une redéfinition de la puissance d'agir (I). Nous discuterons ensuite les limites d'une liberté réduite à la seule rationalité, au risque d'exclure ceux que la raison ne guide pas encore (II). Nous montrerons enfin que cette conception engage une éthique de la libération par la connaissance, seule voie vers une autonomie effective (III). »
\end{exoresolu}

\begin{methode}[Construire le développement : expliquer -- illustrer -- discuter (40--50 min)]
\begin{enumerate}
    \item \textbf{Structure dialectique (trois parties, deux ou trois sous-parties chacune) :}
    \begin{itemize}
        \item \textbf{Partie I : EXPLIQUER (fidélité au texte).} Reprendre le fil de la démonstration de l'auteur. \emph{Mots-clés :} « Spinoza montre que... », « Il distingue... », « L'argument repose sur... ». \textbf{Aucun avis personnel.}
        \item \textbf{Partie II : DISCUTER, NUANCER (esprit critique).} Tester la solidité de la thèse. Apporter \textbf{une objection précise} (auteur adverse, contre-exemple, limite logique). \emph{Mots-clés :} « Cependant... », « On peut objecter que... (Kant, Sartre, Marx...) », « Cette thèse suppose que... ».
        \item \textbf{Partie III : DÉPASSER, SYNTHÉTISER (approfondissement).} Résoudre la tension. Montrer comment l'objection enrichit ou corrige la thèse. Ouvrir sur une implication éthique, politique ou anthropologique. \emph{Mots-clés :} « Au terme de cette analyse... », « Il apparaît que... », « Cela implique que... ».
    \end{itemize}
    \item \textbf{Règle d'or du paragraphe (MEI) :} \textbf{M} -- idée directrice (1 phrase) ; \textbf{E} -- explication, citation courte ou exemple concret (le Cameroun est bienvenu) ; \textbf{I} -- interprétation, lien avec la problématique, transition.
    \item \textbf{Connecteurs obligatoires :} \emph{en effet, plus précisément, c'est pourquoi, toutefois, néanmoins, par conséquent, en définitive.}
\end{enumerate}
\end{methode}

\begin{exoresolu}[Rédaction guidée de la Partie I (explication)]
\textbf{Consigne :} rédiger la première partie du développement : « La distinction entre indépendance (extériorité) et liberté (intériorité) ».

\tcblower
\textbf{Proposition de rédaction (niveau Probatoire technique) :}

\textbf{I. La liberté comme maîtrise de soi : rupture avec le sens commun de l'indépendance}

\textbf{A. L'indépendance : une définition négative et extérieure}

Le texte opère d'emblée une \textbf{distinction conceptuelle rigoureuse}. L'indépendance se définit \emph{négativement} : c'est l'\emph{« absence de contrainte extérieure »}. Elle relève de l'\textbf{extériorité} : l'indépendant n'est pas empêché par autrui. C'est une condition \textbf{juridique ou physique} (statut d'affranchi, absence de chaînes). Or, le texte montre que cette condition est \textbf{insuffisante} : l'esclave affranchi jouit d'une indépendance formelle, mais s'il demeure « esclave de ses passions », il subit une contrainte \emph{intérieure} qui le prive de liberté. L'indépendance est donc une liberté purement formelle, qui ne garantit pas l'exercice effectif de l'action.

\textbf{B. La liberté : une définition positive et intérieure}

À l'inverse, la liberté est une \textbf{puissance intérieure} (\emph{« une puissance intérieure, une maîtrise de soi »}). Elle se définit \emph{positivement} comme \textbf{autonomie} : se donner à soi-même sa loi (\emph{auto-nomos}). L'exemple de Socrate est décisif : en prison (dépendance physique maximale), il obéit à sa raison (autonomie spirituelle maximale). La liberté n'est donc pas le caprice (\emph{« faire ce qu'on veut »}), mais la \textbf{détermination rationnelle de soi} (\emph{« vouloir ce qu'on fait »}). C'est le passage de la \textbf{servitude} (subir des causes extérieures : les passions) à la \textbf{liberté} (agir selon sa propre raison).

\textbf{\textit{Transition :} si cette distinction fondatrice permet de penser la liberté comme une conquête (la sagesse) et non comme un donné (le statut social), elle soulève une difficulté majeure : que devient la liberté de celui qui n'a pas encore accédé à la raison ?}
\end{exoresolu}

\begin{methode}[Rédiger la conclusion et relire (10 min)]
\begin{enumerate}
    \item \textbf{Bilan synthétique (3 à 4 lignes) :} reprendre la problématique. Répondre de manière \textbf{affirmative mais nuancée}. \emph{« Au terme de cette analyse, il apparaît que la liberté ne se réduit pas à l'indépendance... elle est conquête rationnelle... »}
    \item \textbf{Ouverture (1 à 2 lignes) :} prolonger la réflexion sans ouvrir un nouveau dossier. Lien avec une autre notion du programme (droit, politique, morale, technique) ou avec l'actualité. \emph{« Cette conception éclaire nos luttes actuelles : l'indépendance nationale n'a de sens que si elle s'accompagne d'une libération des mentalités... »}
    \item \textbf{Relecture active (check-list) :}
    \begin{itemize}
        \item \textbf{Respect du texte :} ai-je cité le texte (entre guillemets) pour l'expliquer ? Ai-je évité le hors-texte en Partie I ?
        \item \textbf{Cohérence :} le plan annoncé est-il respecté ? Les transitions sont-elles logiques ?
        \item \textbf{Vocabulaire :} les concepts philosophiques sont-ils utilisés correctement (autonomie, hétéronomie, servitude, passion, raison) ?
        \item \textbf{Orthographe et grammaire :} accords, majuscules (auteur, œuvre), ponctuation.
        \item \textbf{Présentation :} sauts de ligne clairs (introduction, I, A, B, transition, II, III, conclusion), paragraphes en retrait.
    \end{itemize}
\end{enumerate}
\end{methode}

\begin{exoresolu}[Rédaction de la conclusion et ouverture camerounaise]
\textbf{Sujet :} même extrait (texte d'inspiration spinoziste).

\tcblower
\textbf{Proposition de conclusion :}

« En définitive, le commentaire de cet extrait permet d'affirmer que \textbf{la liberté ne saurait se confondre avec l'indépendance}. Là où l'indépendance n'est qu'une \textbf{condition négative et extérieure} (l'absence d'obstacles), la liberté est une \textbf{réalité positive et intérieure} : la maîtrise de soi par la raison. Le texte nous enseigne que l'on peut être libre dans une prison (Socrate) et esclave en pleine ville (l'affranchi dominé par ses désirs).

Cette analyse invite à revisiter l'histoire de notre pays : les \textbf{indépendances} du Cameroun (1\ier{} janvier 1960 pour le Cameroun oriental, 1\ier{} octobre 1961 pour le Cameroun occidental) ont brisé les chaînes coloniales extérieures. Mais, comme le suggère le texte, cette indépendance nationale reste inachevée si elle ne s'accompagne pas d'une \textbf{libération intérieure} : éducation à la raison critique, lutte contre les "passions" collectives (tribalisme, corruption, résignation), construction d'une citoyenneté active. \textbf{La vraie souveraineté, comme la vraie liberté, se conquiert par la connaissance.} »

\textbf{Check-list de relecture validée :} texte cité (Socrate, passions) ; plan respecté (I. Expliquer, II. Discuter, III. Synthétiser) ; concepts employés (autonomie, hétéronomie, servitude) ; ouverture contextuelle (indépendances du Cameroun).
\end{exoresolu}

\begin{methode}[Gestion du temps et stratégie pour l'épreuve (4 h)]
\begin{enumerate}
    \item \textbf{0 h 00 -- 0 h 20 :} choix du sujet (commentaire ou dissertation). \textbf{Conseil :} si le texte est court et clair, choisir le commentaire ; si le texte est obscur ou aride, préférer la dissertation (lorsque le choix est offert).
    \item \textbf{0 h 20 -- 0 h 45 :} \textbf{brouillon intensif.} Lecture analytique (trois lectures), tableau d'analyse, problématique, plan détaillé (idées directrices, citations, exemples, auteurs pour la discussion).
    \item \textbf{0 h 45 -- 2 h 45 :} \textbf{rédaction au propre (2 h).} Introduction (15 min) ; Partie I (35 min) ; Partie II (35 min) ; Partie III (25 min) ; conclusion (10 min).
    \item \textbf{2 h 45 -- 3 h 00 :} \textbf{relecture finale} (voir la méthode précédente). Ne pas recopier le brouillon : \textbf{écrire directement au propre} en suivant le brouillon structuré.
    \item \textbf{Matériel :} deux stylos (bleu ou noir), correcteur (si autorisé), montre, dictionnaire (si autorisé). Pas de téléphone.
    \item \textbf{Présentation de la copie :} saut de ligne \textbf{obligatoire} entre introduction, I, A, B, transition, II, III et conclusion. Numéroter les pages.
\end{enumerate}
\end{methode}

\begin{exoresolu}[Simulation de planning de l'épreuve (4 h)]
\textbf{Situation :} sujet de commentaire sur un texte de vingt lignes (par exemple un extrait sur la technique et la condition humaine).

\tcblower
\textbf{Chronogramme réaliste pour un élève de Première technique :}

\begin{center}
\begin{tabularx}{\linewidth}{|c|X|X|}\hline
\textbf{Tranche horaire} & \textbf{Activité précise (brouillon / propre)} & \textbf{Objectif / pièges à éviter} \\ \hline
0 h 00 -- 0 h 10 & \textbf{Choix et lecture 1} (découverte) & Lire les deux sujets. Choisir celui dont le vocabulaire est maîtrisé. Ne pas choisir « par facilité apparente ». \\ \hline
0 h 10 -- 0 h 35 & \textbf{Brouillon : analyse structurée} & Tableau § / idée / concepts. \textbf{Repérer la thèse} et \textbf{l'enchaînement des idées}. \textit{Piège :} passer une heure sur le brouillon. \\ \hline
0 h 35 -- 0 h 50 & \textbf{Brouillon : problématique et plan} & Formuler la question. Tracer le plan dialectique (I, II, III avec A, B). Chercher \textbf{un auteur adverse} (Partie II) et \textbf{un exemple concret camerounais}. \\ \hline
0 h 50 -- 1 h 10 & \textbf{Propre : introduction} & Rédiger au net. Soigner l'accroche et l'annonce du plan. \\ \hline
1 h 10 -- 1 h 55 & \textbf{Propre : Partie I (explication)} & Suivre le texte pas à pas. Citer court (guillemets). Expliquer les concepts. \textit{Piège :} faire du hors-texte. \\ \hline
1 h 55 -- 2 h 40 & \textbf{Propre : Partie II (discussion)} & Objection argumentée (par exemple Marx : la liberté exige des conditions matérielles, pas seulement la raison). Nuancer la thèse du texte. \\ \hline
2 h 40 -- 3 h 10 & \textbf{Propre : Partie III (synthèse) et conclusion} & Dépasser l'opposition. Ouverture sur le programme (technique, politique). Bilan puis ouverture. \\ \hline
3 h 10 -- 3 h 30 & \textbf{Relecture globale} & Check-list (orthographe, structure, concepts, respect du texte). Corriger \textbf{proprement} (raturer d'un seul trait). \\ \hline
\end{tabularx}
\end{center}

\textbf{Bilan :} 50 min de choix et de brouillon, 2 h 20 de rédaction, 20 min de relecture, soit 3 h 30 effectives, avec une marge de 30 min pour les imprévus, la fatigue ou l'approfondissement.
\end{exoresolu}

\begin{methode}[Exploiter un exemple concret camerounais (contextualisation)]
\begin{enumerate}
    \item \textbf{Pourquoi ?} Le programme exige d'« appliquer les notions de l'unité à des situations concrètes de la vie courante ». Cela prouve la \textbf{pertinence philosophique} de la réflexion et ancre l'abstraction dans le réel.
    \item \textbf{Comment choisir l'exemple ?}
    \begin{itemize}
        \item \textbf{Thème de la liberté :} indépendances de 1960--1961 et libertés publiques aujourd'hui ; éducation (contrainte scolaire et autonomie).
        \item \textbf{Thème de la technique :} téléphone portable (libération de la communication ou asservissement aux écrans) ; agriculture moderne (machine et savoir paysan).
        \item \textbf{Thème du droit et de la politique :} décentralisation (autonomie locale et unité nationale) ; justice traditionnelle et justice étatique.
    \end{itemize}
    \item \textbf{Comment l'intégrer ?} \textbf{Jamais en « bloc » à la fin.} L'exemple sert d'\textbf{illustration} dans un paragraphe (Partie I pour expliquer, Partie II pour objecter, Partie III pour synthétiser).
    \item \textbf{Formule type :} \emph{« Cette analyse trouve une illustration frappante dans le contexte camerounais : [fait précis] montre que [concept]... »}
\end{enumerate}
\end{methode}

\begin{exoresolu}[Intégration d'exemples camerounais dans le développement]
\textbf{Sujet :} toujours le même extrait (liberté et indépendance). Intégrer deux exemples distincts.

\tcblower
\textbf{Exemple 1 (Partie I -- illustration de la thèse : indépendance $\neq$ liberté) :}

« L'exemple de l'esclave affranchi "esclave de ses passions" trouve un écho saisissant dans la situation de \textbf{jeunes diplômés en recherche d'emploi à Douala ou à Yaoundé}. Ils jouissent d'une \textbf{indépendance juridique} (citoyens, diplômés, sans maître), mais beaucoup subissent une \textbf{servitude intérieure} : la passion du "gain facile" (arnaques, jeux de hasard), la résignation ("cela ne sert à rien"), ou la dépendance aux réseaux d'influence pour survivre. Ils sont "indépendants" sur le papier, mais non "libres" au sens du texte, car ils ne sont pas maîtres de leur existence : ils subissent la conjoncture sans la comprendre ni la transformer par la raison. »

\textbf{Exemple 2 (Partie II -- objection et nuance : les conditions matérielles de la liberté) :}

« Cependant, \textbf{Marx} objecterait que la "maîtrise de soi" est un luxe impossible sans \textbf{indépendance matérielle}. Prenons l'exemple des \textbf{paysans du Nord-Cameroun} confrontés aux aléas climatiques. Ils possèdent peut-être un savoir agricole ancestral, mais l'absence d'infrastructures (retenues d'eau, routes, crédit agricole) -- c'est-à-dire un manque d'\textbf{indépendance économique et technique} -- les réduit à la survie. On ne peut "vouloir ce qu'on fait" (liberté) si les conditions matérielles de l'action font défaut. La liberté définie par le texte risque de rester une "liberté de sages" (comme Socrate) si elle ignore les conditions sociales de son exercice. »

\textbf{Qualité attendue au Bac :} des exemples datés, localisés, reliés aux concepts (servitude, maîtrise de soi, conditions sociales), utilisés pour \textbf{argumenter} et non pour « décorer » la copie.
\end{exoresolu}

\begin{attention}[Les cinq erreurs qui coûtent des points au Bac]
\begin{enumerate}
    \item \textbf{La paraphrase (hors-sujet majeur) :} raconter le texte avec ses propres mots au lieu de l'\textbf{analyser} (expliquer la logique, les concepts, les enjeux). \textit{Sanction :} note plafonnée même si la copie est bien écrite.
    \item \textbf{L'introduction « catalogue » ou « vague » :} pas de problématique réelle (question fermée, hors-sujet ou simple résumé) ; plan non annoncé, ou plan « thématique » (I. La liberté, II. L'indépendance) au lieu d'un plan \textbf{dialectique} (I. Thèse, II. Limites, III. Synthèse).
    \item \textbf{L'oubli du texte en Partie I :} développer \textbf{son propre cours} sur la liberté sans s'appuyer sur les arguments, citations et distinctions \textbf{spécifiques de l'extrait}. Le commentaire n'est pas une dissertation sur le thème.
    \item \textbf{La « discussion » sans auteur ni argument :} écrire « on n'est pas d'accord » ou « c'est trop idéaliste » sans citer un \textbf{philosophe précis} (Kant, Sartre, Marx, Bergson...) et son \textbf{argument structuré}. L'avis personnel non fondé ne rapporte aucun point.
    \item \textbf{La conclusion « bâclée » ou « nouveau développement » :} pas de réponse à la problématique, pas de bilan, ouverture absurde (par exemple : « pour finir, parlons de la mort »). Pire encore : introduire un \textbf{nouvel argument majeur} qui n'a pas été traité dans le développement.
\end{enumerate}
\end{attention}

\newpage

\coursec{Exercices}

\courssub{Je vérifie mes connaissances}

\begin{exercice}[QCM : Nature du commentaire philosophique]
Parmi les propositions suivantes, laquelle définit le mieux l'objectif principal d'un commentaire philosophique de texte ?
\begin{enumerate}
\item Résumer le texte en utilisant ses propres mots.
\item Exposer son opinion personnelle sur le sujet traité par l'auteur.
\item Expliquer la pensée de l'auteur en dégageant sa thèse, ses arguments et sa structure, puis en en discutant la portée.
\item Recenser toutes les références culturelles et historiques présentes dans le texte.
\end{enumerate}
\end{exercice}

\begin{exercice}[Vrai / Faux justifié : Étapes de la méthode]
Indiquez si les affirmations suivantes sont vraies ou fausses. Justifiez brièvement chaque réponse.
\begin{enumerate}
\item La problématique doit être posée dès la première phrase de l'introduction, avant la présentation du texte.
\item L'explication du texte consiste à le paraphraser phrase par phrase en changeant le vocabulaire.
\item La discussion (ou évaluation critique) porte sur la validité des arguments de l'auteur et non sur son style littéraire.
\item La conclusion doit obligatoirement ouvrir sur une autre référence philosophique ou une citation célèbre.
\end{enumerate}
\end{exercice}

\begin{exercice}[Définitions précises]
Donnez une définition rigoureuse (2 à 3 lignes) pour chacun des termes méthodologiques suivants :
\begin{enumerate}
\item La thèse d'un texte philosophique.
\item Le thème (ou sujet) d'un texte philosophique.
\item La problématique.
\item L'argument (dans le cadre d'un raisonnement philosophique).
\end{enumerate}
\end{exercice}

\begin{exercice}[Repérage structurel : Les moments du raisonnement]
Un texte philosophique classique propose souvent un raisonnement en trois temps (moment dialectique ou analytique). Nommez ces trois moments types et décrivez en une phrase la fonction de chacun dans l'économie du texte.
\end{exercice}

\courssub{Je m'entraîne}

\begin{exercice}[Analyse guidée : Texte sur la liberté (extrait type Bac)]
\textbf{Consigne :} Lisez l'extrait suivant (environ 15 lignes) puis répondez aux questions.

\textit{« L'homme naît libre, et partout il est dans les fers. Celui qui se croit le maître des autres ne cesse pas d'être plus esclave qu'eux. Comment ce changement s'est-il fait ? Je l'ignore. Qu'est-ce qui peut le rendre légitime ? C'est là ce que je crois pouvoir résoudre. [...] L'ordre social est un droit sacré qui sert de base à tous les autres. Cependant ce droit ne vient pas de la nature, il est donc fondé sur des conventions. »} \\
\textit{(J.-J. Rousseau, \emph{Du contrat social}, Livre I, chap. 1 et 6)}

\begin{enumerate}
\item Dégagez le \cle{thème} général de cet extrait.
\item Formulez la \cle{thèse} que l'auteur défend ici.
\item Identifiez les \textbf{trois moments} du raisonnement rousseauiste présents dans cet extrait (constat, questionnement, annonce de solution/fondation).
\item Formulez la \cle{problématique} à laquelle ce texte répond (sous forme de question unique et précise).
\end{enumerate}
\end{exercice}

\begin{exercice}[Rédaction d'une introduction complète]
\textbf{Sujet :} « Le travail n'est-il qu'une contrainte ? »
Rédigez une introduction complète en respectant scrupuleusement les \textbf{quatre temps canoniques} :
\begin{enumerate}
\item L'accroche (amorce) : contextualisation philosophique ou référence à une expérience commune.
\item La présentation du sujet : définition des termes clés (\emph{travail}, \emph{contrainte}) et mise en tension.
\item La problématique : reformulation interrogative du sujet (autre que le sujet lui-même).
\item L'annonce du plan : présentation synthétique des deux ou trois grandes parties du développement envisagé.
\end{enumerate}
\end{exercice}

\begin{exercice}[Chasse aux erreurs méthodologiques]
Le paragraphe suivant est extrait d'une copie d'élève tentant d'expliquer un argument d'Aristote sur l'amitié. Identifiez \textbf{trois fautes méthodologiques distinctes} (parmi : paraphrase, hors-sujet, opinion non justifiée, contresens, absence d'articulation) et proposez une correction pour chacune.

\textit{« Aristote dit que l'amitié est nécessaire à la vie. C'est vrai car sans amis on est triste. Moi, par exemple, mon meilleur ami s'appelle Paul et on joue à la PlayStation ensemble, donc c'est nécessaire. D'ailleurs, tout le monde aime avoir des amis, c'est normal. Aristote a raison sur toute la ligne car l'amitié c'est le bonheur. Donc il faut se faire des amis pour être heureux. »}
\end{exercice}

\begin{exercice}[Paragraphe de développement : Explication et illustration camerounaise]
\textbf{Consigne :} Rédigez \textbf{un unique paragraphe structuré} (environ 15-20 lignes) qui \textbf{explique} l'argument suivant de John Locke : « La propriété privée naît du mélange du travail de l'homme avec les ressources communes de la nature » (\emph{Second Traité du gouvernement civil}, §27).
Votre paragraphe doit contenir :
\begin{enumerate}
\item La reformulation claire de l'argument de l'auteur.
\item L'analyse du concept clé (\emph{mélange}, \cle{appropriation}).
\item Une \textbf{illustration concrète prise de la vie courante au Cameroun} (ex. : un cultivateur à Bafia, un artisan menuisier à Douala, un développeur d'application à Yaoundé).
\item Une phrase de transition ou de mini-conclusion sur la portée de cet argument.
\end{enumerate}
\end{exercice}

\courssub{Je me prépare au Bac}

\begin{exercice}[Sujet complet type Probatoire / Bac Technique : Commentaire guidé — Le Devoir]
\textbf{Durée : 2 heures} \hfill \textbf{Coefficient : 3}

\textbf{Texte à commenter :}
\textit{« Le devoir n'est pas une contrainte extérieure qui s'imposerait à nous par la force. Il est la loi que notre raison nous donne à nous-mêmes. Agir par devoir, c'est agir par respect pour la loi morale que nous nous sommes prescrite en tant qu'êtres raisonnables. L'action a alors sa valeur morale non pas dans le but qu'elle vise, mais dans la maxime selon laquelle elle est entreprise. Ce n'est pas parce que l'action est utile ou agréable qu'elle est bien, mais parce qu'elle est faite par devoir. L'autonomie de la volonté est le principe suprême de la morale. »} \\
\textit{(D'après I. Kant, \emph{Fondements de la métaphysique des mœurs})}

\textbf{Grille de lecture fournie (à utiliser pour structurer la copie) :}
\begin{enumerate}
\item \textbf{Analyse du texte (Explication) :}
      \begin{itemize}
      \item Thème et Thèse.
      \item Vocabulaire clé : \emph{contrainte}, \emph{loi morale}, \emph{autonomie}, \emph{maxime}, \emph{respect}.
      \item Structure du raisonnement (3 moments) : Opposition contrainte/devoir $\rightarrow$ Définition de l'autonomie $\rightarrow$ Critère de la valeur morale (désintéressement).
      \end{itemize}
\item \textbf{Problématique :} En quoi le devoir, loin d'être une contrainte hétéronome, est-il l'expression de la liberté rationnelle (autonomie) ?
\item \textbf{Discussion (Évaluation critique) :}
      \begin{itemize}
      \item La thèse kantienne : la rigueur du devoir pur (force : universalité, dignité ; limite : formalisme, abstraction du concret).
      \item Objections / Nuances : Le devoir social (Durkheim), l'utilitarisme (Mill), la psychologie morale (inclination vs devoir), le contexte camerounais (devoirs traditionnels vs devoirs modernes).
      \end{itemize}
\end{enumerate}

\textbf{Travail demandé :}
Rédigez le commentaire complet en suivant la méthode (Introduction, Développement en 2 ou 3 parties équilibrées, Conclusion). L'explication du texte doit être fondues dans le développement, pas traitée à part.
\end{exercice}

\begin{exercice}[Comparaison de problématiques : Choix et justification]
\textbf{Texte de référence (non reproduit, connu du cours) :} Extrait de \emph{L'Être et le Néant} de Sartre sur la \cle{mauvaise foi} (le garçon de café qui « joue » à être un garçon de café).

Deux élèves proposent la problématique suivante pour ce même texte :
\begin{itemize}
\item \textbf{Problématique A :} « Qu'est-ce que la mauvaise foi chez Sartre ? »
\item \textbf{Problématique B :} « La mauvaise foi est-elle une fuite inévitable devant l'angoisse de la liberté ou une possibilité structurelle de la conscience ? »
\end{itemize}

\textbf{Consigne :} En \textbf{cinq lignes maximum}, comparez ces deux formulations et justifiez laquelle est la plus pertinente pour un commentaire philosophique de niveau Bac Technique. Utilisez les critères : \emph{ouverture au débat}, \emph{ancrage dans le texte}, \emph{portée philosophique}.
\end{exercice}

\begin{exercice}[Sujet de dissertation corrélée (Entraînement Bac) — Optionnel si temps restant]
\textbf{Sujet :} « Obéir, est-ce renoncer à sa liberté ? »

\textbf{Consigne :} Traitez ce sujet en dissertation (Introduction, Développement en 3 parties, Conclusion). Vous devez obligatoirement :
\begin{enumerate}
\item Mobiliser la distinction kantienne entre \emph{autonomie} et \emph{hétéronomie} (vue dans l'exercice 1).
\item Mobiliser la distinction entre \emph{liberté naturelle} et \emph{liberté civile} (Rousseau, \emph{Contrat social}).
\item Donner un exemple concret d'obéissance légitime dans le contexte camerounais actuel (code de la route, règlement intérieur d'un lycée technique, loi sur la protection de l'environnement).
\end{enumerate}
\end{exercice}

\newpage

\coursec{Corrigés des exercices}

\coursec{Consolidation des acquis de la méthodologie du commentaire philosophique}

\courssub{Rappel des fondamentaux : qu'est-ce que commenter un texte philosophique ?}

\begin{exoresolu}[QCM : Nature du commentaire philosophique]
Énoncé : Parmi les propositions suivantes, laquelle définit le mieux l'activité de commentaire philosophique ?
\begin{enumerate}
\item Résumer le texte de l'auteur.
\item Exprimer son opinion personnelle sur le sujet.
\item Expliquer le texte (thèse, arguments, structure) puis en discuter la portée et la validité.
\item Recenser les références savantes liées au thème.
\end{enumerate}
\tcblower
\textbf{Bonne réponse : proposition 3.}

\begin{itemize}
\item \textbf{Proposition 1 (fausse) :} résumer le texte n'est qu'une étape préparatoire (la lecture analytique). Le résumé ne dégage ni la structure argumentative, ni la portée philosophique ; il reste à la surface du texte.
\item \textbf{Proposition 2 (fausse) :} le commentaire n'est pas un essai d'opinion. Dire « je pense que... » sans s'appuyer sur le texte est une faute méthodologique majeure : le commentaire part du texte et y revient constamment.
\item \textbf{Proposition 3 (juste) :} c'est la définition canonique : \emph{expliquer} (thèse, arguments, structure), puis \emph{discuter} (portée, validité, limites). Le commentaire est à la fois fidélité au texte et distance critique.
\item \textbf{Proposition 4 (fausse) :} l'érudition n'est pas la philosophie. Recenser des références ne prouve rien et ne fait pas progresser la compréhension de la pensée de l'auteur.
\end{itemize}

\textbf{Point de vigilance :} au Bac, le correcteur vérifie dès l'introduction que le candidat a compris que le commentaire est une \cle{explication de texte} et non un résumé ni une dissertation déguisée.
\end{exoresolu}

\begin{exoresolu}[Vrai / Faux justifié : Étapes de la méthode]
Énoncé : Juste ou faux ? Justifiez brièvement.
\begin{enumerate}
\item La problématique se pose avant la présentation du texte dans l'introduction.
\item Expliquer un texte, c'est le paraphraser avec ses propres mots.
\item La discussion porte sur la validité des arguments et non sur le style littéraire.
\item L'ouverture (citation, perspective) est obligatoire en conclusion.
\end{enumerate}
\tcblower
\begin{enumerate}
\item \textbf{Faux.} L'introduction suit un ordre précis : accroche, présentation du texte (auteur, \oe uvre, thème), \emph{puis} problématique, puis annonce du plan. Poser la problématique avant de présenter le texte revient à interroger un texte que le lecteur ne connaît pas encore : la question serait suspendue dans le vide.
\item \textbf{Faux.} La paraphrase est le défaut le plus sanctionné. Expliquer, c'est dégager la \emph{structure} du raisonnement (thèse, arguments, articulations), définir les concepts clés et montrer comment l'auteur prouve ce qu'il avance — non répéter le texte avec d'autres mots.
\item \textbf{Vrai.} La discussion porte sur la \cle{validité} des arguments (sont-ils solides ? universels ? quelles limites ?) et sur la portée de la thèse. Le style littéraire relève du commentaire littéraire, pas du commentaire philosophique, qui juge des idées et non de l'élégance de la plume.
\item \textbf{Faux.} L'ouverture est facultative, jamais obligatoire. Une conclusion doit d'abord \emph{bilan} (rappel de la thèse et du chemin parcouru) et \emph{réponse} à la problématique. Une citation célèbre plaquée en fin de copie sans lien avec le texte est une faute de goût méthodologique.
\end{enumerate}

\textbf{Point de vigilance :} « justifier brièvement » exige une raison, pas seulement « vrai » ou « faux » : une réponse non justifiée ne vaut aucun point.
\end{exoresolu}

\begin{exoresolu}[Définitions précises : Thèse, Thème, Problématique, Argument]
Énoncé : Définissez rigoureusement les quatre notions suivantes en les distinguant clairement : Thèse, Thème (Sujet), Problématique, Argument.
\tcblower
\begin{enumerate}
\item \textbf{La thèse :} c'est la position doctrinale que l'auteur défend dans le texte, la réponse qu'il apporte au problème posé. Elle s'énonce en une phrase affirmative claire (ex. : pour Kant, « la valeur morale de l'action réside dans le devoir, non dans ses résultats »).
\item \textbf{Le thème (ou sujet) :} c'est le domaine général, la notion ou la question dont traite le texte (la liberté, le devoir, l'État...). Le thème est \emph{commun} à plusieurs auteurs ; la thèse est \emph{propre} à l'auteur commenté. Ne jamais confondre les deux.
\item \textbf{La problématique :} c'est la question précise, formulée interrogativement, à laquelle le texte répond et qui guide tout le commentaire. Elle doit être unique, philosophique (ouvrant un débat) et fidèle au texte.
\item \textbf{L'argument :} c'est une raison avancée pour établir ou justifier la thèse. Un raisonnement philosophique articule plusieurs arguments (logiques, d'expérience, par l'absurde...) qui convergent vers la conclusion de l'auteur.
\end{enumerate}

\textbf{Point de vigilance :} une définition rigoureuse distingue le terme de ses voisins : thème $\neq$ thèse $\neq$ problématique. Cette distinction est la base de toute introduction réussie.
\end{exoresolu}

\begin{exoresolu}[Repérage structurel : Les moments du raisonnement]
Énoncé : Quels sont les trois moments types d'un raisonnement philosophique classique ? Quelle est leur fonction dans l'économie du texte ?
\tcblower
Les trois moments types d'un raisonnement philosophique classique sont :

\begin{enumerate}
\item \textbf{Le moment du constat ou de la position du problème :} l'auteur part d'un fait, d'une opinion commune ou d'un paradoxe qui appelle l'examen (ex. : « L'homme naît libre, et partout il est dans les fers »).
\item \textbf{Le moment de l'analyse ou de la réfutation :} l'auteur examine les solutions possibles, critique les opinions reçues ou les thèses adverses, et dégage les concepts nécessaires à la résolution du problème.
\item \textbf{Le moment de la synthèse ou de la solution doctrinale :} l'auteur énonce sa propre thèse, fondée sur les analyses précédentes, et en tire les conséquences.
\end{enumerate}

\textbf{Fonction dans l'économie du texte :} ces trois moments donnent au texte sa progression : on ne comprend la thèse finale que si l'on a suivi le chemin qui y conduit. Repérer ces moments, c'est dégager le \emph{plan interne} du texte, indispensable à l'explication.

\textbf{Point de vigilance :} certains textes sont dialectiques (thèse / antithèse / synthèse), d'autres purement analytiques ; le candidat doit s'adapter au texte réel et ne pas forcer un schéma.
\end{exoresolu}

\courssub{La lecture analytique du texte}

\begin{exoresolu}[Analyse guidée : Texte sur la liberté (Rousseau, \emph{Du contrat social}, L. I, ch. 1)]
Énoncé : À partir de l'extrait : « L'homme naît libre, et partout il est dans les fers. [...] Comment ce changement s'est-il fait ? Je l'ignore. Qu'est-ce qui peut le rendre légitime ? [...] ce droit ne vient pas de la nature, il est donc fondé sur des conventions. », identifier : 1. Le thème général. 2. La thèse défendue. 3. Les trois moments du raisonnement. 4. Une problématique pertinente.
\tcblower
\begin{enumerate}
\item \textbf{Thème général :} la liberté de l'homme et son fondement politique — plus précisément, le passage de la liberté naturelle à la liberté civile et la légitimité de l'ordre social.

\item \textbf{Thèse défendue :} l'ordre social, qui enchaîne l'homme naturellement libre, ne peut être légitime que s'il repose non sur la nature ou la force, mais sur des \emph{conventions} librement consenties (le contrat social). La légitimité politique est donc d'origine conventionnelle, non naturelle.

\item \textbf{Les trois moments du raisonnement :}
\begin{itemize}
\item \emph{Constat paradoxal :} « L'homme naît libre, et partout il est dans les fers » — l'écart entre ce qui est (la servitude) et ce qui devrait être (la liberté).
\item \emph{Questionnement :} « Comment ce changement s'est-il fait ? Je l'ignore. Qu'est-ce qui peut le rendre légitime ? » — Rousseau écarte la question historique pour poser la question de droit.
\item \emph{Annonce de la solution :} « ce droit ne vient pas de la nature, il est donc fondé sur des conventions » — la fondation conventionnelle de l'ordre social, que développera la théorie du contrat.
\end{itemize}

\item \textbf{Problématique :} « Comment l'ordre social, qui prive l'homme de sa liberté naturelle, peut-il être rendu légitime ? » (ou : « Sur quel fondement une autorité politique peut-elle légitimement enchaîner des hommes nés libres ? »).
\end{enumerate}

\textbf{Points de vigilance :} (a) la problématique doit porter sur la \emph{légitimité}, non sur l'origine historique — Rousseau dit explicitement « Je l'ignore » ; (b) la thèse doit mentionner les \emph{conventions}, mot-clé du texte ; (c) ne pas écrire que Rousseau condamne toute société : il cherche au contraire à la \emph{fonder} légitimement.
\end{exoresolu}

\courssub{La problématisation et l'introduction}

\begin{exoresolu}[Rédaction d'une introduction complète : Sujet « Le travail est-il une contrainte ? »]
Énoncé : Rédigez une introduction complète (accroche, présentation, problématique, annonce du plan) sur le sujet : « Le travail est-il une contrainte ? ».
\tcblower
\textbf{Corrigé type (les quatre temps sont signalés) :}

\emph{[1. Accroche]} Dès l'aube, dans les quartiers de Douala comme dans les champs de l'Ouest, des millions d'hommes et de femmes se mettent au travail, souvent sans joie, parce qu'il « faut bien vivre ». Le travail apparaît ainsi d'emblée comme ce à quoi l'on est obligé, ce qui s'impose à nous du dehors.

\emph{[2. Présentation du sujet]} Le \emph{travail} désigne l'activité par laquelle l'homme transforme la nature et produit les biens nécessaires à sa vie ; la \emph{contrainte} est ce qui s'impose à la volonté contre son gré, par la force ou la nécessité. Dire que le travail est une contrainte, c'est affirmer qu'on ne travaille que parce qu'on y est forcé — par la faim, par la loi, par un employeur. Mais le travail ne serait-il que cela ? N'est-il pas aussi activité de création, d'accomplissement et de reconnaissance sociale ?

\emph{[3. Problématique]} Le travail se réduit-il à une nécessité subie, ou peut-il être, au-delà de la contrainte, un moyen d'émancipation et de réalisation de soi ?

\emph{[4. Annonce du plan]} Nous verrons d'abord que le travail se présente bien, dans l'expérience commune, comme une contrainte extérieure ; nous montrerons ensuite qu'il est aussi transformation de soi et source de liberté ; nous demanderons enfin sous quelles conditions le travail peut cesser d'être pure servitude pour devenir accomplissement.

\textbf{Points de vigilance :} (a) la problématique ne doit pas recopier le sujet mot pour mot ; (b) les termes clés doivent être définis avant d'être discutés ; (c) l'annonce du plan doit être sobre (une phrase) et fidèle au développement qui suivra réellement.
\end{exoresolu}

\begin{exoresolu}[Comparaison de problématiques : Sartre et la mauvaise foi]
Énoncé : Deux élèves proposent une problématique pour commenter l'extrait de \emph{L'Être et le Néant} sur le garçon de café (mauvaise foi).
\begin{itemize}
\item Problématique A : « Qu'est-ce que la mauvaise foi ? »
\item Problématique B : « La mauvaise foi est-elle une fuite inévitable de la conscience ou une possibilité structurelle de la liberté ? »
\end{itemize}
Laquelle choisir et pourquoi ?
\tcblower
La problématique A (« Qu'est-ce que la mauvaise foi ? ») est une simple question de définition : elle appelle un exposé scolaire, ferme le débat avant de l'ouvrir et n'exige aucune discussion critique. La problématique B est supérieure sur les trois critères : \emph{ouverture au débat} (elle propose une alternative réelle — fuite inévitable ou possibilité structurelle — qui autorise une discussion) ; \emph{ancrage dans le texte} (l'exemple du garçon de café illustre précisément la conscience qui se « joue » elle-même pour fuir sa liberté) ; \emph{portée philosophique} (elle engage le rapport fondamental entre liberté, angoisse et conscience, cœur de l'existentialisme sartrien). C'est donc B qui convient à un commentaire de niveau Bac.

\textbf{Point de vigilance :} une bonne problématique n'est jamais une question de cours (« qu'est-ce que... ? ») mais une question qui met en tension deux thèses possibles et que le texte permet de trancher.
\end{exoresolu}

\courssub{Le développement : expliquer, illustrer, discuter}

\begin{exoresolu}[Chasse aux erreurs méthodologiques : Commentaire sur l'amitié (Aristote)]
Énoncé : Relevez et corrigez au moins trois fautes méthodologiques dans le paragraphe suivant :
« Moi, mon meilleur ami s'appelle Paul et on joue à la PlayStation ensemble, donc l'amitié c'est super important. Tout le monde aime avoir des amis, c'est normal. Aristote a raison sur toute la ligne quand il dit que l'amitié est nécessaire à la vie. »
\tcblower
Le paragraphe contient (au moins) trois fautes distinctes :

\begin{enumerate}
\item \textbf{Opinion non justifiée et exemple personnel hors de propos :} « Moi, mon meilleur ami s'appelle Paul et on joue à la PlayStation ensemble ». L'anecdote personnelle ne prouve rien philosophiquement : elle est particulière, non universalisable. \emph{Correction :} remplacer par une analyse du concept : Aristote distingue trois formes d'amitié (utile, agréable, vertueuse) ; seule l'amitié vertueuse, fondée sur la vertu des amis, est parfaite et nécessaire à la vie bonne.

\item \textbf{Paraphrase et lieu commun :} « tout le monde aime avoir des amis, c'est normal » répète l'opinion commune sans l'expliquer ni la discuter. \emph{Correction :} dégager l'argument d'Aristote : l'homme est un animal politique ; l'amitié (\emph{philia}) est le lien qui maintient la cité unie, plus encore que la justice — c'est en cela qu'elle est « nécessaire à la vie ».

\item \textbf{Absence de discussion critique (assentiment béat) :} « Aristote a raison sur toute la ligne ». Un commentaire exige une évaluation : la thèse est-elle sans limites ? \emph{Correction :} discuter : l'amitié vertueuse est rare et exigeante ; peut-on fonder le bonheur sur un bien qui dépend d'autrui ? On pourrait objecter, avec les stoïciens, que le sage doit se suffire à lui-même.
\end{enumerate}

(On pouvait aussi relever le \textbf{contresens} : réduire l'amitié aristotélicienne au « bonheur » sentimental, alors qu'elle est d'abord une vertu éthique et politique.)

\textbf{Point de vigilance :} chaque faute identifiée doit être accompagnée d'une correction précise, ancrée dans la doctrine d'Aristote, et non d'une simple remarque de forme.
\end{exoresolu}

\begin{exoresolu}[Paragraphe de développement type : Locke et la propriété (\emph{Second Traité du gouvernement}, §27)]
Énoncé : Rédigez un paragraphe de développement expliquant et illustrant l'argument du « mélange » du travail et de la nature chez Locke. Intégrez une illustration camerounaise.
\tcblower
\textbf{Corrigé type :}

Locke affirme que la propriété privée trouve son origine légitime dans le travail : les ressources de la nature sont données à tous en commun, mais chaque homme possède au moins une chose en propre — son corps et donc son travail. Le concept décisif est celui de \emph{mélange} : lorsqu'un homme « mêle » son travail à une portion de nature — en défrichant une terre, en cueillant des fruits —, il y incorpore quelque chose qui lui appartient indiscutablement ; la chose travaillée devient ainsi sienne, car on ne peut la séparer du labeur qui s'y est déposé. L'\cle{appropriation} n'est donc ni un vol ni un privilège arbitraire : elle est l'extension naturelle de la propriété que chacun a de sa propre personne. Une illustration camerounaise l'éclaire : un cultivateur de Bafia qui défriche une parcelle en friche, y plante ses arachides et l'entretient pendant des mois fait de cette terre, par son travail, \emph{sa} plantation — nul ne conteste dans le village que la récolte lui revienne, précisément parce que la terre porte la marque de son labeur. De même, le menuisier de Douala qui transforme un tronc d'iroko en meuble est propriétaire de l'objet, car il a joint son travail à la matière première. Ainsi, pour Locke, le travail est le titre originel et légitime de toute propriété — thèse fondatrice, mais dont on pourra discuter les limites : que devient ce droit lorsque la terre commune vient à manquer pour les autres ?

\textbf{Points de vigilance :} (a) le paragraphe doit suivre l'ordre : reformulation $\rightarrow$ analyse des concepts $\rightarrow$ illustration $\rightarrow$ mini-conclusion ; (b) l'exemple camerounais doit \emph{illustrer} l'argument (montrer le « mélange » du travail et de la nature), pas le remplacer ; (c) la phrase finale ouvre la discussion sans la développer — c'est le rôle de la transition.
\end{exoresolu}

\courssub{La conclusion et la relecture}

\begin{aretenir}[La conclusion réussie : Bilan, Réponse, (Ouverture)]
Une conclusion de commentaire philosophique se structure en deux temps obligatoires et un facultatif :
\begin{enumerate}
\item \textbf{Le bilan (synthèse) :} en une ou deux phrases, rappeler la thèse de l'auteur et le chemin parcouru dans le développement (les grandes étapes de l'explication et de la discussion). Pas de nouvelle idée.
\item \textbf{La réponse à la problématique :} formuler explicitement la réponse apportée à la question posée en introduction, en tenant compte des nuances de la discussion.
\item \textbf{L'ouverture (facultative) :} seulement si elle surgit naturellement de la discussion (prolongement vers un autre auteur, un problème voisin, une actualité). Une citation plaquée artificiellement pénalise la copie.
\end{enumerate}
\textbf{Relecture (5 min) :} vérifier l'orthographe des noms propres (Kant, Rousseau, Aristote...), l'accord des participes passés, la cohérence des connecteurs logiques (donc, cependant, par conséquent), la présence des quatre temps en introduction, l'adéquation plan/annonce.
\end{aretenir}

\courssub{Entraînement intégré : commentaire complet guidé (complément Bac)}

\begin{exoresolu}[Sujet complet type Probatoire / Bac : Commentaire du texte de Kant sur le devoir (\emph{Fondements de la métaphysique des mœurs}, 1785)]
Énoncé : Commenter le texte suivant : « Le devoir est la nécessité d'une action par respect pour la loi. [...] L'action a sa valeur morale non pas dans le but qu'elle vise, mais dans la maxime selon laquelle elle est entreprise. [...] L'autonomie de la volonté est le principe suprême de la morale. »
\tcblower
\textbf{Barème indicatif (sur 20) :} Introduction (4 pts) ; Explication du texte fondue dans le développement (8 pts) ; Discussion critique (5 pts) ; Conclusion (2 pts) ; Langue et présentation (1 pt).

\textbf{Corrigé type (plan détaillé rédigé en éléments) :}

\emph{Introduction.} [Accroche] Nous obéissons chaque jour à des règles — au code de la route, aux règlements scolaires — parce qu'une autorité extérieure nous y oblige. Mais y a-t-il des obligations qui ne viennent pas du dehors ? [Présentation] Le texte, extrait des \emph{Fondements de la métaphysique des mœurs} (1785) de Kant, traite du \emph{devoir} et du fondement de la morale. [Problématique] En quoi le devoir, loin d'être une contrainte hétéronome, est-il l'expression de la liberté rationnelle, c'est-à-dire de l'autonomie ? [Plan] Nous expliquerons d'abord la thèse kantienne (le devoir comme loi que la raison se donne), nous en dégagerons ensuite le critère de la valeur morale (le désintéressement), puis nous discuterons la portée et les limites de cette morale du devoir pur.

\emph{Développement — Partie 1 : le devoir n'est pas une contrainte extérieure.} Kant oppose d'emblée deux conceptions : le devoir comme \emph{hétéronomie} (loi reçue du dehors : force, tradition, intérêt) et le devoir comme \emph{autonomie} (« la loi que notre raison nous donne à nous-mêmes »). Obéir au devoir, c'est donc obéir à sa propre raison : loin d'aliéner la liberté, le devoir en est la manifestation la plus haute. Le mot \emph{respect} est décisif : agir par devoir, c'est agir par respect de la loi morale, sentiment unique car il est produit par la raison elle-même et non par une inclination sensible.

\emph{Partie 2 : le critère de la valeur morale.} Kant déplace la valeur morale des \emph{résultats} vers l'\emph{intention} : « l'action a sa valeur morale non pas dans le but qu'elle vise, mais dans la maxime selon laquelle elle est entreprise ». Une action utile ou agréable n'est pas pour autant bonne ; seule est morale l'action faite \emph{par devoir}, et non simplement \emph{conforme} au devoir. Le commerçant honnête par calcul agit conformément au devoir, mais non par devoir : sa conduite n'a pas de valeur morale proprement dite. L'autonomie de la volonté est ainsi « le principe suprême de la morale » : la dignité de l'homme tient à sa capacité de se légiférer à soi-même.

\emph{Partie 3 : discussion.} \emph{Force de la thèse :} elle fonde l'universalité de la morale (la loi vaut pour tout être raisonnable) et la dignité humaine, indépendamment des circonstances — repère exigeant pour juger, par exemple, la corruption : peu importe que le détournement « réussisse », il est mal parce que sa maxime est immorale. \emph{Limites :} (a) le formalisme : Kant ne dit pas le contenu concret des devoirs ; le rigorisme peut conduire à des cas difficiles (faut-il dire la vérité au meurtrier ?) ; (b) l'abstraction : Durkheim objecte que le devoir est d'abord \emph{social}, intériorisé par l'éducation ; Mill et les utilitaristes soutiennent que la valeur d'une action se mesure à ses conséquences ; (c) la psychologie : nul n'agit jamais sans aucune inclination — le « devoir pur » est-il humainement possible ? Dans le contexte camerounais, on peut interroger la tension entre devoirs traditionnels (envers la famille élargie, la communauté) et devoirs modernes (envers l'État, la loi) : la raison suffit-elle à arbitrer ?

\emph{Conclusion.} [Bilan] Kant renverse la conception commune : le devoir n'enchaîne pas, il libère, car il est la loi que la raison se prescrit ; la valeur morale réside dans l'intention désintéressée. [Réponse à la problématique] Le devoir est bien l'expression de la liberté rationnelle — à condition d'admettre que la vraie liberté n'est pas de faire ce qu'on veut, mais de vouloir ce que la raison ordonne. [Ouverture possible] Reste à savoir si une morale de l'intention peut suffire à juger les actions dans la complexité de la vie sociale.

\textbf{Points de vigilance :} (a) l'explication ne doit \emph{jamais} être un bloc séparé suivi d'une discussion déconnectée : elle se fond dans le développement ; (b) définir \emph{autonomie} (se donner sa loi) contre \emph{hétéronomie} (recevoir sa loi du dehors) est obligatoire ; (c) la discussion doit citer au moins une objection doctrinale (Durkheim, Mill) et une application concrète ; (d) interdiction de paraphraser le texte ligne à ligne.
\end{exoresolu}

\begin{exoresolu}[Dissertation corrélée : « Obéir, est-ce renoncer à sa liberté ? »]
Énoncé : Traiter le sujet de dissertation : « Obéir, est-ce renoncer à sa liberté ? » En mobilisant Kant (autonomie/hétéronomie), Rousseau (liberté naturelle/liberté civile) et un exemple camerounais précis.
\tcblower
\textbf{Barème indicatif :} Introduction (4 pts) ; trois parties équilibrées (12 pts) ; Conclusion (3 pts) ; langue (1 pt). Les trois références imposées (Kant, Rousseau, exemple camerounais) sont chacune attendues explicitement.

\textbf{Corrigé type (plan détaillé) :}

\emph{Introduction.} [Accroche] L'enfant obéit à ses parents, l'élève au règlement, le citoyen à la loi : l'obéissance semble le contraire de la liberté, puisque j'agis selon la volonté d'un autre. [Définitions] \emph{Obéir} : se conformer à un ordre ou une loi ; \emph{liberté} : pouvoir d'agir selon sa propre volonté. La tension est nette : obéir, n'est-ce pas toujours céder sa volonté ? [Problématique] Toute obéissance est-elle aliénation, ou existe-t-il une obéissance qui réalise la liberté au lieu de la détruire ? [Plan] Nous verrons d'abord que l'obéissance hétéronome est bien renoncement à la liberté ; puis qu'il existe une obéissance autonome qui est liberté même ; enfin nous demanderons à quelles conditions l'obéissance sociale reste légitime.

\emph{Partie 1 — L'obéissance comme hétéronomie : le renoncement.} Quand j'obéis par peur, par intérêt ou par habitude, ma volonté est gouvernée du dehors : c'est l'\emph{hétéronomie} (Kant). L'esclave obéit, le sujet soumis à la propagande obéit : ils agissent, mais ne décident pas. Rousseau le constate : « l'homme naît libre, et partout il est dans les fers » — la \emph{liberté naturelle} (faire ce qu'on veut selon sa seule force) est perdue dans l'état social. Obéir à la force n'est jamais un devoir : « la force est une puissance physique, je ne vois pas quelle moralité peut résulter de ses effets » (Rousseau). L'obéissance servile est donc bien renoncement.

\emph{Partie 2 — L'obéissance autonome : la liberté réalisée.} Mais Kant montre qu'obéir à la loi que ma raison se donne, c'est l'\emph{autonomie} : je n'obéis alors qu'à moi-même en tant qu'être raisonnable. De même Rousseau : par le contrat social, le citoyen n'obéit qu'à la volonté générale, c'est-à-dire à une loi qu'il s'est prescrite ; il échange la liberté naturelle contre la \emph{liberté civile} et morale, qui est « obéissance à la loi qu'on s'est prescrite ». Obéir n'est donc pas renoncer à sa liberté quand la loi est rationnellement voulue : c'est au contraire cesser d'être l'esclave de ses passions.

\emph{Partie 3 — Conditions d'une obéissance légitime (application).} L'obéissance est légitime si la loi vise le bien commun, si elle s'applique à tous également et si le citoyen peut la reconnaître comme sienne. Exemples camerounais : respecter le code de la route à Yaoundé n'aliène pas ma liberté — il protège ma vie et celle d'autrui, et je peux en vouloir rationnellement les règles ; le règlement intérieur d'un lycée technique oblige, mais il rend possible l'apprentissage de tous ; la loi sur la protection de l'environnement limite certaines pratiques (déforestation abusive) au nom d'un bien commun. En revanche, obéir à un ordre injuste ou absurde reste une servitude : la désobéissance civile peut alors être un devoir.

\emph{Conclusion.} Obéir n'est renoncer à sa liberté que lorsque la loi m'est étrangère (hétéronomie). Lorsque j'obéis à une loi que ma raison reconnaît comme juste et que je me suis en quelque sorte prescrite (autonomie kantienne, volonté générale rousseauiste), l'obéissance est la forme sociale de la liberté. Tout dépend donc non du fait d'obéir, mais de \emph{à quoi} et \emph{pourquoi} l'on obéit.

\textbf{Points de vigilance :} (a) les deux distinctions imposées (autonomie/hétéronomie ; liberté naturelle/liberté civile) doivent être nommées et définies, pas seulement évoquées ; (b) l'exemple camerounais doit être précis et analysé (en quoi cette obéissance est-elle légitime ?), non plaqué ; (c) la troisième partie doit trancher la problématique par une position nuancée et justifiée.
\end{exoresolu}

\newpage

\coursec{Fiche bilan}

\begin{popfigure}
\begin{tikzpicture}[
    mindmap,
    grow cyclic,
    every node/.style={concept, execute at begin node=\hskip0pt},
    concept color=popblue,
    level 1/.append style={level distance=4.5cm, sibling angle=60},
    level 2/.append style={level distance=3cm, sibling angle=45},
    texte court/.style={text width=2.8cm, align=center, font=\small\bfseries},
    branch1/.style={concept color=poporange},
    branch2/.style={concept color=popgreen},
    branch3/.style={concept color=poppurple},
    branch4/.style={concept color=poppink},
    branch5/.style={concept color=popgold},
    branch6/.style={concept color=popblue!70!popgreen}
]
\node[texte court, minimum size=2.2cm, align=center]{Commentaire\\ Philosophique}
    child[branch1] { node[texte court, align=center]{1. Lecture\\ Analytique} }
    child[branch2] { node[texte court, align=center]{2. Problématisation\\ \& Introduction} }
    child[branch3] { node[texte court, align=center]{3. Développement\\ Expliquer} }
    child[branch4] { node[texte court, align=center]{4. Développement\\ Discuter} }
    child[branch5] { node[texte court, align=center]{5. Conclusion\\ \& Relecture} }
    child[branch6] { node[texte court, align=center]{6. Entraînement\\ Intégré} };
\end{tikzpicture}
\legende{Carte mentale : Les 6 étapes clés du commentaire philosophique en Première Technique.}
\end{popfigure}

\begin{bilan}
\courssub{Définitions clés}
\begin{itemize}
    \item \cle{Commentaire philosophique} : Exercice consistant à \textbf{expliquer}, \textbf{analyser} et \textbf{discuter} la pensée d'un auteur à partir d'un texte donné, sans y substituer sa propre opinion.
    \item \cle{Problématique} : Question philosophique précise à laquelle le texte tente de répondre ; elle relie le \emph{sujet} (thème) à la \emph{thèse} (réponse de l'auteur).
    \item \cle{Thèse} : Réponse affirmée par l'auteur au problème posé ; elle doit être formulée en une phrase claire.
    \item \cle{Structure du texte} : Enchaînement logique des arguments (ex: thèse $\rightarrow$ arguments $\rightarrow$ objections $\rightarrow$ réponse $\rightarrow$ conséquence).
    \item \cle{Concept} : Notion philosophique opératoire dans le texte (ex: liberté, conscience, État, technique, bonheur) ; il faut le définir dans le contexte de l'auteur.
\end{itemize}

\courssub{Méthode express : Les 4 temps du commentaire (Bac/Probatoire)}
\begin{center}
\begin{tabularx}{\linewidth}{|>{\bfseries}l|X|}\hline
\textbf{Temps} & \textbf{Actions clés} \\\hline
1. \textbf{Lecture \& Brouillon} (30--40 min) & Identifier : Auteur, œuvre, date, thème. Repérer : Thèse, concepts, structure, connecteurs logiques. Formuler : Problématique. \\\hline
2. \textbf{Introduction} (10--15 min) & Accroche $\rightarrow$ Présentation texte $\rightarrow$ \textbf{Problématique} $\rightarrow$ \textbf{Thèse} $\rightarrow$ Annonce du plan (2--3 parties). \\\hline
3. \textbf{Développement} (1h30--2h) & \textbf{Partie A : Expliquer} (suivre le fil du texte, définir concepts, reformuler arguments).\\
& \textbf{Partie B : Discuter} (Portée/Limites : exemples concrets CM, contre-exemples, autres auteurs au programme). \\\hline
4. \textbf{Conclusion} (10 min) & Réponse synthétique à la problématique $\rightarrow$ Ouverture (autre auteur, actualité, autre notion). \\\hline
\end{tabularx}
\end{center}

\courssub{Principes d'or à connaître par cœur}
\begin{itemize}
    \item \textbf{Fidélité au texte} : Ne jamais importer des connaissances extérieures dans la partie \emph{Expliquer}.
    \item \textbf{Unité de paragraphe} : 1 paragraphe = 1 idée principale = 1 argument du texte (partie A) ou 1 critique/approfondissement (partie B).
    \item \textbf{Transition} : Obligatoire entre chaque grande partie (rappeler la problématique, annoncer le virage).
    \item \textbf{Exemples camerounais} : Privilégier le contexte local (coutumes, droit OHADA, bilingue, agriculture, urbanisation Yaoundé/Douala) pour illustrer/discuter.
    \item \textbf{Relecture} : Vérifier l'orthographe des noms d'auteurs, la cohérence thèse/problématique/plan, la syntaxe.
\end{itemize}
\end{bilan}

\begin{aretenir}[Checklist avant le Bac / Probatoire]
$\square$ J'ai lu le texte \textbf{3 fois minimum} (global, analytique, structurel).\\
$\square$ J'ai formulé une \textbf{problématique} en forme de question fermée (ouverte mais précise) issue du texte.\\
$\square$ J'ai identifié la \textbf{thèse} de l'auteur en une phrase affirmative.\\
$\square$ Mon \textbf{plan détaillé} (2 ou 3 parties, sous-parties) répond \textbf{exactement} à la problématique.\\
$\square$ Dans la partie \textbf{Expliquer} : je suis le texte pas à pas, je définis les concepts \emph{in situ}, je cite le texte (guillemets).\\
$\square$ Dans la partie \textbf{Discuter} : j'évalue la portée ET les limites, j'apporte \textbf{2 exemples concrets camerounais minimum}.\\
$\square$ Mes \textbf{transitions} font le lien logique (pas de "Ensuite", "De plus" seuls).\\
$\square$ Ma \textbf{conclusion} répond à la problématique et propose une \textbf{ouverture} pertinente.\\
$\square$ J'ai relu : \textbf{orthographe auteurs/concepts}, \textbf{accords}, \textbf{ponctuation}, \textbf{lisibilité}.\\
$\square$ J'ai géré mon \textbf{temps} : Brouillon $\le$ 40 min, Rédaction $\le$ 2h30, Relecture $\ge$ 10 min.\\
$\square$ Je maîtrise les \textbf{références au programme} (auteurs, notions de l'année) pour la discussion.
\end{aretenir}

\findecours

\end{document}
